\chapter{Galois representations of modular forms, level lowering and congruences}\label{ch:modgalrep}
\chaptermark{Galois representations of modular forms}

Section 3 of the paper proves Theorem thm:qual by a chain of four moves: attach a Galois representation to a newform, lower the
level of its reduction modulo $\ell$, twist, and turn a congruence between eigenforms into a divisibility of the congruence number.
Section 4 of the paper reads the same mechanism quantitatively: a depleted form is a form whose level has been raised, the change of the
congruence ideal is the level-raising Euler factor, and the twist map transports congruences. This chapter supplies all of it. It
constructs the representation $\rho_{f,\lambda}$ of a newform of weight $2$ and compares it with $\rho_{E,\ell}$; it describes the restriction of
$\rho_{f,\lambda}$ to a decomposition group at every prime (local-global compatibility, the local types, the conductor of a twist); it states Serre's
conjecture and the theorems of Mazur, Ribet, Carayol and Diamond that lower the level, and Ribet's theorem that raises it; and it proves in
full the paper's Lemma lem:crit, Lemma lem:old and Theorem thm:qual, with the paper's Example. The algebra is \Cref{ch:modules} and
\Cref{ch:reps}, the local Galois theory \Cref{ch:local}, everything specific to $\rho_{E,\ell}$ \Cref{ch:galrep}, the $q$-expansion side of Hecke
operators \Cref{ch:hecke}, and the Jacobian $J_0(N)$, modularity and the modular degree \Cref{ch:moddeg}. The deep imported theorems
(Deligne, Carayol, Khare and Wintenberger, Ribet, Diamond, Diamond, Flach and Guo) are stated precisely, with the idea of the proof and a
reference.

\section*{Core results at a glance}

Throughout, $f$ is a normalised newform of weight $2$ on $\Gamma_0(N)$ with coefficient field $K_f$, $\lambda\mid\ell$ a prime of $K_f$, $\T$ the full Hecke algebra of level $N$, and
$\phi_f\colon\T\to\mathcal{O}_f$ the eigencharacter; $\Frob_p$ is an arithmetic Frobenius; $q$ denotes a prime except in $q$-expansions.

\begin{enumerate}[label=\textbf{\thechapter.\arabic*}]
\item \textbf{Conventions.} $\rho_{f,\lambda}$ is normalised by $\det(X-\rho_{f,\lambda}(\Frob_p))=X^2-a_pX+p$ ($p\nmid N\ell$), $\det\rho_{f,\lambda}=\chi_\ell$; $H^1$ realises the dual. $\chi_{q^*}$ is the quadratic
character of conductor $q$; $\chi_\ell$ is cyclotomic. \emph{(convention)} \Cref{conv:modgalrep:notation}.
\item \textbf{Twists.} For $\chi$ quadratic of conductor $r$: the naive twist $g\otimes\chi=\sum\chi(n)a_n(g)q^n$, the newform twist $g^\chi$ ($\rho_{g^\chi}\simeq\rho_g\otimes\chi$), the depletion $h^\Sigma$;
$g\otimes\chi=(g^\chi)^{\Sigma(r)}$, and $g\otimes\chi=g^\chi$ iff $a_q(g^\chi)=0$ for all $q\mid r$. \emph{(convention)} \Cref{conv:modgalrep:twists}.
\item \textbf{Newforms.} $a_{mn}=a_ma_n$ for $(m,n)=1$, $a_{p^{k+1}}=a_pa_{p^k}-pa_{p^{k-1}}$ ($p\nmid N$), $a_{p^k}=a_p^k$ ($p\mid N$); $a_p\in\{\pm1\}$ if $p\,\|\,N$, $a_p=0$ if $p^2\mid N$;
strong multiplicity one; $S_2(\Gamma_0(N))=\bigoplus_h V_h$ with $V_h=\langle h(dz):d\mid N/M\rangle$ $\T$-stable and mutually orthogonal. \emph{(stated; proved in \Cref{ch:hecke})}
\Cref{prop:modgalrep:newforms-recap}.
\item \textbf{Eichler-Shimura.} $J_f=J_0(N)/I_fJ_0(N)$ has dimension $[K_f:\Q]$ and good reduction outside $N$; $V_\ell(J_f)$ is free of rank $2$ over $K_f\otimes\Ql$, and
$\rho_{f,\lambda}$ on $V_\lambda(f)$ is continuous, unramified outside $N\ell$, with $\det(X-\rho_{f,\lambda}(\Frob_p))=X^2-a_p(f)X+p$ for $p\nmid N\ell$, $\det=\chi_\ell$, odd; $H^1_{\mathrm{et}}(X_0(N))[I_f]$ realises
$V_\lambda(f)^\vee\simeq V_\lambda(f)(-1)$. \emph{(proof sketched here; Shimura, Deligne)} \Cref{thm:modgalrep:eichler-shimura}.
\item \textbf{Irreducibility.} $\rho_{f,\lambda}$ is absolutely irreducible. \emph{(proof sketched here; Ribet)} \Cref{prop:modgalrep:irreducible}.
\item \textbf{Deligne.} A newform of weight $k\ge2$ and character $\varepsilon$ has $\rho_{f,\lambda}$ unramified outside $N\ell$ with $\det(X-\rho(\Frob_p))=X^2-a_pX+\varepsilon(p)p^{k-1}$, $\det=\varepsilon\chi_\ell^{k-1}$,
odd. \emph{(stated; reference)} \Cref{thm:modgalrep:deligne}.
\item \textbf{Lattices.} Stable lattices exist; $\bar\rho^{\mathrm{ss}}$ is independent of the lattice; for $\ell$ odd $\bar\rho$ is irreducible iff absolutely irreducible; if $\bar\rho$ is
irreducible all stable lattices are homothetic. \emph{(proved here, given \Cref{ch:reps,ch:galois})} \Cref{prop:modgalrep:lattices}.
\item \textbf{Congruent newforms.} $a_p(f)\equiv a_p(h)\pmod\lambda$ on a density-one set gives $\bar\rho_f^{\mathrm{ss}}\simeq\bar\rho_h^{\mathrm{ss}}$, and $\bar\rho_h\simeq\bar\rho_f$ if $\ell$ is odd and $\bar\rho_f$
irreducible; conversely; and equality on a density-one set gives $f=h$. \emph{(proved here, given \Cref{ch:galois})} \Cref{cor:modgalrep:congruent-newforms}.
\item \textbf{$\rho_{f_E,\ell}$ and $\rho_{E,\ell}$.} $J_f=E_0$ and $\rho_{f,\ell}\simeq\rho_{E,\ell}\otimes\Ql$. The lattice $M_{f,\lambda}=H^1(X_0(N),\Zl)[\ker\phi_f]$ equals $\varphi_{E_0}^*H^1(E_0,\Zl)$, which is isomorphic to
$T_\ell E_0(-1)$, without hypotheses; it is homothetic to $(T_\ell E)^\vee$ when $\bar\rho$ is irreducible; $\ad M_{f,\lambda}\simeq\ad T_\ell E_0$. \emph{(proved here, given \Cref{ch:moddeg})}
\Cref{thm:modgalrep:comparison-E}.
\item \textbf{Local Langlands for $\GL_2(\Q_p)$.} A bijection $\pi\mapsto\operatorname{rec}(\pi)$ to two-dimensional Frobenius-semisimple Weil-Deligne representations, compatible with
twists, preserving conductors and matching principal series, twists of Steinberg and supercuspidals with sums of characters, $\Sp(2)\otimes\mu$ and irreducible
representations. \emph{(stated; reference)} \Cref{thm:modgalrep:local-langlands}.
\item \textbf{Local components.} $\pi_f=\bigotimes'\pi_v$; $\pi_p$ has trivial central character and conductor exponent $v_p(N)$. \emph{(definition; facts stated with references)}
\Cref{def:modgalrep:local-component}.
\item \textbf{Carayol.} For $p\ne\ell$: $\WD_p(f)$ corresponds to $\pi_p$ and is independent of $\ell$; $a(\rho_{f,\lambda}|_{G_{\Q_p}})=v_p(N)$, so the conductor of $\rho_{f,\lambda}$ is $N$
($\ell\nmid N$); $L_p(f,s)=L(V_\lambda(f)|_{G_{\Q_p}},s-1)$. \emph{(stated; reference)} \Cref{thm:modgalrep:carayol}.
\item \textbf{Local types.} Exactly one of (U), (St), (RSt), (RPS), (SC) holds at $p\ne\ell$; $\delta_p=2,1,0$ for $p\nmid N$, $p\,\|\,N$, $p^2\mid N$; $a_p=0$ and $V^{I_p}=0$ if $p^2\mid N$; in case
(St), $\rho|_{G_{\Q_p}}\simeq\sm{\varepsilon\chi_\ell}{*}{0}{\varepsilon}$ on every stable lattice with $\varepsilon(\Frob_p)=a_p=\pm1$, the reduction is non-split with a unique stable line when
ramified, and $\bar\rho(\Frob_p)$ has eigenvalues $a_pp$, $a_p$ when unramified. \emph{(proved here, given \Cref{thm:modgalrep:carayol})} \Cref{thm:modgalrep:local-types}.
\item \textbf{Kodaira types.} For $f=f_E$ at $q\ge5$: type $\kod{I}_n$ gives (St), $\kod{I}_n^*$ gives (RSt), $\kod{I}_0^*$ gives (RPS); types $\kod{II}$, $\kod{II}^*$, $\kod{IV}$, $\kod{IV}^*$ give (RPS) or (SC) as
$q\equiv1$ or $2\pmod3$, and $\kod{III}$, $\kod{III}^*$ give (RPS) or (SC) as $q\equiv1$ or $3\pmod4$. \emph{(proved here, given \Cref{ch:galrep})} \Cref{prop:modgalrep:kodaira-dictionary}.
\item \textbf{$\delta_p$ and the naive factor.} For $p\mid N$: $a_p=\pm1$ iff $\delta_p=1$, $a_p=0$ iff $\delta_p=0$; the naive factor equals the true one iff $\delta_p>0$; the true
factor is twist-invariant, the naive one is not. \emph{(proved here)} \Cref{cor:modgalrep:delta-naive}.
\item \textbf{Conductor of a quadratic twist.} For $q$ odd and $\chi$ ramified quadratic: (U)$\mapsto2$, (St)$\mapsto2$, (RSt)$\mapsto1$, (RPS) with $\mu_i|_{I}=\chi|_I\mapsto0$, tame (RPS)/(SC) of
inertial order $3,4,6\mapsto2$. \emph{(proved here, given \Cref{ch:local})} \Cref{prop:modgalrep:twist-conductor}.
\item \textbf{Atkin-Li (Lemma lem:twist(iii)).} If $q$ is odd and $q^2\nmid M$, then $g\otimes\chi_{q^*}=g^{\chi_{q^*}}$ is a newform of level $\lcm(M,q^2)$, with $a_q=0$. \emph{(proved here, given
\Cref{thm:modgalrep:carayol})} \Cref{cor:modgalrep:atkin-li}.
\item \textbf{Twists of newforms.} $g^\chi$ exists and is unique, of level $M\prod_{q\mid r}q^{a(\WD_q\otimes\chi)-a(\WD_q)}$ dividing $\lcm(M,r^2)$; $a_n(g^\chi)=\chi(n)a_n(g)$ for $(n,r)=1$;
$g\otimes\chi=(g^\chi)^{\Sigma(r)}$; $(g^\chi)^\chi=g$; $f_{E^{(d)}}=(f_E)^{\chi_d}$. \emph{(proved here, given \Cref{thm:modgalrep:carayol}, \Cref{ch:hecke})} \Cref{prop:modgalrep:twists}.
\item \textbf{At $\ell\nmid N$.} $J_f[\ell^k]$ and its subquotients are finite flat, so $\bar\rho|_{D_\ell}$ is finite (good); $\rho|_{G_{\Q_\ell}}$ is crystalline with Hodge-Tate weights $0,1$ and
crystalline Frobenius polynomial $X^2-a_\ell X+\ell$; Fontaine-Laffaille range for $\ell\ge3$; Serre weight $2$. \emph{(proof sketched here; references)} \Cref{thm:modgalrep:crystalline}.
\item \textbf{Deligne, Fontaine.} For $\ell$ odd and $\ell\nmid N$: if $a_\ell\not\equiv0$, $\bar\rho|_{D_\ell}\simeq\sm{\bar\chi_\ell\mu(\bar a_\ell^{-1})}{*}{0}{\mu(\bar a_\ell)}$; if $a_\ell\equiv0$, $\bar\rho|_{I_\ell}\simeq\omega_2\oplus\omega_2^\ell$ and $\bar\rho|_{D_\ell}$
is irreducible \emph{(stated; reference)}; hence $a_\ell\bmod\lambda$ is $0$ or the Frobenius eigenvalue on the unique unramified quotient, and is determined by $\bar\rho|_{D_\ell}$
\emph{(proved here)}. \Cref{thm:modgalrep:deligne-fontaine}.
\item \textbf{Serre's invariants.} Conductor $N(\bar\sigma)=\prod_{p\ne\ell}p^{a_p(\bar\sigma)}$, weight $k(\bar\sigma)$, character $\varepsilon(\bar\sigma)$; ``arises from'' and ``modular''. \emph{(definition)}
\Cref{def:modgalrep:serre-invariants}.
\item \textbf{$N(\bar\rho_f)\mid N$.} $a_p(\bar\rho^{\mathrm{ss}})\le a_p(\bar\rho)=v_p(N)-(\dim\bar\rho^{I_p}-\delta_p)\le v_p(N)$ for $p\ne\ell$. \emph{(proved here, given \Cref{thm:modgalrep:carayol})}
\Cref{prop:modgalrep:serre-conductor-divides}.
\item \textbf{Serre's conjecture (Khare, Wintenberger).} Every odd absolutely irreducible $\bar\sigma\colon G_\Q\to\GL_2(\Fbar_\ell)$ arises from a newform of weight $k(\bar\sigma)$, level $N(\bar\sigma)$ and
character lifting $\varepsilon(\bar\sigma)$. \emph{(stated; reference)} \Cref{thm:modgalrep:serre-conjecture}.
\item \textbf{Mazur's principle.} $\ell$ odd, $\ell\nmid N$, $\bar\rho$ irreducible, $q\,\|\,N$, $\bar\rho$ unramified at $q$, $q\not\equiv1\pmod\ell$: $\bar\rho$ arises from level $N/q$. \emph{(stated; idea of proof)}
\Cref{thm:modgalrep:mazur-principle}.
\item \textbf{Ribet.} $\ell$ odd, $\ell\nmid N$, $\bar\rho$ irreducible, $q\,\|\,N$, $\bar\rho$ unramified at $q$: $\bar\rho$ arises from level $N/q$. \emph{(stated; idea of proof)} \Cref{thm:modgalrep:ribet}.
\item \textbf{The optimal level (DDT Theorem 3.15).} $\ell$ odd, $\bar\sigma$ irreducible and modular, absolutely irreducible on $\Gal(\Qbar/\Q(\sqrt{-3}))$ if $\ell=3$: $\bar\sigma$ arises in weight $2$,
level $N(\bar\sigma)\ell^{\delta(\bar\sigma)}$, character of order prime to $\ell$, $\delta=0$ if $\bar\sigma|_{D_\ell}$ is good. \emph{(stated; reference)} \Cref{thm:modgalrep:optimal-level}.
\item \textbf{For elliptic curves.} $\ell$ odd, $\ell\nmid N'$, $\bar\rho_{E',\ell}$ irreducible, image not $N_s(\F_3)$ if $\ell=3$: $\bar\rho_{E',\ell}$ arises from a newform on $\Gamma_0(N(\bar\rho))$, trivial character,
and $N(\bar\rho)\mid N'$. \emph{(proved here, given \Cref{thm:modgalrep:optimal-level})} \Cref{cor:modgalrep:optimal-level-E}.
\item \textbf{Level raising (Ribet).} For $q\nmid N\ell$: $\bar\rho$ arises from a form new at $q$ of level dividing $Nq$ iff $a_q^2\equiv(q+1)^2\pmod\lambda$; then $a_q(f)\equiv(q+1)a_q(g)$. \emph{(necessity
proved here; sufficiency stated, reference)} \Cref{thm:modgalrep:level-raising}.
\item \textbf{The level-raising factor.} $(1-\frac\alpha{\beta p})(1-\frac1p)(1-\frac\beta{\alpha p})=\frac{p-1}p\cdot\frac{(p+1)^2-a_p^2}{p^2}=L_p(A_f,1)^{-1}$ at good $p$; $1-p^{-2}$ at $p\,\|\,N$. \emph{(proved here)}
\Cref{prop:modgalrep:level-raising-factor}.
\item \textbf{Theorem thm:dfg14.} $\eta^\Sigma_{f,\lambda}=\eta^\emptyset_{f,\lambda}\prod_{p\in\Sigma}L_p^{\nv}(A_f,1)^{-1}$ for $\ell\nmid2M$, $\bar\rho$ irreducible. \emph{(stated; proof idea in \Cref{ch:congruence})}
\Cref{thm:modgalrep:dfg14}.
\item \textbf{Ihara's lemma.} For $q\nmid N\ell$ and $\mathfrak m$ non-Eisenstein, $H^1(X_0(N),\F_\ell)^2_{\mathfrak m}\to H^1(X_0(Nq),\F_\ell)$, $(x,y)\mapsto\alpha_1^*x+\alpha_q^*y$, is injective. \emph{(stated;
reference)} \Cref{lem:modgalrep:ihara}.
\item \textbf{The congruence number.} $r_f$ is the largest $r$ with $f\equiv g\pmod r$, $g\in f^\perp\cap S_2(\Gamma_0(N);\Z)$. \emph{(definition)} \Cref{def:modgalrep:congruence-number}.
\item \textbf{Agashe, Ribet, Stein.} $\deg\varphi_{E_0}\mid r$, with equal $p$-adic valuations for $p^2\nmid N$; $r(\curve{99a1})=12$, $\deg\varphi=4$. \emph{(stated; reference; data computed here)}
\Cref{thm:modgalrep:ars}.
\item \textbf{Lemma lem:crit.} $I=\ker\phi_f$, $J=\Ann_\T(f^\perp)$: $I\cap J=0$, $J=\Ann_\T I$, $I=\Ann_\T J$, $\T/(I+J)\simeq\Z/r$, $r\Z=\phi_f(J)$, $[S_2(\Z):\Z f\oplus L]=r$, and the congruence moduli
are the divisors of $r$; an eigenform $g\in f^\perp$ of all of $\T$ with $a_n(g)\equiv a_n(f)\pmod\lambda$ for every $n$ forces $\ell\mid r$; the same over $\Zl$. \emph{(proved here, given
\Cref{ch:hecke})} \Cref{lem:modgalrep:crit}.
\item \textbf{$U_p$ on oldforms.} On $\langle h(p^jz)\rangle$: eigenvalues $\alpha_p,\beta_p$ (and $0$ if $e_p\ge2$) for $p\nmid M$, $a_p(h)$ and $0$ for $p\,\|\,M$, $0$ for $p^2\mid M$, with explicit
eigenvectors; any choice gives a normalised eigenform of all of $\T$ in $V_h$. \emph{(proved here)} \Cref{lem:modgalrep:Up-oldspace}.
\item \textbf{Lemma lem:old.} $\ell$ odd, $\ell\nmid N$, $\bar\rho_{E,\ell}$ irreducible; a newform $h\ne f$ of level $M\mid N$ with $a_p(h)\equiv a_p(f)$ for $p\nmid\ell N$, and $p\nmid M$ at every $p\,\|\,N$ where $\bar\rho$ is
unramified: $V_h$ contains an eigenform of all of $\T$ congruent to $f$ at every $n$, so $\ell\mid r$. \emph{(proved here, given \Cref{thm:modgalrep:local-types,thm:modgalrep:deligne-fontaine})}
\Cref{lem:modgalrep:old}.
\item \textbf{Theorem thm:qual.} $\ell$ odd, $\ell\nmid N$, $\bar\rho_{E,\ell}$ irreducible (and not $N_s(\F_3)$ if $\ell=3$), and at $q\ne\ell$ type (a) $\kod{IV},\kod{IV}^*$, $q\ge5$ or $q=2$, $\ell=3$; (b)
$\kod{II},\kod{II}^*$, $q\ge5$, $\ell=3$; or (c) $\kod{I}_n^*$, $\ell\mid n$, $q$ odd: then $\ell\mid\deg\varphi_{E_0}$. \emph{(proved here, given the stated theorems
\Cref{thm:modgalrep:carayol,thm:modgalrep:deligne-fontaine,thm:modgalrep:optimal-level,thm:modgalrep:ars})} \Cref{thm:modgalrep:qual}.
\item \textbf{The twist map.} $\iota_S\colon h\mapsto h\otimes\chi_S$ twists generalised eigensystems by $\chi_S$, maps $g^\perp$ into $f^\perp$ when $f=g^{\chi_S}$, and $\iota(g)=(g^{\chi_S})^S=\iota(g^S)$. \emph{(proved
here, given \Cref{ch:hecke})} \Cref{prop:modgalrep:twist-map}.
\item \textbf{Levels of $E_S$ and $g^S$.} $N_S=N/D$, $N/D^2$, $N$ for types $\kod{I}_n^*$, $\kod{I}_0^*$, potentially good at $S$; $\iota(f_{E_S})=f_E$; $g^S$ lies at level $N_S\prod q^{\delta_q}=N$.
\emph{(proved here)} \Cref{prop:modgalrep:twist-levels}.
\item \textbf{Theorem thm:twist.} If $q^2\mid N$ for $q\in S$ and $\lcm(N_S,D^2)\mid N$, then $r_{E_S}\mid r_E$, with equality if also $\lcm(N,D^2)\mid N_S$. \emph{(stated; proved in \Cref{ch:congruence})}
\Cref{thm:modgalrep:twist-monotonicity}.
\item \textbf{Theorem thm:quant.} $\ell\nmid2N$, $E[\ell]$ irreducible, $S$ the odd $\kod{I}_0^*$ and $\kod{I}_n^*$ primes: $v_\ell(\degphi)\ge v_\ell(r_{E_S})+\sum_{q\in S}e_\ell(E,q)$, with equality if $\ell\notin S_f$.
\emph{(stated; proved in \Cref{ch:congruence})} \Cref{thm:modgalrep:quant}.
\item \textbf{Tamagawa exponent.} $t_q$ is the largest $t$ with $\rho\bmod\lambda^t$ unramified at $q$. \emph{(definition)} \Cref{def:modgalrep:tamagawa-exponent}.
\item \textbf{$t_q$ for $f_E$.} At $q\,\|\,N$, $t_q=v_\ell(n_q)=t_\ell(E,q)$. \emph{(proved in \Cref{ch:galrep}; recalled)} \Cref{prop:modgalrep:tamagawa-exponent}.
\item \textbf{Pollack and Weston.} Moving $q$ into the new part changes the congruence ideal by $\lambda^{t_q}$ (weight $2$, square-free level, $\bar\rho$ ramified at two primes or more).
\emph{(stated as the paper summarises it; reference)} \Cref{thm:modgalrep:pollack-weston}.
\item \textbf{Kim and Ota.} $\ord_\lambda\eta_f(N)=\ord_\lambda\eta_f(N^+,N^-)+\sum_{q\mid N^-}t_f(q)$ under their hypotheses \emph{(stated; reference)}; hence $v_\ell(r_E)\ge\sum_{q\mid N^-}t_\ell(E,q)$
\emph{(proved here)}. \Cref{thm:modgalrep:kim-ota}.
\end{enumerate}

\section*{Where the paper uses this chapter}

\begin{itemize}
\item \textbf{Section 1.1} ($V$, $T$, $\rho$; $\delta_p=2,1,0$; naive and true factors): \Cref{thm:modgalrep:comparison-E,thm:modgalrep:local-types,cor:modgalrep:delta-naive}.
\item \textbf{Section 1.5} (history; Theorem thm:qual; Theorem thm:quant; Kim and Ota; ``by level lowering and twisting rather than by the Taylor-Wiles method''):
\Cref{thm:modgalrep:qual,thm:modgalrep:quant,thm:modgalrep:kim-ota,thm:modgalrep:pollack-weston}, \Cref{rem:modgalrep:ribet-takahashi,rem:modgalrep:taylor-wiles}.
\item \textbf{Lemma lem:twist(iii)} (Atkin and Li): \Cref{cor:modgalrep:atkin-li}, from \Cref{prop:modgalrep:twist-conductor} and \Cref{thm:modgalrep:carayol}.
\item \textbf{Lemma lem:euler} (``The local component of $f$ is an unramified twist of the Steinberg representation''; ``$V=V'\otimes\chi_{q^*}$ with $V'$ Steinberg''):
\Cref{prop:modgalrep:kodaira-dictionary,thm:modgalrep:local-types}.
\item \textbf{Section 3, Lemma lem:crit} (congruence criterion): \Cref{lem:modgalrep:crit}, with \Cref{def:modgalrep:congruence-number,thm:modgalrep:ars}.
\item \textbf{Section 3, Lemma lem:old} (oldforms): \Cref{lem:modgalrep:old,lem:modgalrep:Up-oldspace}, using \Cref{cor:modgalrep:congruent-newforms} ($p\nmid\ell N$),
\Cref{thm:modgalrep:deligne-fontaine} ($p=\ell$), \Cref{thm:modgalrep:local-types} ($p^2\mid N$; $p\,\|\,N$ via Theorem 3.1(e) of \cite{DDT}).
\item \textbf{Section 3, the optimal-level paragraph}: \Cref{thm:modgalrep:optimal-level,cor:modgalrep:optimal-level-E}, \Cref{thm:modgalrep:crystalline} (``good''),
\Cref{prop:modgalrep:serre-conductor-divides} ($N(\bar\rho)\mid N$), \Cref{def:modgalrep:serre-invariants}.
\item \textbf{Section 3, cases (a), (b), (c)}: \Cref{thm:modgalrep:qual} and its proof (\Cref{sec:modgalrep:proof-qual}), with \Cref{thm:modgalrep:ribet} in case (c) and
\Cref{prop:modgalrep:twists}.
\item \textbf{Section 3, the Example and the remark after it} (\curve{52a1}, \curve{121a1}, \curve{539c1}, \curve{980i1}; 99A1): \Cref{ex:modgalrep:paper-example,rem:modgalrep:l-divides-N,ex:modgalrep:level-lowering}.
\item \textbf{Section 4.1} ($M^\Sigma$, $f^\Sigma$) and \textbf{Theorem thm:dfg14} with the paragraph after it (Wiles's level-raising factor):
\Cref{thm:modgalrep:dfg14,prop:modgalrep:level-raising-factor,thm:modgalrep:level-raising}.
\item \textbf{Section 4.2} (the lattice $M_{f,\lambda}$ of DFG04 is $H^1(E,\Z_\ell)$; $A=\ad T$): \Cref{thm:modgalrep:comparison-E}.
\item \textbf{Theorem thm:twist} and the data paragraph after it: \Cref{prop:modgalrep:twist-map,prop:modgalrep:twist-levels,thm:modgalrep:twist-monotonicity}.
\item \textbf{Theorem thm:quant}, its proof (Steps 1 and 3) and Remark rem:afterquant: \Cref{prop:modgalrep:twist-levels}(iii), \Cref{lem:modgalrep:Up-oldspace},
\Cref{thm:modgalrep:quant}, \Cref{rem:modgalrep:quant-level-raising}.
\item \textbf{Remark rem:ko(2)}: \Cref{thm:modgalrep:kim-ota}, \Cref{prop:modgalrep:tamagawa-exponent}.
\item \textbf{Theorem thm:dfg37} (the hypothesis $\ell\notin S_f$): \Cref{rem:modgalrep:taylor-wiles}.
\end{itemize}

\section{Conventions and recollections}\label{sec:modgalrep:conventions}

This section fixes the notation for newforms, Hecke algebras, Frobenius elements, characters and twists, and recalls from
\Cref{ch:hecke} the facts about newforms that every later proof uses. The one convention that matters for the paper is the distinction
between the naive twist $f\otimes\chi$ of a $q$-expansion and the newform twist $f^\chi$: the paper writes $f\otimes\chi$ for both, and its proof of
Theorem thm:qual is correct only when the two agree or when the difference is handled by strong multiplicity one.

\begin{convention}[newforms, Hecke algebras, Frobenius, characters]\label{conv:modgalrep:notation}
\begin{enumerate}[label=(\alph*)]
\item \emph{Newforms.} Unless said otherwise, $f=\sum_{n\ge1}a_n(f)q^n$ is a normalised newform of weight $2$ on $\Gamma_0(N)$ (trivial character); here $q=e^{2\pi
iz}$, and a prime is written $q$ only when no $q$-expansion is in the same formula. $K_f=\Q(a_n(f):n\ge1)$ is its coefficient field, a totally real
number field (\Cref{ch:hecke}), $\mathcal{O}_f$ its ring of integers, $\lambda\mid\ell$ a prime of $\mathcal{O}_f$, or of the ring of integers $\mathcal{O}$ of a finite
extension of $K_f$ (``after enlarging $\mathcal{O}$''), $K_{f,\lambda}$ and $\mathcal{O}_\lambda$ the completions, $k_\lambda=\mathcal{O}_\lambda/\lambda$ the residue field (\Cref{ch:ant}).
For an elliptic curve $E/\Q$ of conductor $N$, $f_E$ is its newform (modularity, \Cref{ch:moddeg}), and then $K_f=\Q$, $\lambda=\ell$.
\item \emph{Hecke algebras.} $\T=\T_N$ is the \emph{full} Hecke algebra of level $N$, the subring of $\End_\C S_2(\Gamma_0(N))$ generated by all $T_n$, $n\ge1$,
the operators $U_p=T_p$ for $p\mid N$ included; $\T_\ell=\T\otimes\Zl$. The \emph{eigencharacter} of $f$ is $\phi_f\colon\T\to\mathcal{O}_f$, $T_n\mapsto a_n(f)$ (upright
$\phi$; a curly $\varphi$ is always a map of curves, such as the modular parametrisation $\varphi_E$). $J_f=J_0(N)/I_fJ_0(N)$, $I_f=\ker\phi_f$, is
Shimura's abelian variety attached to $f$; the letter $A_f$ is reserved for the paper's adjoint $\ad V$.
\item \emph{Frobenius.} $\Frob_p$ is an \emph{arithmetic} Frobenius at $p$ and $\Frob_p^{-1}$ (or $\Phi$, at the prime of a local field) the geometric one.
$\chi_\ell(\Frob_p)=p$ for $p\ne\ell$. Local $L$-factors use the geometric Frobenius: $L(M,s)=\det(1-\Frob_p^{-1}p^{-s}\mid M^{I_p})^{-1}$ for a representation $M$ of
$G_{\Q_p}$ over a field of characteristic $0$, $p\ne\ell$ (\Cref{ch:local}).
\item \emph{Normalisation of $\rho_{f,\lambda}$.} The representation $\rho_{f,\lambda}$ on $V_\lambda(f)$ is normalised by $\det(X-\rho_{f,\lambda}(\Frob_p))=X^2-a_p(f)X+p$ for $p\nmid N\ell$
and $\det\rho_{f,\lambda}=\chi_\ell$, so that $\rho_{f_E,\ell}\simeq\rho_{E,\ell}$ on $V_\ell E$ (\Cref{thm:modgalrep:comparison-E}). The cohomological realisation
$H^1(X_0(N),\cdot)$ gives the dual $V_\lambda(f)^\vee\simeq V_\lambda(f)(-1)$, on which the \emph{geometric} Frobenius at $p\nmid N\ell$ has trace $a_p(f)$.
\item \emph{Characters.} A Dirichlet character $\chi$ is also the Galois character with $\chi(\Frob_p)=\chi(p)$ for $p\nmid\cond\chi$. $\chi_\ell$ is the $\ell$-adic
cyclotomic character and $\bar\chi_\ell$ its reduction; a letter subscript on $\chi$ always means a cyclotomic character. $\chi_d$ is the quadratic
character of $\Q(\sqrt d)$ for a fundamental discriminant $d$; for an odd prime $q$, $q^*=(-1)^{(q-1)/2}q$ and $\chi_{q^*}$ is the quadratic character of
conductor $q$ (so $\chi_{-3}$, $\chi_{-7}$, $\chi_{-11}$ appear below, never $\chi_3$, $\chi_7$, $\chi_{11}$). For a finite set $S$ of odd primes,
$\chi_S=\prod_{q\in S}\chi_{q^*}$, of conductor $D=\prod_{q\in S}q$.
\item \emph{Local components.} For a prime $p\ne\ell$, $\WD_p(f)$ is the Frobenius-semisimplification of the Weil-Deligne representation of
$\rho_{f,\lambda}|_{G_{\Q_p}}$ (\Cref{ch:local}); in that context $N$ also denotes the monodromy operator, and the context always makes clear which $N$
is meant. $\delta_p=\dim V_\lambda(f)^{I_p}$.
\item \emph{Labels.} Curve labels are Cremona's: \curve{11a1} is a curve, 11a an isogeny class, $f_{11a}$ its newform; the paper's (and ARS12's)
99A1 is \curve{99a1}. ``Optimal'' means $X_0(N)$-optimal, and $E_0$ is the optimal curve in the class of $E$.
\end{enumerate}
\end{convention}

\begin{convention}[naive twists, newform twists and depletions]\label{conv:modgalrep:twists}
Let $\chi$ be a quadratic Dirichlet character of conductor $r$, $\Sigma(r)$ the set of primes dividing $r$, and $g$ a newform of weight $2$ with
trivial character.
\begin{enumerate}[label=(\alph*)]
\item The \emph{naive twist} is the $q$-series $g\otimes\chi=\sum_n\chi(n)a_n(g)q^n$. This is the paper's $g\otimes\chi$ (Lemma lem:twist(iii), cases (b), (c) of the
proof of Theorem thm:qual) and its twist map $\iota$ (Theorems thm:twist, thm:quant).
\item The \emph{newform twist} $g^\chi$ is the unique newform of weight $2$ with trivial character such that $\rho_{g^\chi,\lambda}\simeq\rho_{g,\lambda}\otimes\chi$, equivalently
$a_p(g^\chi)=\chi(p)a_p(g)$ for almost all $p$ (\Cref{prop:modgalrep:twists}).
\item For a set $\Sigma$ of primes the \emph{depletion} of a $q$-series $h$ is $h^\Sigma=\sum_{(n,\Sigma)=1}a_n(h)q^n$ (the paper's $f^\Sigma$), and $h^{\{q\}}$ for one prime.
A parenthesised superscript never means a depletion: $E^{(q)}$ is the twist of $E$ by $\Q(\sqrt{q^*})$.
\end{enumerate}
Then (\Cref{prop:modgalrep:twists}(ii))
\[
g\otimes\chi=(g^\chi)^{\Sigma(r)},\qquad g\otimes\chi=g^\chi\iff a_q(g^\chi)=0\ \text{for every } q\mid r,\qquad (g\otimes\chi)\otimes\chi=g^{\Sigma(r)} .
\]
The last identity holds because $\chi(n)^2=1$ for $(n,r)=1$ and both sides vanish at $n$ with $(n,r)>1$.
\end{convention}

\begin{example}[the twist dictionary on two curves]\label{ex:modgalrep:twist-dictionary}
(a) The twist of \curve{11a1} by $\Q(\sqrt{-11})$ is \curve{121d2}, so $f_{11a}^{\chi_{-11}}=f_{121d}$, and $a_{11}(f_{121d})=0$ (type $\kod{I}_5^*$ at $11$). Hence
$f_{11a}\otimes\chi_{-11}=f_{121d}$ exactly, as the paper says in its Example. (b) The twist of \curve{539c1} by $\Q(\sqrt{-7})$ is \curve{77c1}, so
$f_{539c}^{\chi_{-7}}=f_{77c}$. But $a_7(f_{77c})=-1$ (type $\kod{I}_3$ at $7$), so $f_{539c}\otimes\chi_{-7}=(f_{77c})^{\{7\}}$, the $7$-depletion of $f_{77c}$, which is not a newform:
it is an eigenform of level $539$ for the full Hecke algebra there, with $U_7$-eigenvalue $0$, in the span of $f_{77c}(z)$ and $f_{77c}(7z)$. Both
statements were checked with PARI/GP on the coefficients $a_n$, $n\le3000$ (for (b), $a_p(f_{539c})\chi_{-7}(p)=a_p(f_{77c})$ for all primes $p<400$, $p\ne7$).
\end{example}

\begin{proposition}[newforms; recollection]\label{prop:modgalrep:newforms-recap}
Let $f$ be as in \Cref{conv:modgalrep:notation}(a).
\begin{enumerate}[label=(\roman*)]
\item $\phi_f(T_n)=a_n(f)$ for all $n$; $a_{mn}=a_ma_n$ for $(m,n)=1$; $a_{p^{k+1}}=a_pa_{p^k}-p\,a_{p^{k-1}}$ for $p\nmid N$; and $a_{p^k}=a_p^k$ for $p\mid N$. In general $a_1(T_ng)=a_n(g)$ for
every $g\in S_2(\Gamma_0(N))$, so an eigenform of all of $\T$ has $a_1\ne0$ and can be normalised.
\item $a_p(f)\in\{\pm1\}$ if $p\,\|\,N$, and $a_p(f)=0$ if $p^2\mid N$ (Atkin and Lehner \cite[Theorem 3]{AtkinLehner70}; reproved from the Galois side in
\Cref{thm:modgalrep:local-types}).
\item (Strong multiplicity one.) Two newforms of weight $2$ with trivial character, of any levels, with $a_p$ equal for all but finitely many $p$ are
equal; in particular they have the same level. (\Cref{cor:modgalrep:congruent-newforms}(iv) relaxes ``all but finitely many'' to ``a set of
density one''.)
\item $S_2(\Gamma_0(N))=\bigoplus_{M\mid N}\bigoplus_h V_h$, where $h$ runs over the newforms of level $M$ and $V_h$ is the span of the $h(dz)$, $d\mid N/M$, which are
linearly independent. Each $V_h$ is $\T$-stable, $T_n$ acts on $V_h$ by the scalar $a_n(h)$ for $(n,N)=1$, and the $V_h$ are mutually orthogonal for the
Petersson product. Consequently the generalised eigenspace in $S_2(\Gamma_0(N))$ of the operators $T_p$, $p\nmid N$, for the eigenvalues $(a_p(h))_{p\nmid N}$ is
$V_h$.
\item $L(f,s)=\sum a_n(f)n^{-s}=\prod_{p\nmid N}(1-a_pp^{-s}+p^{1-2s})^{-1}\prod_{p\mid N}(1-a_pp^{-s})^{-1}$.
\end{enumerate}
\end{proposition}

\begin{proof}
All of this is proved in \Cref{ch:hecke} (Atkin-Lehner-Li theory, the Petersson product, the Hecke relations), except (ii) and (iii), which
are quoted there from \cite{AtkinLehner70} and \cite[Chapter 5]{DiamondShurman}. The last sentence of (iv) follows from the rest of (iv) and from (iii):
on $V_{h'}$ the $T_p$, $p\nmid N$, act by the scalars $a_p(h')$, and by (iii) these agree with the $a_p(h)$ for all $p\nmid N$ only if $h'=h$.
\end{proof}

\begin{paperlink}
Lemma lem:old uses (iv) in its first sentence (``The forms $h(dz)$, $d\mid N/M$, span a $\T$-stable subspace $V$ of $S_2(\Gamma_0(N))$ orthogonal to the newform
$f$''), and (i) and (ii) in its last paragraph and in the case $p^2\mid N$. Lemma lem:crit uses (iii) and (iv) for the decomposition $\T\otimes\Q=\Q\times\T'$
(\Cref{lem:modgalrep:crit}, Step 1), and Theorem thm:quant uses (iii) and (iv) in Step 1.
\end{paperlink}

\section{The Galois representation of a newform}\label{sec:modgalrep:galrep}

This section attaches to a newform $f$ of weight $2$ and a prime $\lambda$ of its coefficient field a two-dimensional $\lambda$-adic representation
$\rho_{f,\lambda}$ of $G_\Q$, its lattices and its reduction $\bar\rho_{f,\lambda}$, and compares it with $\rho_{E,\ell}$ when $f=f_E$. The paper uses these objects
everywhere in Section 3: the residual representation $\bar\rho_h$ of a newform $h$ in Lemma lem:old, the phrase ``arises from a newform'' in the
optimal-level theorem, and ``$\bar\rho_g\simeq\bar\rho$, that is, $a_p(g)\equiv a_p(f)\pmod\lambda$'' in cases (a), (b), (c); and in Section 4.2 the lattice
$M_{f,\lambda}$ of Diamond, Flach and Guo is compared with $T_\ell E$. \Cref{ch:galois} used the representations of newforms as a black box
(its corollary on congruent traces, part (d)); \Cref{thm:modgalrep:eichler-shimura} discharges it.

\begin{theorem}[Eichler, Shimura; Deligne and Rapoport]\label{thm:modgalrep:eichler-shimura}
Let $f$ be a newform of weight $2$ on $\Gamma_0(N)$, $I_f=\ker\phi_f$, $J_f=J_0(N)/I_fJ_0(N)$, $d=[K_f:\Q]$.
\begin{enumerate}[label=(\roman*)]
\item $J_f$ is an abelian variety over $\Q$ of dimension $d$, with good reduction at every $p\nmid N$, and $\T/I_f$, an order in $K_f$, acts on it by
endomorphisms defined over $\Q$.
\item For every prime $\ell$, $V_\ell(J_f)=T_\ell(J_f)\otimes\Ql$ is free of rank $2$ over $K_f\otimes\Ql=\prod_{\lambda\mid\ell}K_{f,\lambda}$, and $G_\Q$ acts $K_f\otimes\Ql$-linearly. For
$\lambda\mid\ell$ put $V_\lambda(f)=V_\ell(J_f)\otimes_{K_f\otimes\Ql}K_{f,\lambda}$, a two-dimensional $K_{f,\lambda}$-space, and let $\rho_{f,\lambda}\colon G_\Q\to\GL(V_\lambda(f))\simeq\GL_2(K_{f,\lambda})$ be the
action. Then:
\begin{enumerate}[label=(\alph*)]
\item $\rho_{f,\lambda}$ is continuous and unramified at every $p\nmid N\ell$;
\item for $p\nmid N\ell$, $\det(X-\rho_{f,\lambda}(\Frob_p))=X^2-a_p(f)X+p$;
\item $\det\rho_{f,\lambda}=\chi_\ell$;
\item $\rho_{f,\lambda}$ is odd: $\det\rho_{f,\lambda}(c)=-1$ for every complex conjugation $c$;
\item $\rho_{f,\lambda}$ is absolutely irreducible (\Cref{prop:modgalrep:irreducible}).
\end{enumerate}
\item The $I_f$-torsion $H^1_{\mathrm{et}}(X_0(N)_{\Qbar},\Ql)[I_f]$ is free of rank $2$ over $K_f\otimes\Ql$, and its $\lambda$-component is $V_\lambda(f)^\vee\simeq V_\lambda(f)(-1)$; the
geometric Frobenius at $p\nmid N\ell$ has trace $a_p(f)$ on it.
\end{enumerate}
\end{theorem}

\begin{proof}[Sketch of proof]
(i) $J_0(N)$ has good reduction at every $p\nmid N$ (Igusa; Deligne and Rapoport \cite{DeligneRapoport73}), hence so has its quotient $J_f$. The
cotangent space of $J_f$ is the part of $S_2(\Gamma_0(N))$ killed by $I_f$, which is spanned by the $d$ Galois conjugates of $f$ (a newform occurs with
multiplicity one, \Cref{prop:modgalrep:newforms-recap}(iv)), so $\dim J_f=d$. These facts, and the action of $\T$ on $J_0(N)$ by endomorphisms over
$\Q$, are in \Cref{ch:moddeg}.

(ii) $H_1(J_f(\C),\Q)$ is a vector space over the field $K_f=(\T/I_f)\otimes\Q$ of $\Q$-dimension $2d$, so it is free of rank $2$ over $K_f$; tensoring with $\Ql$
and using $V_\ell(J_f)\simeq H_1(J_f(\C),\Q)\otimes\Ql$ ($\T$-equivariantly) gives the freeness, and the Galois action commutes with $\T$ because the Hecke
correspondences are defined over $\Q$. Continuity is that of the action on Tate modules (\Cref{ch:galois}). (a) For $p\nmid N\ell$, $J_f$ has good reduction
at $p$, so inertia at $p$ acts trivially on $V_\ell(J_f)$ (the criterion of N\'eron, Ogg and Shafarevich for abelian varieties, due to Serre and Tate
\cite{SerreTate68}; \Cref{ch:galrep} proves the case of elliptic curves). (c) The Weil pairing attached to a polarisation of $J_f$ for which
the Hecke operators are self-adjoint up to the Atkin-Lehner involution gives a $K_f\otimes\Ql$-bilinear alternating perfect pairing $V_\ell(J_f)\times
V_\ell(J_f)\to K_f\otimes\Ql(1)$ (Shimura \cite[Chapter 7]{Shimura71}), so $\wedge^2_{K_{f,\lambda}}V_\lambda(f)\simeq K_{f,\lambda}(1)$ and $\det\rho_{f,\lambda}=\chi_\ell$. (b) Reduction modulo $p$
identifies $V_\ell(J_f)$ with $V_\ell((J_f)_{\F_p})$, $\Frob_p$ acting as the Frobenius endomorphism $\pi_p$. The Eichler-Shimura congruence relation
$(T_p)_*=\pi_p+\hat\pi_p$ with $\pi_p\hat\pi_p=p$ (\Cref{ch:moddeg}) gives $A^2-a_pA+p=0$ for $A=\rho_{f,\lambda}(\Frob_p)$, since $T_p$ acts on $V_\lambda(f)$ by $a_p=a_p(f)$. If $A$ is
not a scalar, its characteristic polynomial is its minimal polynomial $X^2-a_pX+p$. If $A=c$ is a scalar, then $c^2=\det A=p$ by (c), and $c^2-a_pc+p=0$
gives $a_p=2c=\tr A$; so again $\det(X-A)=X^2-\tr A\,X+\det A=X^2-a_pX+p$. (d) $\chi_\ell(c)=-1$, as $c$ inverts roots of unity. (iii) The \'etale cohomology
$H^1_{\mathrm{et}}(X_0(N)_{\Qbar},\Zl)$ is $\Hom(T_\ell J_0(N),\Zl)\simeq T_\ell J_0(N)(-1)$, Hecke and Galois equivariantly (\Cref{ch:curves}); its $I_f$-torsion is
dual to the $I_f$-coinvariants $V_\ell(J_0(N))/I_f$, which is $V_\ell(J_f)$; and $V^\vee\simeq V\otimes\det^{-1}=V(-1)$ for two-dimensional $V$ with $\det=\chi_\ell$. An
eigenvalue $\alpha$ of the arithmetic Frobenius on $V$ is the eigenvalue $\alpha$ of the geometric Frobenius on $V^\vee$, which gives the last claim.
For details see Shimura \cite[Chapter 7]{Shimura71} (Theorem 7.14 in the edition we know; an unverified pointer), Diamond and Shurman
\cite[Chapter 9]{DiamondShurman}, Diamond and Im \cite{DiamondIm} and \cite[Theorem 3.1]{DDT}; the general case is Deligne's \cite{Deligne}.
\end{proof}

\begin{proposition}[irreducibility]\label{prop:modgalrep:irreducible}
For every newform $f$ of weight $2$ with trivial character and every $\lambda$, $\rho_{f,\lambda}$ is absolutely irreducible.
\end{proposition}

\begin{proof}[Sketch of proof]
Suppose that $\rho_{f,\lambda}\otimes\overline K_{f,\lambda}$ has a stable line, with character $\psi_1$ and quotient character $\psi_2$; both are continuous characters
$G_\Q\to\overline\Q_\ell^\times$ unramified outside $N\ell$, with $\psi_1\psi_2=\chi_\ell$. They are subquotients of $V_\ell(J_f)$, which is Hodge-Tate with weights $0$
and $1$ (Faltings; Tate for abelian varieties with good reduction), so they are Hodge-Tate, and an abelian Hodge-Tate character of $G_\Q$ is
of the form $\varepsilon\chi_\ell^k$ with $\varepsilon$ of finite order and $k\in\Z$ (Serre \cite{SerreAbelian}). As $\psi_1\psi_2=\chi_\ell$, the exponents are $\{0,1\}$:
$\{\psi_1,\psi_2\}=\{\varepsilon_1,\varepsilon_2\chi_\ell\}$. Then $a_p(f)=\varepsilon_1(p)+\varepsilon_2(p)p$ for all $p\nmid N\ell$ outside the conductors of the $\varepsilon_i$, an identity between
algebraic numbers, so under every complex embedding $\abs{a_p(f)}\ge p-1>2\sqrt p$ for $p\ge7$. This contradicts the bound $\abs{a_p(f)}\le2\sqrt p$ of
Eichler, Shimura and Weil (\Cref{ch:modforms}). For the general case (any weight and character) see Ribet \cite{Ribet77}.
\end{proof}

\begin{theorem}[Deligne; Deligne and Serre]\label{thm:modgalrep:deligne}
Let $f$ be a normalised newform of weight $k\ge2$, level $N$ and character $\varepsilon$, and $\lambda\mid\ell$ a prime of $K_f$. There is a continuous
irreducible representation $\rho_{f,\lambda}\colon G_\Q\to\GL_2(K_{f,\lambda})$, unramified outside $N\ell$, with $\det(X-\rho_{f,\lambda}(\Frob_p))=X^2-a_p(f)X+\varepsilon(p)p^{k-1}$
for $p\nmid N\ell$ and $\det\rho_{f,\lambda}=\varepsilon\chi_\ell^{k-1}$; it is odd. In weight $1$ Deligne and Serre attach an Artin representation with the same
properties ($k=1$).
\end{theorem}

We do not prove this; see Deligne \cite{Deligne} and \cite[Chapter 9]{DiamondShurman}. Only weight $2$ and trivial character occur in the paper's
proofs; the optimal-level theorem (\Cref{thm:modgalrep:optimal-level}) produces a priori a newform with a character, and the case of
characters is needed only to see that this character is trivial (\Cref{cor:modgalrep:optimal-level-E}).

\begin{proposition}[lattices and reduction]\label{prop:modgalrep:lattices}
Let $f$ be a newform of weight $2$ with trivial character and $\lambda\mid\ell$.
\begin{enumerate}[label=(\roman*)]
\item $\rho_{f,\lambda}$ stabilises an $\mathcal{O}_\lambda$-lattice $L\subset V_\lambda(f)$. The reduction $\bar\rho_L\colon G_\Q\to\GL(L/\lambda L)\simeq\GL_2(k_\lambda)$ has a
semisimplification $\bar\rho_{f,\lambda}^{\mathrm{ss}}$ independent of $L$, and $\tr\bar\rho_L(\Frob_p)\equiv a_p(f)$, $\det\bar\rho_L(\Frob_p)\equiv p\pmod\lambda$ for $p\nmid N\ell$.
\item $\bar\rho_L$ is odd, and for $\ell$ odd it is irreducible over $k_\lambda$ if and only if it is absolutely irreducible.
\item If $\bar\rho_L$ is irreducible, every $G_\Q$-stable $\mathcal{O}_\lambda$-lattice in $V_\lambda(f)$ is homothetic to $L$, so $\bar\rho_{f,\lambda}:=\bar\rho_L$ is well defined
up to isomorphism.
\end{enumerate}
\end{proposition}

\begin{proof}
(i) $G_\Q$ is compact and acts continuously, so the $\mathcal{O}_\lambda$-span of the orbit of any lattice is a stable lattice (\Cref{ch:galois}, the proposition
on continuous representations). The independence of the semisimplification is the theorem of Brauer and Nesbitt in the form of
\Cref{ch:reps}: two stable lattices give reductions with the same characteristic polynomials at every element. The congruences follow from
\Cref{thm:modgalrep:eichler-shimura}(ii)(b). (ii) $\bar\rho_L(c)$ has order dividing $2$ and determinant $-1$, so for $\ell$ odd its eigenvalues are $1$
and $-1$, which are distinct in $k_\lambda$. A line over $\bar k_\lambda$ stable under $\bar\rho_L(G_\Q)$ is stable under $\bar\rho_L(c)$, hence is one of the two
eigenlines of $\bar\rho_L(c)$, which are defined over $k_\lambda$; it is then a $k_\lambda$-rational stable line. (iii) This is the homothety lemma of
\Cref{ch:reps} (if $L'$ is stable, scale it so that $L'\subseteq L$, $L'\not\subseteq\lambda L$; its image in $L/\lambda L$ is a nonzero stable subspace, hence all of $L/\lambda L$,
and Nakayama's lemma gives $L'=L$).
\end{proof}

\begin{corollary}[congruent newforms]\label{cor:modgalrep:congruent-newforms}
Let $f$, $h$ be newforms of weight $2$ with trivial character, of levels $N_f$, $N_h$, with coefficients in the ring of integers $\mathcal{O}$ of a number field
$K$, and $\lambda\mid\ell$ a prime of $\mathcal{O}$; write $\bar\rho_f$, $\bar\rho_h$ for reductions over $k_\lambda$ of lattices in $\rho_{f,\lambda}$, $\rho_{h,\lambda}$ (formed with the
completion of $K$ at $\lambda$).
\begin{enumerate}[label=(\roman*)]
\item If $a_p(f)\equiv a_p(h)\pmod\lambda$ for all $p$ in a set of primes of density one, then $\bar\rho_f^{\mathrm{ss}}\simeq\bar\rho_h^{\mathrm{ss}}$.
\item If moreover $\ell$ is odd and $\bar\rho_f$ is irreducible, then $\bar\rho_h\simeq\bar\rho_f$ for every choice of lattices.
\item Conversely, if $\bar\rho_f^{\mathrm{ss}}\simeq\bar\rho_h^{\mathrm{ss}}$, then $a_p(f)\equiv a_p(h)\pmod\lambda$ for every $p\nmid N_fN_h\ell$.
\item If $a_p(f)=a_p(h)$ for all $p$ in a set of density one, then $f=h$.
\end{enumerate}
\end{corollary}

\begin{proof}
Determinants agree: both are $\chi_\ell$ (\Cref{thm:modgalrep:eichler-shimura}(ii)(c)). (i), (ii), (iii) are the corollary on congruent traces of
\Cref{ch:galois}, parts (a), (b), (c), applied to $\rho_1=\rho_{f,\lambda}$, $\rho_2=\rho_{h,\lambda}$ with $S$ the primes dividing $N_fN_h\ell$; in (ii), $\bar\rho_f$ is absolutely
irreducible by \Cref{prop:modgalrep:lattices}(ii), which is what that corollary needs. (iv) By the same corollary in characteristic $0$ (Chebotarev
and Brauer-Nesbitt), the semisimple representations $\rho_{f,\lambda}$, $\rho_{h,\lambda}$ (irreducible, \Cref{prop:modgalrep:irreducible}) are isomorphic, so $a_p(f)=a_p(h)$
for all $p\nmid N_fN_h\ell$, and strong multiplicity one (\Cref{prop:modgalrep:newforms-recap}(iii)) gives $f=h$.
\end{proof}

\begin{paperlink}
Lemma lem:old, case $p\nmid\ell N$: ``Consequently $\bar\rho_h\simeq\bar\rho$ over a common finite extension of $\F_\ell$, by Chebotarev and Brauer-Nesbitt, $\bar\rho$
being irreducible, hence absolutely irreducible.'' This is \Cref{cor:modgalrep:congruent-newforms}(ii) with $f=f_E$ and $\bar\rho=\bar\rho_{E,\ell}$ (so
$\bar\rho_f\simeq\bar\rho\otimes_{\F_\ell}k_\lambda$ by \Cref{thm:modgalrep:comparison-E}); the common finite extension is $k_\lambda$. The irreducibility of $\bar\rho$ alone would
give only $\bar\rho_h^{\mathrm{ss}}\simeq\bar\rho^{\mathrm{ss}}$; absolute irreducibility (\Cref{prop:modgalrep:lattices}(ii)) is what identifies $\bar\rho_h$ itself. The proof of Theorem
thm:qual uses (i) to (iii) in the form ``$\bar\rho_g\simeq\bar\rho$, that is, $a_p(g)\equiv a_p(f)\pmod\lambda$ for all $p\nmid 3N$''.
\end{paperlink}

\begin{theorem}[$\rho_{f_E,\ell}$ and $\rho_{E,\ell}$]\label{thm:modgalrep:comparison-E}
Let $E/\Q$ have conductor $N$, $f=f_E$, $E_0$ the optimal curve of the isogeny class, $\varphi_{E_0}\colon X_0(N)\to E_0$ the modular parametrisation, and
$\ell$ a prime.
\begin{enumerate}[label=(\roman*)]
\item $J_f=E_0$, so $T_\ell E_0$ is a $G_\Q$-stable $\Zl$-lattice in $V_\ell(f)$ and $\rho_{E_0,\ell}=\rho_{f,\ell}$ on it. For every $E$ in the class, $V_\ell E\simeq V_\ell E_0$ through
an isogeny, so $\rho_{E,\ell}\otimes\Ql\simeq\rho_{f,\ell}$.
\item The lattice $M_{f,\lambda}:=H^1(X_0(N),\Zl)[\ker\phi_f]$ is equal to the image $\varphi_{E_0}^*H^1(E_0,\Zl)$, and $\varphi_{E_0}^*$ is a $G_\Q$-isomorphism from
$H^1(E_0,\Zl)=\Hom_{\Zl}(T_\ell E_0,\Zl)$, which is isomorphic to $T_\ell E_0(-1)$, onto it. No irreducibility is needed.
\item If $\bar\rho_{E,\ell}$ is irreducible, every $G_\Q$-stable lattice in $V_\ell(E)^\vee$, in particular $M_{f,\lambda}$ and $(T_\ell E)^\vee$ for every $E$ in the class,
is homothetic to $(T_\ell E)^\vee\simeq T_\ell E(-1)$.
\item In all cases $\ad M_{f,\lambda}\simeq\ad T_\ell E_0$ as $\Zl[G_\Q]$-modules, and $\ad M_{f,\lambda}\simeq\ad T_\ell E$ for every $E$ in the class when $\bar\rho_{E,\ell}$ is
irreducible.
\end{enumerate}
\end{theorem}

\begin{proof}
(i) That $J_f$ is the optimal quotient $E_0$ of $J_0(N)$ is \Cref{ch:moddeg} (modularity and the definition of optimality: the kernel of
$J_0(N)\to E_0$ is $I_fJ_0(N)$, connected). Then $V_\ell(f)=V_\ell(J_f)=V_\ell(E_0)$ with $K_f=\Q$. An isogeny $E\to E_0$ over $\Q$ induces a $G_\Q$-isomorphism $V_\ell E\to V_\ell E_0$.
(ii) This is proved in \Cref{ch:moddeg} (the proposition identifying $M_{f,\lambda}$ with the pull-back, which is the paper's Proposition prop:pullback).
In outline: the kernel of $J_0(N)\to E_0$ is connected, so $\varphi_{E_0*}\colon H_1(X_0(N),\Z)\to H_1(E_0,\Z)$ is surjective and $\varphi_{E_0}^*$ is injective with
torsion-free cokernel; $\T$ acts on its image through $\phi_f$; the image and $M_{f,\lambda}$ are saturated of the same rank $2$ in the same plane (the
$\phi_f$-eigenspace of $H^1(X_0(N),\Q)$ is two-dimensional, as a newform occurs once); so they are equal. Galois equivariance comes from the
\'etale comparison (\Cref{ch:curves}). (iii) $V_\ell(E)^\vee$ reduces to $\bar\rho_{E,\ell}^\vee\simeq\bar\rho_{E,\ell}\otimes\bar\chi_\ell^{-1}$, irreducible; apply
\Cref{prop:modgalrep:lattices}(iii) (in the form of \Cref{ch:reps} for any irreducible residual representation). (iv) $\ad(L^\vee)\simeq\ad L$ by
$x\mapsto-x^{\mathsf T}$, and $\ad(L(n))=\ad L$, $\ad(cL)=\ad L$ for a scalar $c$ (\Cref{ch:reps}); apply (ii), and (iii) with the isomorphism
$T_\ell E\simeq T_\ell E_0$ given by an isogeny of degree prime to $\ell$ when $E[\ell]$ is irreducible (\Cref{ch:galrep}).
\end{proof}

\begin{paperlink}
Section 4.2 of the revised paper, in the paragraph after Proposition prop:pullback: ``The lattice $M_{f,\lambda}$ of \cite{DFG04} is therefore
$H^1(E,\Z_\ell)=\operatorname{Hom}(T_\ell E,\Z_\ell)\simeq T_\ell E(-1)$, by Galois-equivariance of $\varphi^*$ in \'etale cohomology, and its adjoint is $\ad T_\ell E$, the adjoint
being insensitive to dualising and twisting.'' This is \Cref{thm:modgalrep:comparison-E}(ii) and (iv) for the optimal curve $E=E_0$. The original
version of the paper said instead that $M_{f,\lambda}$ and $T_\ell E$ are homothetic; by (iii) that is right only after replacing $T_\ell E$ by its dual
$T_\ell E(-1)$, and only when $\bar\rho$ is irreducible. Only the adjoint enters the paper's arguments (Theorem thm:main, with $A=\ad T$), and it is the same in
every reading.
\end{paperlink}

\section[Local-global compatibility]{Local-global compatibility and local components}\label{sec:modgalrep:local}

This section describes $\rho_{f,\lambda}|_{G_{\Q_p}}$ for every prime $p\ne\ell$. The paper uses this description in four places: $\delta_p=\dim V^{I_p}=2,1,0$ in Section
1.1; the phrases ``The local component of $f$ is an unramified twist of the Steinberg representation'' and ``$V=V'\otimes\chi_{q^*}$ with $V'$ Steinberg at $q$''
in the proof of Lemma lem:euler; the proof of Lemma lem:twist(iii) (Atkin and Li) by the conductor of a twisted Steinberg representation; and the
shape of $\bar\rho|_{D_p}$ at a prime $p\,\|\,N$ in Lemma lem:old (Theorem 3.1(e) of \cite{DDT}). \Cref{ch:local} has done all the Weil-Deligne
computations; this section adds the automorphic names, the theorem of Carayol that connects the two sides, and the consequences for the
levels of twists.

\begin{remark}[recollection from \Cref{ch:local}]\label{rem:modgalrep:wd-recap}
Let $K=\Q_p$, $p\ne\ell$, and let $V$ be a continuous representation of $G_K$ on a finite-dimensional space over a finite extension of $\Ql$. Then:
$V$ has a Weil-Deligne representation $\WD(V)=(r,N)$, with $r(I_K)$ finite and $N$ nilpotent, $r(w)Nr(w)^{-1}=\norm wN$, $\norm{\Frob}=p$, and $V\mapsto\WD(V)$ is an
equivalence compatible with sums, tensor products, duals and twists; $V^{I_K}=(\ker N)^{r(I_K)}$; $L(V,s)=\det(1-\Frob^{-1}p^{-s}\mid V^{I_K})^{-1}$; the conductor
exponent is $a(V)=\dim V-\dim V^{I_K}+\operatorname{sw}(V)$, additive, invariant under unramified twists, zero exactly when $V$ is unramified, equal to
$\dim V-\dim V^{I_K}$ when $V$ is tame. The special representation $\Sp(2)$ has basis $e_0,e_1$ with $Ne_0=e_1$, $Ne_1=0$, $r(w)e_i=\norm w^ie_i$; a two-dimensional
$(r,N)$ with $N\ne0$ is $\Sp(2)\otimes\mu$ for a unique character $\mu$ (the character on $V/\ker N$); $a(\Sp(2)\otimes\mu)=1$ for $\mu$ unramified and $2a(\mu)$ for $\mu$
ramified; a two-dimensional Frobenius-semisimple $(r,0)$ is a sum of two characters or is irreducible, and for $p$ odd an irreducible one is
induced from a quadratic extension. The $\epsilon$-factors of Langlands and Deligne are defined there as well; the paper does not use them.
On the Galois side, $\Sp(2)\otimes\mu$ is the representation $\sm{\mu\chi_\ell}{*}{0}{\mu}$ of $G_K$ (after identifying characters of $W_K$ and of $G_K$ with values
in $\Ql^\times$); on $I_K$ it is $\mu$ times $\sm{1}{c\,t_\ell}{0}{1}$ with $c\ne0$ and $t_\ell$ the tame character, and it is non-split.
\end{remark}

\begin{theorem}[the local Langlands correspondence for $\GL_2(\Q_p)$]\label{thm:modgalrep:local-langlands}
Fix an isomorphism $\overline\Q_\ell\simeq\C$. There is a bijection $\pi\mapsto\operatorname{rec}(\pi)$ from isomorphism classes of irreducible admissible
infinite-dimensional representations of $\GL_2(\Q_p)$ to isomorphism classes of two-dimensional Frobenius-semisimple Weil-Deligne
representations of $W_{\Q_p}$, characterised by the compatibility with twists by characters (identified through local class field theory) and
the preservation of $L$- and $\epsilon$-factors of pairs. It preserves conductors, $a(\operatorname{rec}\pi)=a(\pi)$, where $p^{a(\pi)}$ is the conductor of Casselman's
newvector theory \cite{Casselman73}, and it matches the types: principal series $\pi(\mu_1,\mu_2)$ with sums of two characters, twists $\St\otimes\mu$ of
the Steinberg representation with $\Sp(2)\otimes\mu$, and supercuspidal representations with irreducible representations ($N=0$); in each case
up to one fixed unramified twist, so that ``$\mu$ unramified'' has the same meaning on both sides.
\end{theorem}

We do not prove this; see Bushnell and Henniart \cite{BushnellHenniart06}, and Rohrlich \cite{Rohrlich94} for the Weil-Deligne side of elliptic
curves. Two normalisations are in common use, differing by the twist by $\abs{\det}^{1/2}$, and some sources attach to a newform the dual of
our $\rho_{f,\lambda}$; we fix none of them, because everything this book uses (conductors, the type, $L$-factors) is invariant under unramified twists and
duality or is stated directly on the Galois side (\Cref{thm:modgalrep:carayol}(iii)). The automorphic names below are labels and enter no proof.

\begin{definition}[local components]\label{def:modgalrep:local-component}
A newform $f$ of weight $2$ on $\Gamma_0(N)$ generates a cuspidal automorphic representation $\pi_f=\bigotimes'_v\pi_v$ of $\GL_2(\A_\Q)$ (Jacquet and Langlands
\cite{JacquetLanglands70}). Its factor $\pi_p$ is the \emph{local component of $f$ at $p$}. It has trivial central character and conductor exponent
$a(\pi_p)=v_p(N)$ (Casselman \cite{Casselman73}), and the eigenvalue $a_p(f)$ is determined by $\pi_p$ (it is, after normalisation, the eigenvalue of the
Hecke operator at $p$ on the newvector). We say that $f$ has \emph{type} (U), (St), (RSt), (RPS) or (SC) at $p$ according to the rows of
\Cref{thm:modgalrep:local-types}; Loeffler and Weinstein \cite{LoefflerWeinstein12} explain how to compute $\pi_p$ in practice.
\end{definition}

\begin{theorem}[local-global compatibility at $p\ne\ell$; Langlands, Deligne, Carayol]\label{thm:modgalrep:carayol}
Let $f$ be a newform of weight $2$ on $\Gamma_0(N)$, $p$ a prime and $\lambda\mid\ell$ with $\ell\ne p$; let $V=V_\lambda(f)|_{G_{\Q_p}}$.
\begin{enumerate}[label=(\roman*)]
\item $\WD_p(f)$ corresponds to $\pi_p$ under \Cref{thm:modgalrep:local-langlands} (in the normalisation of the source). In particular $\WD_p(f)$, viewed
over $\C$, does not depend on $\ell\ne p$ or on $\lambda$, and its type is that of $\pi_p$.
\item The conductor exponent of $V$ is $a(V)=v_p(N)$. Hence the conductor of $\rho_{f,\lambda}$ is $N$ when $\ell\nmid N$, and its prime-to-$\ell$ part is that of
$N$ in general.
\item $L_p(f,s)=L(V,s-1)=\det(1-\Frob_p^{-1}p^{1-s}\mid V^{I_p})^{-1}$, where $L_p(f,s)$ is the Euler factor of \Cref{prop:modgalrep:newforms-recap}(v). Equivalently:
for $p\nmid N$ the arithmetic Frobenius has eigenvalues $\alpha$, $\beta$ on $V$ with $\alpha+\beta=a_p(f)$, $\alpha\beta=p$; for $p\mid N$, either $V^{I_p}=0$ and $a_p(f)=0$, or
$V^{I_p}$ is a line on which $\Frob_p$ acts by $p\,a_p(f)$ and $a_p(f)\ne0$.
\end{enumerate}
\end{theorem}

We do not prove (i), (ii), or (iii) at $p\mid N$; see Carayol \cite{Carayol86} (the theorem stated in its introduction), after Langlands \cite{Langlands73} and
Deligne \cite{Deligne73}, and \cite[Theorem 3.1]{DDT}; (iii) is the equality of $L$-factors under \Cref{thm:modgalrep:local-langlands}, read through (i). The two forms
of (iii) are equivalent by a one-line computation:
for $p\nmid N$, the geometric Frobenius has eigenvalues $\alpha^{-1}$, $\beta^{-1}$, and $(1-\alpha^{-1}p\cdot p^{-s})(1-\beta^{-1}p\cdot p^{-s})=(1-\beta p^{-s})(1-\alpha p^{-s})$ because $p/\alpha=\beta$; for a
line on which $\Frob_p$ acts by $c$, $1-c^{-1}p\cdot p^{-s}=1-a_pp^{-s}$ exactly when $c=p/a_p$, which is $p\,a_p$ when $a_p=\pm1$. For $p\nmid N\ell$, (iii) is
\Cref{thm:modgalrep:eichler-shimura}(ii)(b). On $V^\vee$, the cohomological realisation, (iii) reads: $a_p(f)$ is the trace of the geometric
Frobenius on $(V^\vee)^{I_p}$ when $\delta_p\ne1$, and its eigenvalue on that line when $\delta_p=1$.

\begin{theorem}[the local types of a newform of weight $2$]\label{thm:modgalrep:local-types}
Let $f$ be a newform of weight $2$ on $\Gamma_0(N)$, $p\ne\ell$, $V=V_\lambda(f)|_{G_{\Q_p}}$, $\WD_p(f)=(r,N)$ and $\delta_p=\dim V^{I_p}$. Exactly one of the following holds.
\begin{center}\small
\begin{tabular}{@{}lcp{4.4cm}p{2.5cm}cp{2.1cm}@{}}
\toprule
type & $v_p(N)$ & $\WD_p(f)$ & $\pi_p$ & $\delta_p$ & $a_p(f)$\\
\midrule
(U) & $0$ & $\mu_1\oplus\mu_2$ unramified, $\mu_1\mu_2=\norm{\cdot}$ & unramified principal series & $2$ & $\mu_1(\Frob)+\mu_2(\Frob)$\\
(St) & $1$ & $\Sp(2)\otimes\mu$, $\mu$ unramified, $\mu^2=1$ & unramified twist of $\St$ & $1$ & $\mu(\Frob_p)=\pm1$\\
(RSt) & $2a(\chi)$ & $\Sp(2)\otimes\mu\chi$, $\mu$ unramified, $\chi$ ramified, $(\mu\chi)^2=1$ & ramified twist of $\St$ & $0$ & $0$\\
(RPS) & $2a(\mu_1)$ & $\mu_1\oplus\mu_2$, both ramified, $\mu_1\mu_2=\norm{\cdot}$ & ramified principal series & $0$ & $0$\\
(SC) & $\ge2$ & irreducible, $N=0$ & supercuspidal & $0$ & $0$\\
\bottomrule
\end{tabular}
\end{center}
In particular $a_p(f)=0$ and $V^{I_p}=0$ whenever $p^2\mid N$, and $\delta_p=2,1,0$ according as $p\nmid N$, $p\,\|\,N$, $p^2\mid N$. In case (St), with $\varepsilon=\mu$ viewed as an
unramified character of $G_{\Q_p}$ ($\varepsilon(\Frob_p)=a_p(f)$):
\begin{enumerate}[label=(\alph*)]
\item every $G_{\Q_p}$-stable $\mathcal{O}_\lambda$-lattice $L\subset V$ has a basis in which $\rho_{f,\lambda}|_{G_{\Q_p}}=\sm{\varepsilon\chi_\ell}{*}{0}{\varepsilon}$, the corner being nonzero on $I_p$;
\item the reduction $\bar\rho_L|_{D_p}$ is an extension of $\bar\varepsilon$ by $\bar\varepsilon\bar\chi_\ell$; it is unramified if and only if the corner vanishes on $I_p$ modulo $\lambda$,
and when it is ramified it is non-split and has exactly one stable line;
\item if $\bar\rho_L$ is unramified at $p$, the eigenvalues of $\bar\rho_L(\Frob_p)$ are $a_p(f)p$ and $a_p(f)$ modulo $\lambda$.
\end{enumerate}
\end{theorem}

\begin{proof}
By \Cref{thm:modgalrep:carayol}(ii), $a(V)=v_p(N)$, and $\det V=\chi_\ell|_{G_{\Q_p}}$ is unramified, corresponding to $\norm{\cdot}$ on $W_{\Q_p}$.

\emph{$a(V)=0$.} Then $V$ is unramified, $N=0$, and $r(\Frob)$ has the eigenvalues $\alpha,\beta$ of \Cref{thm:modgalrep:carayol}(iii); this is (U).

\emph{$a(V)=1$.} If $V^{I_p}=V$ then $a(V)=0$. So $\dim V^{I_p}\le1$, and $1=a(V)=2-\dim V^{I_p}+\operatorname{sw}(V)$ forces $\operatorname{sw}(V)=0$ and $\dim V^{I_p}=1$; so $V$ is tame. If $N=0$,
then $r|_{I_p}$, a representation of a finite group in characteristic $0$, is semisimple with a fixed line: $r|_{I_p}\simeq1\oplus\psi$ with $\psi\ne1$ (as $\dim V^{I_p}=1$).
But $\det r|_{I_p}=\psi$ must be trivial, because $\det V$ is unramified. So $N\ne0$, and $V\simeq\Sp(2)\otimes\mu$ for the character $\mu$ on $V/\ker N$
(\Cref{rem:modgalrep:wd-recap}). As $a(\Sp(2)\otimes\mu)=2a(\mu)\ge2$ for ramified $\mu$, $\mu$ is unramified. Comparing determinants, $\norm{\cdot}\mu^2=\norm{\cdot}$, so $\mu^2=1$. Now
$V^{I_p}=\ker N$ is the line on which $W_{\Q_p}$ acts by $\norm{\cdot}\mu$, so $\Frob_p$ acts on it by $p\,\mu(\Frob_p)$, and \Cref{thm:modgalrep:carayol}(iii) gives
$a_p(f)=\mu(\Frob_p)=\pm1$. This is (St).

\emph{$a(V)\ge2$.} $\dim V^{I_p}=2$ is impossible ($V$ would be unramified). Suppose $\dim V^{I_p}=1$. If $N=0$, the argument just given shows that
$r|_{I_p}$ is trivial, so $V^{I_p}=V$, a contradiction. If $N\ne0$, $V\simeq\Sp(2)\otimes\mu$ and $V^{I_p}=(\ker N)^{r(I_p)}$ is a line only if $\mu$ is unramified,
and then $a(V)=1$, a contradiction. Hence $V^{I_p}=0$, $\delta_p=0$, $L(V,s)=1$, and $a_p(f)=0$ by \Cref{thm:modgalrep:carayol}(iii). If $N\ne0$,
$V\simeq\Sp(2)\otimes\nu$ with $\nu$ ramified and $\nu^2=1$ (determinants again); writing $\nu=\mu\chi$ with $\chi$ a ramified character and $\mu$ unramified gives (RSt), with
$a(V)=2a(\nu)=2a(\chi)$. If $N=0$, $r$ is Frobenius-semisimple of dimension $2$, so it is a sum of two characters, both ramified as $V^{I_p}=0$, with product
$\norm{\cdot}$ (so $\mu_2|_{I_p}=\mu_1^{-1}|_{I_p}$ and $a(\mu_2)=a(\mu_1)$, $a(V)=2a(\mu_1)$): this is (RPS); or it is irreducible: (SC). In (SC), $a(V)\ge2$ since $V^{I_p}=0$.

The types are mutually exclusive (they are distinguished by $\delta_p$, by $N$ and by reducibility), and the cases $v_p(N)=0$, $1$, $\ge2$ are (U), (St),
and the last three, which proves the statements on $a_p$ and $\delta_p$.

(a) In case (St), $V$ is the representation $\sm{\varepsilon\chi_\ell}{*}{0}{\varepsilon}$ of \Cref{rem:modgalrep:wd-recap} with $\mu=\varepsilon$. The stable line $\ker N$ meets $L$ in
a saturated rank-one submodule $L_1$, which is a direct summand of $L$; a basis $e_1$ of $L_1$ extended by $e_2$ to a basis of $L$ gives the matrix
form, with entries in $\mathcal{O}_\lambda$. (b) Reduce modulo $\lambda$; the corner restricted to $I_p$ is a homomorphism (the diagonal is trivial on $I_p$), so the
reduction is unramified exactly when it vanishes on $I_p$ modulo $\lambda$. If the reduction is ramified it is non-split (a split extension of two
unramified characters is unramified), and a non-split extension of two characters has a unique stable line (\Cref{ch:reps}; this holds also
when the two characters coincide, which happens when $p\equiv1\pmod\ell$). (c) Read off the diagonal: $\bar\varepsilon\bar\chi_\ell(\Frob_p)=a_p(f)p$ and $\bar\varepsilon(\Frob_p)=a_p(f)$.
\end{proof}

\begin{paperlink}
Section 1.1 of the paper: ``for a prime $p\ne\ell$ let $\delta_p=\dim V^{I_p}$, so $\delta_p=2$, $1$, $0$ for good, multiplicative, additive reduction''. For $f=f_E$ this is
\Cref{thm:modgalrep:local-types} (and, on the elliptic-curve side, \Cref{ch:galrep}). The case $p^2\mid N$ of Lemma lem:old (``a newform on $\Gamma_0(M)$ with $p^2\mid M$
has $a_p=0$ \cite[Theorem 3]{AtkinLehner70}'') is the rows (RSt), (RPS), (SC) of the table, here derived from the Galois side; and the case $p\,\|\,N$ of
Lemma lem:old (``$\rho_{E,\ell}|_{D_p}$ is an extension of an unramified character $\varepsilon$ by $\varepsilon\chi_\ell$ with $\varepsilon(\Frob_p)=a_p(f)=\pm1$ \cite[Theorem 3.1(e)]{DDT}'') is
part (a), applied to $f$ and to the newform $h$ of that lemma.
\end{paperlink}

\begin{proposition}[the local type of $f_E$ by Kodaira type]\label{prop:modgalrep:kodaira-dictionary}
Let $E/\Q$ have bad reduction at a prime $q\ge5$, $f=f_E$, $\ell\ne q$. The type of $f$ at $q$ is given by the Kodaira type of $E$ at $q$ as follows.
\begin{center}\small
\begin{tabular}{@{}lll@{}}
\toprule
Kodaira type at $q$ & $\rho_{E,\ell}(I_q)$ & type of $f$ at $q$\\
\midrule
$\kod{I}_n$, $n\ge1$ & unipotent, infinite & (St)\\
$\kod{I}_n^*$, $n\ge1$ & $\{\pm1\}\times$unipotent, infinite & (RSt), with $\chi=\chi_{q^*}$ on $I_q$\\
$\kod{I}_0^*$ & $\{\pm1\}$ & (RPS), $\mu_i=\chi_{q^*}\cdot$(unramified)\\
$\kod{II},\kod{II}^*,\kod{IV},\kod{IV}^*$ & cyclic of order $6$, $6$, $3$, $3$ & (RPS) if $q\equiv1\pmod3$, (SC) if $q\equiv2\pmod3$\\
$\kod{III},\kod{III}^*$ & cyclic of order $4$ & (RPS) if $q\equiv1\pmod4$, (SC) if $q\equiv3\pmod4$\\
\bottomrule
\end{tabular}
\end{center}
In the row $\kod{I}_0^*$, $V_\ell E|_{G_{\Q_q}}\simeq V'\otimes\chi_{q^*}$ with $V'=V_\ell E^{(q)}$ unramified, and the two characters are $\chi_{q^*}$ times the unramified characters sending
$\Frob_q$ to the Frobenius eigenvalues of $E^{(q)}$. In the row $\kod{I}_n^*$, $V_\ell E|_{G_{\Q_q}}\simeq V'\otimes\chi_{q^*}$ with $V'=V_\ell E^{(q)}$ of type (St).
\end{proposition}

\begin{proof}
The images of inertia are those of \Cref{ch:galrep} (the Tate shape at multiplicative primes, the potentially multiplicative case, and the
theorem on additive potentially good reduction at $q\ge5$, from Kraus's theorem). Type $\kod{I}_n$: $\delta_q=1$, so (St) by \Cref{thm:modgalrep:local-types}.
Types $\kod{I}_n^*$ and $\kod{I}_0^*$: $E^{(q)}$ has type $\kod{I}_n$, resp.\ good reduction (the paper's Lemma lem:twist(ii), proved in \Cref{ch:reduction}), and
$V_\ell E\simeq V_\ell E^{(q)}\otimes\chi_{q^*}$ (\Cref{ch:ellcurves}); so $\WD_q(f)$ is $(\Sp(2)\otimes\mu)\otimes\chi_{q^*}$, resp.\ $(\nu_{\alpha'}\oplus\nu_{\beta'})\otimes\chi_{q^*}$ with $\nu_x$ the
unramified character sending $\Frob_q$ to $x$. The other rows: inertia acts through a finite cyclic group $\langle g\rangle$ of order $m\in\{3,4,6\}$, so $N=0$; for $q\ge5$
the action is tame and $\Frob_qg\Frob_q^{-1}=g^q$ (the tame relation, \Cref{ch:local}). The eigenvalues of $g$ are two distinct primitive $m$-th roots of unity
$\zeta,\zeta^{-1}$. As $q\ge5$ is prime to $m$ and the units modulo $3$, $4$, $6$ are $\pm1$, either $q\equiv1\pmod m$ or $q\equiv-1\pmod m$. In the first case $\Frob_q$
commutes with $g$, hence preserves its two eigenlines (over $\overline\Q_\ell$), and $\WD_q(f)$ is a sum of two characters, ramified since $\zeta\ne1$: (RPS). In
the second case $\Frob_q$ conjugates $g$ to $g^{-1}$ and exchanges the two eigenlines, so no line is stable under $W_{\Q_q}$ (a stable line would be an
eigenline of $g$), and $\WD_q(f)$ is irreducible: (SC). Finally $q\equiv1\pmod 6$ iff $q\equiv1\pmod3$ for odd $q$.
\end{proof}

\begin{example}[local types of the paper's curves]\label{ex:modgalrep:local-types}
(Checked with PARI/GP's \texttt{elllocalred}.) \curve{11a1} at $11$: $\kod{I}_5$, (St), $a_{11}=1$. \curve{121d2} at $11$: $\kod{I}_5^*$, (RSt). \curve{539d1} at $7$: $\kod{I}_0^*$, (RPS),
the twist being \curve{11a3}. \curve{121c1} at $11$: $\kod{IV}$, and $11\equiv2\pmod3$, so (SC). \curve{121a1} at $11$: $\kod{II}$, (SC). \curve{726a1} at $11$: $\kod{III}$, and $11\equiv3\pmod
4$, so (SC). \curve{52a1} at $2$ has type $\kod{IV}^*$; the table is for $q\ge5$, but by \Cref{ch:galrep} the image of $I_2$ is cyclic of order $3$ here as well, and
since $2\equiv-1\pmod3$ the argument of the proof gives (SC) at $2$.
\end{example}

\begin{paperlink}
The proof of Lemma lem:euler says, for type $\kod{I}_n$: ``The local component of $f$ is an unramified twist of the Steinberg representation'', and
for type $\kod{I}_n^*$: ``$V=V'\otimes\chi_{q^*}$ with $V'$ Steinberg at $q$''. These are the first two rows of \Cref{prop:modgalrep:kodaira-dictionary}; the
proof of Lemma lem:euler uses them only through $\WD_q(f)$ and $(\ad V)^{I_q}$, which \Cref{ch:galrep} and \Cref{ch:lfunctions} compute.
\end{paperlink}

\begin{corollary}[$\delta_p$ and the naive factor]\label{cor:modgalrep:delta-naive}
Let $f$ be a newform of weight $2$ on $\Gamma_0(N)$ and $p\mid N$, $p\ne\ell$. Then $a_p(f)\in\{\pm1\}$ iff $\delta_p=1$, and $a_p(f)=0$ iff $\delta_p=0$. The naive local factor
$L_p^{\nv}(A_f,s)$ of the paper, which equals the true factor $L_p(A_f,s)$ if $\delta_p>0$ and $1$ if $\delta_p=0$ (equivalently $L_p^{\nv}(A_f,s)^{-1}=1-a_p(f)^2p^{-1-s}$), is
determined by $a_p(f)$; the true factor is determined by $\WD_p(f)$ and is unchanged when $\rho_{f,\lambda}$ is replaced by $\rho_{f,\lambda}\otimes\chi$ for a character $\chi$,
because $\ad(\rho\otimes\chi)=\ad\rho$. The naive factor is not twist-invariant: $f_{11a}$ at $11$ has type (St) and $f_{121d}=f_{11a}^{\chi_{-11}}$ has type (RSt), so
their naive factors at $11$ are $1-11^{-1-s}$ and $1$, while their true factors coincide.
\end{corollary}

\begin{proof}
The first sentence is \Cref{thm:modgalrep:local-types}. If $\delta_p=1$, $(\ad V)^{I_p}$ is the line spanned by $N$, on which the geometric Frobenius acts by
$p^{-1}$ (because $r(w)Nr(w)^{-1}=\norm wN$), so $L_p(A_f,s)^{-1}=1-p^{-1-s}=1-a_p^2p^{-1-s}$; if $\delta_p=0$, $a_p=0$ and the naive factor is $1$ by definition. The
twist-invariance of $\ad$ is \Cref{ch:reps}. The computation of $L_p(A_f,s)$ in every case is \Cref{ch:lfunctions}.
\end{proof}

\begin{proposition}[the conductor of a quadratic twist, type by type]\label{prop:modgalrep:twist-conductor}
Let $q$ be an odd prime, $\chi$ a ramified quadratic character of $G_{\Q_q}$ (so $\chi$ is tame and $a(\chi)=1$, and $\chi|_{I_q}=\chi_{q^*}|_{I_q}$), and $f$ a newform of
weight $2$ with trivial character, of type $t$ at $q$. Then:
\begin{center}\small
\begin{tabular}{@{}p{5.2cm}ccl@{}}
\toprule
type of $f$ at $q$ & $a(\WD_q(f))$ & $a(\WD_q(f)\otimes\chi)$ & type of $\WD_q(f)\otimes\chi$\\
\midrule
(U) & $0$ & $2$ & (RPS)\\
(St) & $1$ & $2$ & (RSt)\\
(RSt) & $2$ & $1$ & (St)\\
(RPS) with $\mu_i|_{I_q}=\chi|_{I_q}$ & $2$ & $0$ & (U)\\
(RPS) or (SC), tame, inertia cyclic of order $3$, $4$ or $6$ & $2$ & $2$ & (RPS) or (SC)\\
\bottomrule
\end{tabular}
\end{center}
At a prime $p\ne q$ the twist by a character unramified at $p$ changes neither the type nor the conductor exponent.
\end{proposition}

\begin{proof}
(U): $\WD\otimes\chi=\mu_1\chi\oplus\mu_2\chi$ with both characters ramified of conductor $1$, so $a=2$. (St): $\Sp(2)\otimes\mu\chi$ with $\mu\chi$ ramified: $a=2a(\chi)=2$. (RSt): for
$q$ odd, every ramified character $\nu$ with $\nu^2=1$ agrees with $\chi$ on $I_q$ (the ramified quadratic characters of $\Q_q^\times$ differ by an unramified one), so
$\WD=\Sp(2)\otimes\nu$ with $\nu\chi$ unramified and $\WD\otimes\chi=\Sp(2)\otimes\nu\chi$ has $a=1$; in particular $v_q(N)=2$ for type (RSt) at odd $q$. (RPS) with
$\mu_i|_{I_q}=\chi|_{I_q}$: both $\mu_i\chi$ are unramified, $a=0$. Last row: inertia acts on $\WD$ through $\langle g\rangle$, $g$ of order $m\in\{3,4,6\}$ with eigenvalues $\zeta^{\pm1}$
primitive $m$-th roots of unity, and on $\WD\otimes\chi$ through $\langle-g\rangle$ (with $\chi(t)=-1$ on a generator $t$ of tame inertia mapping to $g$); $-g$ has order $6$,
$4$, $3$ and eigenvalues $-\zeta^{\pm1}\ne1$, so $(\WD\otimes\chi)^{I_q}=0$, the twist is tame, and $a=2$; its type is (RPS) or (SC) by the argument of
\Cref{prop:modgalrep:kodaira-dictionary}. At $p\ne q$: an unramified twist preserves $a$ and the shape (\Cref{rem:modgalrep:wd-recap}).
\end{proof}

\begin{corollary}[Atkin and Li; the paper's Lemma lem:twist(iii)]\label{cor:modgalrep:atkin-li}
Let $g$ be a newform of weight $2$ on $\Gamma_0(M)$ and $q$ an odd prime with $q^2\nmid M$; let $\chi=\chi_{q^*}$. Then $g^\chi$ has level $\lcm(M,q^2)$, type (RPS) or
(RSt) at $q$, and $a_q(g^\chi)=0$; consequently $g\otimes\chi=g^\chi$ is a newform on $\Gamma_0(\lcm(M,q^2))$: of level $Mq^2$ if $q\nmid M$ and $Mq$ if $q\,\|\,M$.
\end{corollary}

\begin{proof}
The type of $g$ at $q$ is (U) or (St) (\Cref{thm:modgalrep:local-types}), so $\WD_q(g^\chi)=\WD_q(g)\otimes\chi$ has conductor exponent $2$ and type (RPS) or (RSt)
(\Cref{prop:modgalrep:twist-conductor}); at $p\ne q$ the exponents are those of $g$. By \Cref{thm:modgalrep:carayol}(ii) the level of $g^\chi$ is $\lcm(M,q^2)$, and
$a_q(g^\chi)=0$ by \Cref{thm:modgalrep:local-types}. By \Cref{prop:modgalrep:twists}(ii), $g\otimes\chi=(g^\chi)^{\{q\}}=g^\chi$.
\end{proof}

\begin{paperlink}
The proof of Lemma lem:twist(iii) in the paper: ``in the case $q\,\|\,M$, the local component of $g$ at $q$ is an unramified twist of the Steinberg
representation, whose twist by a ramified quadratic character has conductor exponent $2$''. This is the row (St) of \Cref{prop:modgalrep:twist-conductor};
the row (U) is the case $q\nmid M$, and \Cref{cor:modgalrep:atkin-li} is the statement of Lemma lem:twist(iii) (Atkin and Li \cite[Theorem 3.1]{AtkinLi78}, as the
paper cites it). The hypothesis $q^2\nmid M$ cannot be dropped: for $q^2\mid M$ the rows (RSt) and (RPS) show that $g^\chi$ may have smaller level (\curve{539c1},
\curve{539d1} below), and then $g\otimes\chi=(g^\chi)^{\{q\}}$ differs from $g^\chi$ whenever $a_q(g^\chi)\ne0$.
\end{paperlink}

\begin{proposition}[twists of newforms]\label{prop:modgalrep:twists}
Let $g$ be a newform of weight $2$ on $\Gamma_0(M)$ and $\chi$ a quadratic Dirichlet character of conductor $r$.
\begin{enumerate}[label=(\roman*)]
\item $\rho_{g,\lambda}\otimes\chi$ is unramified outside $Mr\ell$ and $\tr(\rho_{g,\lambda}\otimes\chi)(\Frob_p)=\chi(p)a_p(g)$ for $p\nmid Mr\ell$.
\item There is a unique newform $g^\chi$ of weight $2$ with trivial character such that $\rho_{g^\chi,\lambda}\simeq\rho_{g,\lambda}\otimes\chi$ (for one, equivalently every, $\lambda$).
Its level divides $\lcm(M,r^2)$ and equals
\[
M\prod_{q\mid r}q^{a(\WD_q(g)\otimes\chi)-a(\WD_q(g))}.
\]
It satisfies $a_n(g^\chi)=\chi(n)a_n(g)$ for every $n$ prime to $r$, and
$g\otimes\chi=(g^\chi)^{\Sigma(r)}$.
\item $(g^\chi)^\chi=g$.
\item For $E/\Q$ and a fundamental discriminant $d$, $f_{E^{(d)}}=(f_E)^{\chi_d}$.
\end{enumerate}
\end{proposition}

\begin{proof}
(i) is clear, $\chi$ being unramified outside $r$ with $\chi(\Frob_p)=\chi(p)$. (ii) Uniqueness: two such newforms have the same $a_p$ for almost all $p$, so
they are equal (\Cref{prop:modgalrep:newforms-recap}(iii)). Existence: by \Cref{ch:hecke} (for $\chi=\chi_S$, the case used in this book; for any quadratic $\chi$ see Atkin and Li \cite{AtkinLi78}) the naive
twist $g\otimes\chi$ lies in $S_2(\Gamma_0(L))$, $L=\lcm(M,r^2)$ ($\chi^2$ is trivial on $(\Z/L)^\times$), and is an eigenform of $T_p$ with eigenvalue $\chi(p)a_p(g)$ for $p\nmid L$. Decompose it along
\Cref{prop:modgalrep:newforms-recap}(iv); its nonzero components are eigenforms with the same eigenvalues, so it lies in a single $V_h$ with
$a_p(h)=\chi(p)a_p(g)$ for all $p\nmid L$, and $h$ is a newform of level dividing $L$. By (i) and \Cref{cor:modgalrep:congruent-newforms} (Chebotarev and
Brauer-Nesbitt in characteristic $0$; both representations are irreducible), $\rho_{h,\lambda}\simeq\rho_{g,\lambda}\otimes\chi$; put $g^\chi=h$. Its level is the conductor of
$\rho_{g,\lambda}\otimes\chi$ (\Cref{thm:modgalrep:carayol}(ii), with $\ell$ chosen prime to $Mr$), which differs from that of $\rho_{g,\lambda}$ only at the primes dividing $r$.
Coefficients: for $p\nmid r$, $\WD_p(g^\chi)=\WD_p(g)\otimes\chi$ with $\chi$ unramified at $p$, so $\delta_p$ is the same for both and, by
\Cref{thm:modgalrep:carayol}(iii) (applied with an $\ell\ne p$), $a_p(g^\chi)=\chi(p)a_p(g)$ for all $p\nmid r$: if $\delta_p=2$ the Frobenius eigenvalues are multiplied by
$\chi(p)$; if $\delta_p=1$ the eigenvalue $p\,a_p$ on $V^{I_p}$ is; if $\delta_p=0$ both sides vanish. The multiplicativity and the recursions of
\Cref{prop:modgalrep:newforms-recap}(i), which $n\mapsto\chi(n)a_n(g)$ also satisfies on $n$ prime to $r$ (because $\chi(p)^2=1$), give $a_n(g^\chi)=\chi(n)a_n(g)$ for $(n,r)=1$.
Both $g\otimes\chi$ and $(g^\chi)^{\Sigma(r)}$ have $n$-th coefficient $\chi(n)a_n(g)$ for $(n,r)=1$ and $0$ otherwise, so they are equal. (iii) $\rho_{(g^\chi)^\chi}\simeq\rho_g\otimes\chi^2=\rho_g$
and uniqueness. (iv) $\rho_{E^{(d)},\ell}\simeq\rho_{E,\ell}\otimes\chi_d$ (\Cref{ch:ellcurves}), and both $f_{E^{(d)}}$ and $(f_E)^{\chi_d}$ have this representation
(\Cref{thm:modgalrep:comparison-E}(i)); apply uniqueness.
\end{proof}

\begin{example}[levels of twists]\label{ex:modgalrep:twist-levels}
(Each twist computed with PARI/GP's \texttt{elltwist}, reduced to a minimal model and identified in Cremona's tables; each newform identity
checked on $a_p$, $p<400$, $p\ne q$.) With $\chi=\chi_{-11}$ or $\chi_{-7}$ as appropriate:
\begin{center}\small
\begin{tabular}{@{}llllll@{}}
\toprule
curve & type at $q$ & row & twist & level change & newform identity\\
\midrule
\curve{11a1} & $\kod{I}_5$ at $11$, (St) & (St) & \curve{121d2} & $11\to121$ & $f_{11a}^{\chi_{-11}}=f_{11a}\otimes\chi_{-11}=f_{121d}$\\
\curve{11a3} & $\kod{I}_1$ at $11$, (St) & (St) & \curve{121d1} & $11\to121$ & the same class\\
\curve{11a1} & good at $7$, (U) & (U) & \curve{539d2} & $11\to539$ & $f_{11a}^{\chi_{-7}}=f_{11a}\otimes\chi_{-7}=f_{539d}$\\
\curve{20a1} & good at $7$, (U) & (U) & \curve{980g2} & $20\to980$ & $f_{20a}^{\chi_{-7}}=f_{20a}\otimes\chi_{-7}=f_{980g}$\\
\curve{121a1} & $\kod{II}$ at $11$, (SC) & last & \curve{121c2} & $121\to121$ & $f_{121a}^{\chi_{-11}}=f_{121a}\otimes\chi_{-11}=f_{121c}$\\
\curve{539c1} & $\kod{I}_3^*$ at $7$, (RSt) & (RSt) & \curve{77c1} & $539\to77$ & $f_{539c}^{\chi_{-7}}=f_{77c}$, $a_7(f_{77c})=-1$\\
\curve{539d1} & $\kod{I}_0^*$ at $7$, (RPS) & (RPS) & \curve{11a3} & $539\to11$ & $f_{539d}^{\chi_{-7}}=f_{11a}$\\
\curve{980i1} & $\kod{I}_5^*$ at $7$, (RSt) & (RSt) & \curve{140b1} & $980\to140$ & $f_{980i}^{\chi_{-7}}=f_{140b}$, $a_7(f_{140b})=-1$\\
\bottomrule
\end{tabular}
\end{center}
In the three rows where the level drops, the naive twist is a depletion of the newform twist: $f_{539c}\otimes\chi_{-7}=f_{77c}^{\{7\}}$, $f_{980i}\otimes\chi_{-7}=f_{140b}^{\{7\}}$,
$f_{539d}\otimes\chi_{-7}=f_{11a}^{\{7\}}$. The twist of the optimal curve \curve{121c1} by $\Q(\sqrt{-11})$ is \curve{121a2}, not \curve{121a1}: the twist of a curve lands in the
twisted class, but not necessarily on the curve with the same index.
\end{example}

\section{The representation at \texorpdfstring{$\ell$}{l} when \texorpdfstring{$\ell\nmid N$}{l does not divide N}}\label{sec:modgalrep:at-l}

At the prime $\ell$ itself the representation $\rho_{f,\lambda}$ is not described by a Weil-Deligne representation; when $\ell\nmid N$ it is crystalline, and its
reduction is finite flat. The paper uses this twice: in the optimal-level paragraph of Section 3 (``the restriction to $D_\ell$ comes from the finite
flat group scheme $E'[\ell]$ and is good''), and in Lemma lem:old at $p=\ell$, where $a_\ell\bmod\lambda$ must be read off from $\bar\rho|_{D_\ell}$. For $\rho_{E,\ell}$ the
same facts are in \Cref{ch:galrep}; here they are stated for newforms.

\begin{theorem}[the representation at $\ell\nmid N$]\label{thm:modgalrep:crystalline}
Let $f$ be a newform of weight $2$ on $\Gamma_0(N)$, $\lambda\mid\ell$, $\ell\nmid N$.
\begin{enumerate}[label=(\roman*)]
\item $J_f$ has good reduction at $\ell$; for every $k\ge1$ the $G_{\Q_\ell}$-module $J_f[\ell^k]$ is the generic fibre of a finite flat group scheme over $\Zl$ (the
$\ell^k$-torsion of the N\'eron model, an abelian scheme), and so is every $G_{\Q_\ell}$-stable subquotient of it. In particular the reduction $\bar\rho_{f,\lambda}|_{D_\ell}$
of any stable lattice is \emph{finite}, which is ``good'' in the sense of \cite[Theorem 3.15]{DDT}.
\item $\rho_{f,\lambda}|_{G_{\Q_\ell}}$ is crystalline with Hodge-Tate weights $0$ and $1$, and the crystalline Frobenius on $\Dcris(V_\lambda(f))$ has characteristic
polynomial $X^2-a_\ell(f)X+\ell$. For $\ell\ge3$ the weights lie in the Fontaine-Laffaille range $[0,\ell-2]$, so that stable lattices correspond to
Fontaine-Laffaille modules.
\item If $\bar\rho_{f,\lambda}$ is irreducible, its Serre weight is $2$.
\end{enumerate}
\end{theorem}

\begin{proof}[Sketch of proof]
(i) Good reduction of $J_0(N)$ at $\ell\nmid N$ is the theorem of Igusa and of Deligne and Rapoport \cite{DeligneRapoport73}, and it passes to the quotient
$J_f$; the $\ell^k$-torsion of an abelian scheme is finite flat; scheme-theoretic closure in a finite flat group scheme shows that subquotients of
finite flat Galois modules are finite flat \cite{DDT}. (ii) Crystallinity of the Tate module of an abelian variety with good reduction is the
crystalline comparison theorem (Fontaine, Messing, Faltings); see \cite{Berger}. The characteristic polynomial of the crystalline Frobenius is that of
the Frobenius endomorphism of the reduction, which is $X^2-a_\ell X+\ell$ on the $\lambda$-part by the Eichler-Shimura relation, as in
\Cref{thm:modgalrep:eichler-shimura}. The Fontaine-Laffaille correspondence is \cite{FontaineLaffaille}. (iii) Serre's recipe for the weight
\cite{Serre87}, as completed by Edixhoven \cite{Edixhoven92}, gives weight $2$ for an irreducible $\bar\rho$ whose restriction to $D_\ell$ is finite flat;
\Cref{thm:modgalrep:deligne-fontaine} describes $\bar\rho|_{I_\ell}$ explicitly. We do not prove these statements.
\end{proof}

\begin{theorem}[Deligne; Fontaine]\label{thm:modgalrep:deligne-fontaine}
Let $\ell$ be odd, $f$ a newform of weight $2$ on $\Gamma_0(N)$ with $\ell\nmid N$, $\lambda\mid\ell$, and $\bar\rho=\bar\rho_L$ for a stable lattice $L$. Let $\omega=\bar\chi_\ell|_{I_\ell}$ (the
fundamental character of level $1$) and $\omega_2$ a fundamental character of level $2$ (\Cref{ch:galrep}). For $x\in k_\lambda^\times$ let $\mu(x)$ be the unramified
character of $G_{\Q_\ell}$ with $\mu(x)(\Frob_\ell)=x$.
\begin{enumerate}[label=(\roman*)]
\item (Deligne.) If $a_\ell(f)\not\equiv0\pmod\lambda$, then $\bar\rho|_{D_\ell}\simeq\sm{\bar\chi_\ell\,\mu(\bar a_\ell^{-1})}{*}{0}{\mu(\bar a_\ell)}$, where $\bar a_\ell=a_\ell(f)\bmod\lambda$.
\item (Fontaine.) If $a_\ell(f)\equiv0\pmod\lambda$, then $\bar\rho|_{I_\ell}\otimes\bar k_\lambda\simeq\omega_2\oplus\omega_2^\ell$, and $\bar\rho|_{D_\ell}$ is irreducible, even over $\bar k_\lambda$.
\item Consequently $a_\ell(f)\bmod\lambda$ is determined by the isomorphism class of $\bar\rho|_{D_\ell}$: it is $0$ if $\bar\rho|_{D_\ell}$ is irreducible; otherwise
$\bar\rho|_{D_\ell}$ has exactly one one-dimensional unramified quotient, and $a_\ell(f)\bmod\lambda$ is the eigenvalue of $\Frob_\ell$ on it. In particular, if $h$ is a
newform of weight $2$ with trivial character, of level prime to $\ell$, with coefficients in $\mathcal{O}$, and $\bar\rho_h\simeq\bar\rho_f$ (over $k_\lambda$), then
$a_\ell(h)\equiv a_\ell(f)\pmod\lambda$.
\end{enumerate}
\end{theorem}

\begin{proof}
We do not prove (i) and (ii); see Edixhoven \cite[Section 2]{Edixhoven92}, where both are stated with references, Wiles \cite{Wiles88} for the
ordinary case, and Serre \cite[Section 1]{Serre72} for the elliptic-curve case. (iii) In case (ii) $\bar\rho|_{D_\ell}$ is irreducible and in case (i) it is not;
so $\bar\rho|_{D_\ell}$ irreducible forces $a_\ell\equiv0$. If $\bar\rho|_{D_\ell}$ is reducible we are in case (i). A one-dimensional quotient of $\bar\rho|_{D_\ell}$ is the quotient by a
stable line; the stable lines are the given one, with quotient $\mu(\bar a_\ell)$, and, only if the extension splits, a second line, with quotient
$\bar\chi_\ell\mu(\bar a_\ell^{-1})$. The latter is ramified, because $\omega$ has order $\ell-1\ge2$ for $\ell$ odd. So the only unramified one-dimensional quotient is
$\mu(\bar a_\ell)$, on which $\Frob_\ell$ acts by $\bar a_\ell$. Applying this to $f$ and to $h$, whose restrictions to $D_\ell$ are isomorphic, gives $a_\ell(h)\equiv a_\ell(f)$.
\end{proof}

\begin{example}\label{ex:modgalrep:deligne-fontaine}
$a_3(f_{121c})=2$ and $a_3(f_{11a})=-1$, both nonzero modulo $3$ and congruent modulo $3$; this is the case $p=\ell=3$ of Lemma lem:old for the pair
$(f,h)=(f_{121c},f_{11a})$ of \Cref{ex:modgalrep:old-examples}. Likewise $a_3(f_{121a})=2\equiv-1=a_3(f_{121d})\pmod3$, as the paper's Example needs.
\end{example}

\begin{paperlink}
Lemma lem:old, case $p=\ell$: ``By the theorems of Deligne and Fontaine describing $\bar\rho_f|_{D_\ell}$ and $\bar\rho_h|_{D_\ell}$ \cite[\S2]{Edixhoven92}, the residue
class of $a_\ell$ is determined by the restriction to $D_\ell$: it is $0$ when that restriction is irreducible, and otherwise the eigenvalue of Frobenius on
the unramified quotient.'' This is \Cref{thm:modgalrep:deligne-fontaine}(iii); the word ``the'' in ``the unramified quotient'' is justified there (uniqueness
needs $\ell$ odd). The remark after the paper's Example lists Lemma lem:old at $p=\ell$ among the uses of $\ell\nmid N$: \Cref{thm:modgalrep:deligne-fontaine}
needs both levels prime to $\ell$.
\end{paperlink}

\section[Serre's conjecture]{Serre's conjecture: conductor, weight and character}\label{sec:modgalrep:serre}

This section defines the invariants that Serre attached to an odd irreducible mod $\ell$ representation and proves that the Serre conductor of
$\bar\rho_{f,\lambda}$ divides the level of $f$. The paper uses the Serre conductor in the optimal-level paragraph and in cases (a), (b), (c) of the proof
of Theorem thm:qual (``the Serre conductor $M=N(\bar\rho)$ has exponent exactly $1$ at $q$, so $M\mid N/q$''); the divisibility $N(\bar\rho)\mid N$ is what puts the
lowered level inside the original one.

\begin{definition}[Serre's invariants]\label{def:modgalrep:serre-invariants}
Let $\bar\sigma\colon G_\Q\to\GL_2(\Fbar_\ell)$ be continuous, odd ($\det\bar\sigma(c)=-1$) and irreducible.
\begin{enumerate}[label=(\alph*)]
\item The \emph{Serre conductor} is $N(\bar\sigma)=\prod_{p\ne\ell}p^{a_p(\bar\sigma)}$, $a_p(\bar\sigma)=2-\dim\bar\sigma^{I_p}+\operatorname{sw}_p(\bar\sigma)$, the Artin conductor exponent of $\bar\sigma|_{G_{\Q_p}}$
(\Cref{ch:galrep}). It is prime to $\ell$ by definition and prime to every $p$ at which $\bar\sigma$ is unramified.
\item The \emph{Serre weight} $k(\bar\sigma)\ge2$ is defined by Serre's recipe in terms of $\bar\sigma|_{I_\ell}$ \cite[Section 2]{Serre87}, \cite{Edixhoven92}; it is $2$ when
$\bar\sigma|_{D_\ell}$ is finite flat.
\item The \emph{character} $\varepsilon(\bar\sigma)\colon(\Z/N(\bar\sigma))^\times\to\Fbar_\ell^\times$ is defined by $\det\bar\sigma=\varepsilon(\bar\sigma)\bar\chi_\ell^{k(\bar\sigma)-1}$.
\item $\bar\sigma$ \emph{arises from} a newform $g$ (of some weight, level and character) if $\bar\sigma\simeq\bar\rho_{g,\lambda}\otimes_{k_\lambda}\Fbar_\ell$ for some prime $\lambda\mid\ell$ of the
coefficient field of $g$ and some embedding $k_\lambda\into\Fbar_\ell$; when $\bar\sigma$ is absolutely irreducible this is equivalent, for $g$ of weight $2$ with trivial
character and level $M$, to $\tr\bar\sigma(\Frob_p)=a_p(g)\bmod\lambda$ for all $p\nmid M\ell N(\bar\sigma)$ (\Cref{cor:modgalrep:congruent-newforms}). $\bar\sigma$ is \emph{modular} if it
arises from some newform.
\end{enumerate}
\end{definition}

\begin{proposition}[$N(\bar\rho_{f,\lambda})\mid N$; Carayol, Livn\'e]\label{prop:modgalrep:serre-conductor-divides}
Let $f$ be a newform of weight $2$ on $\Gamma_0(N)$, $\lambda\mid\ell$, $L$ a stable lattice, $\bar\rho_L$ its reduction, and $p\ne\ell$. Then
\[
a_p(\bar\rho_L^{\mathrm{ss}})\ \le\ a_p(\bar\rho_L)\ =\ v_p(N)-\bigl(\dim\bar\rho_L^{I_p}-\delta_p\bigr)\ \le\ v_p(N).
\]
Hence the prime-to-$\ell$ part of $N$ is divisible by $N(\bar\rho_{f,\lambda})$ whenever $\bar\rho_{f,\lambda}$ is irreducible, and $N(\bar\rho_{f,\lambda})\mid N$ if $\ell\nmid N$.
\end{proposition}

\begin{proof}
Let $V=L\otimes K_{f,\lambda}$ and $P_p\subset I_p$ the wild inertia group, a pro-$p$ group. The kernel of $\GL_2(\mathcal{O}_\lambda)\to\GL_2(k_\lambda)$ is a pro-$\ell$ group, and a closed
subgroup of it that is also a quotient of a pro-$p$ group is trivial; so $\rho(P_p)$ maps injectively to $\GL_2(k_\lambda)$, and is finite. The Swan conductor of
$V$ (and of $\bar\rho_L$) is computed from the dimensions of the fixed spaces of the ramification groups $G_i$, $i\ge1$, of a finite quotient through which $P_p$
acts (\Cref{ch:local}); these are $p$-groups, of order invertible in $\mathcal{O}_\lambda$. For such a group $G$ the averaging operator $e=\frac1{\#G}\sum_{g\in G}g$ is an
idempotent of $\mathcal{O}_\lambda[G]$, so $L=eL\oplus(1-e)L$ with $eL=L^G$, and reduction modulo $\lambda$ gives $(L/\lambda L)^G=e(L/\lambda L)=L^G/\lambda L^G$. Hence
$\dim V^{G_i}=\operatorname{rank}L^{G_i}=\dim(L/\lambda L)^{G_i}$ for all $i\ge1$, and $\operatorname{sw}(V)=\operatorname{sw}(\bar\rho_L)$. Next, $L^{I_p}=L\cap V^{I_p}$ is saturated in $L$, so
$L^{I_p}\cap\lambda L=\lambda L^{I_p}$ and $L^{I_p}/\lambda L^{I_p}\into(L/\lambda L)^{I_p}$; thus $\dim\bar\rho_L^{I_p}\ge\dim V^{I_p}=\delta_p$. Subtracting the two conductor formulas,
$a(V)-a(\bar\rho_L)=\dim\bar\rho_L^{I_p}-\delta_p\ge0$, and $a(V)=v_p(N)$ by \Cref{thm:modgalrep:carayol}(ii). For the first inequality, let $0\to A\to B\to C\to0$ be an exact
sequence of $k_\lambda[G_{\Q_p}]$-modules with $B=\bar\rho_L$; the Swan conductor is additive (the fixed points of a group of order prime to $\ell$ form an
exact functor), and $\dim B^{I_p}\le\dim A^{I_p}+\dim C^{I_p}$ (invariants are left exact), so $a(B)\ge a(A)+a(C)=a(A\oplus C)$. When $\bar\rho_{f,\lambda}$ is irreducible it
does not depend on $L$ (\Cref{prop:modgalrep:lattices}(iii)), and multiplying over $p\ne\ell$ gives the divisibility.
\end{proof}

For $f=f_E$ the exact value is $a_p(\bar\rho_{E,\ell})=f_p(E)-\dim_{\F_\ell}\Phig_p[\ell]$ (\Cref{ch:galrep}, with examples); Livn\'e \cite{Livne89} lists the possible
exponents in general.

\begin{paperlink}
The optimal-level paragraph of Section 3: ``Recall also that $N(\bar\rho_{E',\ell})$ divides the conductor of $E'$ \cite{Carayol89}, and that $N(\bar\sigma)$ is prime to
every prime at which $\bar\sigma$ is unramified.'' The first clause is \Cref{prop:modgalrep:serre-conductor-divides} for $f=f_{E'}$, the second is
\Cref{def:modgalrep:serre-invariants}(a). Both are used in cases (a), (b), (c) to get $M\mid N/q$ or $M\mid N/q^2$ and the hypothesis of Lemma lem:old.
\end{paperlink}

\begin{theorem}[Serre's conjecture; Khare and Wintenberger]\label{thm:modgalrep:serre-conjecture}
Every odd, absolutely irreducible, continuous $\bar\sigma\colon G_\Q\to\GL_2(\Fbar_\ell)$ arises from a newform of weight $k(\bar\sigma)$, level $N(\bar\sigma)$ and a character
lifting $\varepsilon(\bar\sigma)$.
\end{theorem}

We do not prove this; see Khare and Wintenberger \cite{KhareWintenberger}, Khare \cite{Khare} for level one, Kisin \cite{Kisin09} for an input at
$\ell=2$, and Serre's formulation \cite{Serre87}. The passage from the ``weak'' form (modular of some weight and level) to the ``strong'' form
(the optimal weight and level) is due to Ribet, Carayol, Diamond, Edixhoven, Buzzard and others; \Cref{sec:modgalrep:lowering} states the part
that concerns the level. The paper uses only that part, for $\bar\rho_{E',\ell}$, which is modular because $E'$ is (\Cref{ch:moddeg}); it does not need the
theorem of Khare and Wintenberger.

\section{Lowering the level}\label{sec:modgalrep:lowering}

A mod $\ell$ representation that arises from a newform of level $N$ may also arise at a smaller level. This section states the theorems that
lower the level (Mazur's principle, Ribet's theorem, and the optimal-level theorem of Ribet, Carayol and Diamond in the form of
\cite[Theorem 3.15]{DDT}), proves the specialisation to $\bar\rho_{E',\ell}$ that the paper uses (level exactly $N(\bar\rho)$, trivial character), and works the
paper's examples. This is the step ``lower the level of $\bar\rho$, or of its twist $\bar\rho\otimes\chi_{q^*}$'' of the proof of Theorem thm:qual.

\begin{theorem}[Mazur's principle]\label{thm:modgalrep:mazur-principle}
Let $\ell$ be odd, $f$ a newform of weight $2$ on $\Gamma_0(N)$ with $\ell\nmid N$, $\lambda\mid\ell$ with $\bar\rho=\bar\rho_{f,\lambda}$ irreducible, and $q\,\|\,N$ a prime such that $\bar\rho$ is
unramified at $q$ and $q\not\equiv1\pmod\ell$. Then $\bar\rho$ arises from a newform of weight $2$ on $\Gamma_0(N/q)$.
\end{theorem}

\begin{proof}[Idea of proof]
$\bar\rho$ is realised in $J_0(N)[\mathfrak m]$ for the maximal ideal $\mathfrak m$ of $\T_N$ attached to $(f,\lambda)$. At $q$ the N\'eron model of $J_0(N)$ has semiabelian
reduction whose toric part is governed by the $q$-new quotient (the character group of Grothendieck and Raynaud). If $\bar\rho$ arose only from
$q$-new forms, $J_0(N)[\mathfrak m]$ would be controlled by the toric part, on which inertia acts unipotently and $\Frob_q$ acts by $\pm q$ on the subobject and $\pm1$ on
the quotient, with the monodromy pairing measuring the extension; ``unramified at $q$'' together with $q\not\equiv1\pmod\ell$ (so that the two Frobenius
eigenvalues $\pm q$, $\pm1$ are distinct modulo $\ell$) forces $\mathfrak m$ to be in the support of the $q$-old part, whose Hecke eigensystems come from level $N/q$.
See Ribet \cite{Ribet90} (where the principle is attributed to Mazur), \cite{Mazur77} and \cite{DDT}. We do not prove it.
\end{proof}

\begin{theorem}[Ribet's level-lowering theorem]\label{thm:modgalrep:ribet}
Let $\ell$ be odd, $f$ a newform of weight $2$ on $\Gamma_0(N)$ with $\ell\nmid N$, $\lambda\mid\ell$ with $\bar\rho=\bar\rho_{f,\lambda}$ irreducible, and $q\,\|\,N$ a prime at which $\bar\rho$ is
unramified. Then $\bar\rho$ arises from a newform of weight $2$ on $\Gamma_0(N/q)$.
\end{theorem}

\begin{proof}[Idea of proof]
When $q\not\equiv1\pmod\ell$ this is Mazur's principle. When $q\equiv1\pmod\ell$, Ribet chooses an auxiliary prime $r\nmid N\ell$, $r\not\equiv1\pmod\ell$, at which $\bar\rho$ satisfies the
level-raising condition of \Cref{thm:modgalrep:level-raising}, raises the level to $Nr$, and compares the character groups of the toric parts of
$J_0(Nr)$ at $q$ and of the Jacobian of the Shimura curve attached to the quaternion algebra of discriminant $qr$ at $r$, using the uniformisations
of Cerednik and Drinfeld (``Ribet's exact sequence''); multiplicity one then transfers $\bar\rho$ to the Shimura curve, and back to level $Nr/q$, from
which Mazur's principle removes $r$. See Ribet \cite{Ribet90} (which also treats $\ell\mid N$ under conditions given there; the case stated here, $\ell\nmid N$, is
the one the paper uses), \cite{Ribet91} for a proof without multiplicity one, and \cite{DDT}. We do not prove it.
\end{proof}

\begin{remark}[Carayol's lemma on characters]\label{rem:modgalrep:carayol-lemma}
Carayol \cite{Carayol89} shows that a representation arising from an eigenform of weight $k\ge2$, level $M$ and character $\varepsilon$ also arises in the same
weight and level with any character $\varepsilon'$ having the same reduction modulo $\lambda$ as $\varepsilon$ (under hypotheses stated in the source, which we have
not reproduced). Taking $\varepsilon'$ to be the Teichm\"uller lift of $\bar\varepsilon$, which has order prime to $\ell$, explains the clause ``with nebentypus of order prime
to $\ell$'' in \Cref{thm:modgalrep:optimal-level}. The underlying mechanism is the lifting lemma of Deligne and Serre: a system of Hecke eigenvalues
modulo $\lambda$ occurring in the reduction of a Hecke-stable lattice lifts to an eigenform in characteristic $0$.
\end{remark}

\begin{theorem}[the optimal level; Ribet, Carayol, Diamond]\label{thm:modgalrep:optimal-level}
Let $\ell$ be odd and $\bar\sigma\colon G_\Q\to\GL_2(\Fbar_\ell)$ irreducible and modular, and if $\ell=3$ assume that $\bar\sigma|_{\Gal(\Qbar/\Q(\sqrt{-3}))}$ is absolutely irreducible.
Then $\bar\sigma$ arises from a newform of weight $2$ and level $N(\bar\sigma)\ell^{\delta(\bar\sigma)}$ with character of order prime to $\ell$, where $\delta(\bar\sigma)\in\{0,1\}$ and
$\delta(\bar\sigma)=0$ when $\bar\sigma|_{D_\ell}$ is good in the sense of \cite{DDT}, in particular when it is finite flat.
\end{theorem}

We do not prove this; it is \cite[Theorem 3.15]{DDT}, as the paper quotes it, combining Ribet \cite{Ribet90}, Carayol \cite{Carayol89} and Diamond
\cite{Diamond95}. Two remarks on the hypothesis at $\ell=3$, which is Diamond's. (a) For $\bar\sigma=\bar\rho_{E,3}$ irreducible it holds exactly when the image is not
$N_s(\F_3)$, the normaliser of a split Cartan subgroup (\Cref{ch:galrep}, the proposition on the image $\Ns$). (b) For $\ell\ge5$, $\ell\nmid N(E)$ and $\bar\rho_{E,\ell}$
irreducible, the analogous condition on $\Gal(\Qbar/\Q(\sqrt{\ell^*}))$ holds automatically (\Cref{ch:galrep}), although the theorem does not need it.

\begin{corollary}[the specialisation to elliptic curves]\label{cor:modgalrep:optimal-level-E}
Let $\ell$ be odd and $E'/\Q$ of conductor $N'$ with $\ell\nmid N'$, $\bar\sigma=\bar\rho_{E',\ell}$ irreducible and, if $\ell=3$, of image not $N_s(\F_3)$. Then there are a newform $g$ of
weight $2$ on $\Gamma_0(N(\bar\sigma))$ (trivial character) and a prime $\lambda\mid\ell$ of its coefficient field with $\bar\rho_{g,\lambda}\simeq\bar\sigma\otimes_{\F_\ell}k_\lambda$; so $a_p(g)\equiv a_p(E')\pmod\lambda$
for all $p\nmid\ell N'$. Moreover $N(\bar\sigma)\mid N'$, and $N(\bar\sigma)$ is prime to every prime at which $\bar\sigma$ is unramified.
\end{corollary}

\begin{proof}
$\bar\sigma$ is modular: $\bar\sigma\simeq\bar\rho_{f_{E'},\ell}$ by modularity of $E'$ and \Cref{thm:modgalrep:comparison-E}. It is odd, and irreducible, hence absolutely
irreducible (\Cref{prop:modgalrep:lattices}(ii)). At $\ell=3$ the hypothesis on the image is equivalent to absolute irreducibility on $\Gal(\Qbar/\Q(\sqrt{-3}))$
(\Cref{ch:galrep}). $E'$ has good reduction at $\ell$, so $\bar\sigma|_{D_\ell}=E'[\ell]$ is the generic fibre of the finite flat group scheme $\mathcal E'[\ell]$ of the N\'eron
model, and $\delta(\bar\sigma)=0$ (\Cref{thm:modgalrep:crystalline}(i); \Cref{ch:galrep}). \Cref{thm:modgalrep:optimal-level} gives a newform $g$ of weight $2$, level $N(\bar\sigma)$
and character $\psi$ of order prime to $\ell$, with $\bar\sigma\simeq\bar\rho_{g,\lambda}\otimes\Fbar_\ell$. By \Cref{thm:modgalrep:deligne} (weight $2$), $\det\bar\rho_{g,\lambda}=\bar\psi\bar\chi_\ell$, while
$\det\bar\sigma=\bar\chi_\ell$; so $\bar\psi=1$. The values of $\psi$ are roots of unity of order prime to $\ell$, and reduction modulo $\lambda$ is injective on them ($X^m-1$ is
separable over $\F_\ell$ for $\ell\nmid m$); so $\psi=1$ and $g$ is on $\Gamma_0(N(\bar\sigma))$. Both $\bar\rho_{g,\lambda}$ and $\bar\sigma\otimes k_\lambda$ are absolutely irreducible and become
isomorphic over $\Fbar_\ell$, so they have the same characteristic polynomials and are isomorphic over $k_\lambda$ (Brauer-Nesbitt, \Cref{ch:reps}); taking traces of
Frobenius gives the congruences. Finally $N(\bar\sigma)\mid N'$ by \Cref{prop:modgalrep:serre-conductor-divides} for $f_{E'}$, and the last claim is
\Cref{def:modgalrep:serre-invariants}(a).
\end{proof}

\begin{paperlink}
The optimal-level paragraph of Section 3: ``For $\bar\sigma=\bar\rho_{E',\ell}$ with $E'$ an elliptic curve of conductor prime to $\ell$, the restriction to $D_\ell$ comes from
the finite flat group scheme $E'[\ell]$ and is good, and $\det\bar\sigma=\bar\chi_\ell$ forces the nebentypus, which is trivial modulo $\lambda$ and of order prime to $\ell$,
to be trivial. So in that case $\bar\sigma$ arises from a newform of weight $2$ on $\Gamma_0(N(\bar\sigma))$.'' This is \Cref{cor:modgalrep:optimal-level-E}; its proof
spells out the two compressed steps (the determinant of $\bar\rho_g$, and injectivity of reduction on roots of unity of order prime to $\ell$). The
paper states the extra hypothesis at $\ell=3$ as absolute irreducibility on $\Gal(\Qbar/\Q(\sqrt{-3}))$, exactly as in \Cref{thm:modgalrep:optimal-level}.
\end{paperlink}

\begin{example}[level lowering on the paper's curves]\label{ex:modgalrep:level-lowering}
The Serre conductors below are computed by the formula $a_p(\bar\rho_{E,\ell})=f_p-\dim\Phig_p[\ell]$ of \Cref{ch:galrep}; each congruence $a_p(E)\equiv a_p(g)\pmod\ell$ (twisted by
$\chi$ where indicated) was checked with PARI/GP for all primes $p<1000$ not dividing $\ell N$.
\begin{center}\small
\begin{tabular}{@{}lllll@{}}
\toprule
$E$, $\ell$ & reduction types & $N(\bar\rho_{E,\ell})$ & newform & mechanism\\
\midrule
\curve{121c1}, $3$ & $\kod{IV}$ at $11$ & $11$ & $f_{11a}$ & additive $q^2\to q$ (Lemma lem:IV)\\
\curve{77c1}, $3$ & $\kod{I}_3$ at $7$, $\kod{I}_2$ at $11$ & $11$ & $f_{11a}$ & Ribet at $q=7\equiv1\pmod3$\\
\curve{52a1}, $3$ & $\kod{IV}^*$ at $2$, $\kod{I}_2$ at $13$ & $26$ & $f_{26b}$ & additive $q^2\to q$ at $q=2$\\
\curve{121a1}, $3$ & $\kod{II}$ at $11$ & $121$ & none lower & twist: $N(\bar\rho\otimes\chi_{-11})=11$, $f_{11a}$\\
\curve{539c1}, $3$ & $\kod{I}_3^*$ at $7$, $\kod{I}_2$ at $11$ & $539$ & none lower & twist: $N(\bar\rho\otimes\chi_{-7})=11$, $f_{11a}$\\
\curve{140b1}, $5$ & $\kod{IV}^*$ at $2$, $\kod{I}_1$ at $5$, $\kod{I}_5$ at $7$ & $4$ & $f_{20a}$, level $20$ & $5\mid N$, $\bar\rho|_{D_5}$ not finite\\
\curve{11a1}, $5$ & $\kod{I}_5$ at $11$ & $1$ & none & $\bar\rho$ reducible\\
\bottomrule
\end{tabular}
\end{center}
For \curve{52a1} the other newform of level $26$, $f_{26a}$, is not congruent to $f_{52a}$ modulo $3$; the lowering at $q=2$ uses that the image of $I_2$ on $T_3E$ is
cyclic of order $3$ for type $\kod{IV}^*$ at $2$ (\Cref{ch:galrep}), so $E[3]^{I_2}$ is a line and the Serre exponent at $2$ is $1$. For \curve{140b1} at $\ell=5$, $5\mid N$
and $\bar\rho|_{D_5}$ is not finite ($5\nmid v_5(\Delta)=1$), so \Cref{cor:modgalrep:optimal-level-E} does not apply; the newform has level $N(\bar\rho)\cdot5=20$ (there is
no newform of weight $2$ and level $4$), and indeed $a_p(f_{20a})\equiv a_p(\curve{140b1})\pmod5$ for $p\nmid70$, while at $p=7$ the congruence is the level-raising one,
$a_7(f_{20a})=2\equiv(7+1)a_7(\curve{140b1})=-8\pmod5$. For \curve{11a1} at $\ell=5$, $E[5]\simeq\Z/5\oplus\mu_5$ (image \texttt{5Cs.1.1}), so $a_p\equiv1+p\pmod5$
for all $p\ne5,11$; there is no newform of level $1$, consistent with \Cref{thm:modgalrep:optimal-level}, which needs irreducibility.
\end{example}

\begin{remark}[history]\label{rem:modgalrep:level-lowering-history}
Serre \cite{Serre87} conjectured that an odd irreducible $\bar\sigma$ arises from a newform of the optimal weight, level and character, and observed that the
``epsilon conjecture'' (lowering the level of the representation attached to the Frey curve of a putative solution of Fermat's equation)
would, together with modularity, prove Fermat's Last Theorem. Mazur's principle handled $q\not\equiv1\pmod\ell$; Ribet \cite{Ribet90} proved the epsilon
conjecture, and with it the implication. Carayol \cite{Carayol89} settled the conductor and character side, and Diamond \cite{Diamond95} the remaining
cases of the level, giving \Cref{thm:modgalrep:optimal-level}. Multiplicity one ($J_0(N)[\mathfrak m]$ two-dimensional for non-Eisenstein $\mathfrak m$, \Cref{ch:moddeg}) was used
in Ribet's original argument and removed in \cite{Ribet91}. Khare and Wintenberger proved the existence part of Serre's conjecture
(\Cref{thm:modgalrep:serre-conjecture}).
\end{remark}

\section{Raising the level}\label{sec:modgalrep:raising}

The converse operation adds a prime $q\nmid N\ell$ to the level. It is possible exactly when $a_q^2\equiv(q+1)^2$, and the quantity $(q+1)^2-a_q^2$ reappears in the
paper as the level-raising Euler factor of Theorem thm:dfg14: depleting $f$ at a good prime is ``raising its level'', and the congruence ideal grows
by the factor. The necessity of the congruence is also the step of Lemma lem:old at a prime $p\,\|\,N$ at which $\bar\rho$ is unramified, read backwards.

\begin{theorem}[Ribet's level-raising theorem]\label{thm:modgalrep:level-raising}
Let $f$ be a newform of weight $2$ on $\Gamma_0(N)$, $\lambda\mid\ell$ with $\ell$ odd and $\bar\rho=\bar\rho_{f,\lambda}$ irreducible, and $q\nmid N\ell$ a prime.
\begin{enumerate}[label=(\roman*)]
\item (Necessity.) If $\bar\rho$ arises from a newform $g$ of weight $2$ on $\Gamma_0(M)$ with $q\mid M\mid Nq$, then $a_q(f)\equiv(q+1)a_q(g)\pmod{\lambda'}$ for a prime $\lambda'\mid\lambda$ of a
common coefficient ring, with $a_q(g)=\pm1$; hence $a_q(f)^2\equiv(q+1)^2\pmod\lambda$.
\item (Sufficiency.) Conversely, if $a_q(f)^2\equiv(q+1)^2\pmod\lambda$, then $\bar\rho$ arises from a newform of weight $2$ on $\Gamma_0(M)$ with $q\mid M\mid Nq$.
\end{enumerate}
\end{theorem}

\begin{proof}
(i) $q\,\|\,M$, so $g$ has type (St) at $q$ (\Cref{thm:modgalrep:local-types}). As $\bar\rho_g\simeq\bar\rho$ is unramified at $q$, part (c) of that theorem says that
$\bar\rho(\Frob_q)$ has the eigenvalues $a_q(g)q$ and $a_q(g)$ modulo $\lambda'$; its trace is $a_q(f)$ modulo $\lambda'$ (as $q\nmid N\ell$). So $a_q(f)\equiv(q+1)a_q(g)$, and squaring, with
$a_q(g)^2=1$, $a_q(f)^2\equiv(q+1)^2$ modulo $\lambda'\cap\mathcal{O}_f=\lambda$. (ii) We do not prove this; it is Ribet's theorem \cite{Ribet84}, whose proof uses Ihara's lemma
(\Cref{lem:modgalrep:ihara}) to show that the image of the two degeneracy maps from level $N$ to level $Nq$ meets the $q$-new part in a module whose
Hecke eigenvalues modulo $\lambda$ are those of $f$ exactly when $a_q^2\equiv(q+1)^2$. Diamond and Taylor \cite{DiamondTaylor94} prove a general form, including the
raising of levels by $q^2$. We have not checked the precise hypotheses on $\ell$ in these two sources against the form stated here (the paper uses
only the necessity (i), proved above, and the level-raising factor of Theorem thm:dfg14).
\end{proof}

\begin{proposition}[the level-raising factor]\label{prop:modgalrep:level-raising-factor}
Let $p\nmid N$, and $\alpha$, $\beta$ the roots of $X^2-a_p(f)X+p$. Then
\[
\Bigl(1-\frac{\alpha}{\beta p}\Bigr)\Bigl(1-\frac1p\Bigr)\Bigl(1-\frac{\beta}{\alpha p}\Bigr)=\frac{p-1}{p}\cdot\frac{(p+1)^2-a_p(f)^2}{p^2},
\]
a nonzero element of $K_f$. It is $L_p(A_f,1)^{-1}=L_p^{\nv}(A_f,1)^{-1}$, the eigenvalues of either Frobenius on $\ad V$ being $\alpha/\beta$, $1$, $\beta/\alpha$. For $\ell\ne p$ and
$\lambda\mid\ell$, its $\lambda$-adic valuation is $v_\lambda(p-1)+v_\lambda\bigl((p+1)^2-a_p^2\bigr)$, and the second term is positive exactly when the level-raising condition of
\Cref{thm:modgalrep:level-raising} holds at $p$. At a prime $p\,\|\,N$ the corresponding factor is $L_p(A_f,1)^{-1}=1-p^{-2}$, and at $p^2\mid N$ the naive factor is $1$.
\end{proposition}

\begin{proof}
$\alpha\beta=p$ and $\alpha^2+\beta^2=a_p^2-2p$, so
\[
\Bigl(1-\frac{\alpha}{\beta p}\Bigr)\Bigl(1-\frac{\beta}{\alpha p}\Bigr)=1-\frac{\alpha^2+\beta^2}{\alpha\beta p}+\frac1{p^2}=1-\frac{a_p^2-2p}{p^2}+\frac1{p^2}=\frac{p^2+2p+1-a_p^2}{p^2}=\frac{(p+1)^2-a_p^2}{p^2},
\]
and multiplying by $1-\frac1p=\frac{p-1}p$ gives the identity. It is nonzero because $\abs{a_p}\le2\sqrt p<p+1$ under every complex embedding (\Cref{ch:modforms}). On $\ad V$ the
arithmetic Frobenius has eigenvalues $\alpha/\beta$, $1$, $\beta/\alpha$, an inversion-stable set, so the geometric Frobenius has the same ones, and
$L_p(A_f,s)^{-1}=(1-\frac\alpha\beta p^{-s})(1-p^{-s})(1-\frac\beta\alpha p^{-s})$; at $s=1$ this is the left side. The valuation follows since $p$ is a unit at $\lambda$. The factor at $p\,\|\,N$ is
\Cref{cor:modgalrep:delta-naive} at $s=1$.
\end{proof}

The term $v_\lambda(p-1)$ has a similar meaning, which we do not use: it is positive only when $p\equiv1\pmod\ell$, and then $\bar\rho$ can also arise at level $Np^2$ from forms
of type (RPS) at $p$ whose inertial characters have $\ell$-power order (they reduce to the trivial character); this is part of the theory of non-optimal
levels of Diamond and Taylor \cite{DiamondTaylor94}.

\begin{theorem}[Diamond, Flach and Guo; the paper's Theorem thm:dfg14]\label{thm:modgalrep:dfg14}
Let $f$ be a newform of weight $2$ on $\Gamma_0(M)$ with rational coefficients, $\lambda=\ell$ a prime not dividing $2M$ with $\bar\rho_f$ irreducible, and $\Sigma$ a finite
set of primes. Let $M^\Sigma=M\prod_{p\in\Sigma}p^{\delta_p}$, $f^\Sigma$ the $\Sigma$-depleted eigenform of level $M^\Sigma$, and $\eta^\Sigma_{f,\lambda}$ the congruence ideal of
\cite{DFG04} attached to $f^\Sigma$ (defined in \Cref{ch:congruence}). Then
\[
\eta^\Sigma_{f,\lambda}=\eta^\emptyset_{f,\lambda}\prod_{p\in\Sigma}L_p^{\nv}(A_f,1)^{-1}.
\]
\end{theorem}

We do not prove this here; it is \cite[Proposition 1.4(c)]{DFG04}, as the paper cites it, and \Cref{ch:congruence} explains the proof (Ihara's lemma
and the degeneracy maps). By \Cref{prop:modgalrep:level-raising-factor} and \Cref{cor:modgalrep:delta-naive}, adding to $\Sigma$ a prime with $\delta_p=0$ changes neither the level nor
the ideal, a good prime contributes $\frac{p-1}p\cdot\frac{(p+1)^2-a_p^2}{p^2}$, and a multiplicative prime $1-p^{-2}$. The reading for this chapter: $f^\Sigma$ is ``$f$ with its level
raised at $\Sigma$'' (at a good $p$ it is the eigenform of level $Mp^2$ in the span of $f(z)$, $f(pz)$, $f(p^2z)$ with $U_p$-eigenvalue $0$), and the change of the
congruence ideal is the level-raising factor.

\begin{paperlink}
The paragraph after Theorem thm:dfg14: ``adding a good prime $p$ multiplies the ideal by $(1-\frac\alpha\beta p^{-1})(1-p^{-1})(1-\frac\beta\alpha p^{-1})$, Wiles's level-raising
factor; adding a multiplicative prime multiplies it by $1-p^{-2}$.'' This is \Cref{prop:modgalrep:level-raising-factor}. The same identity, with $a_q(E')$ for the
twist $E'$ of good reduction, gives the $\kod{I}_0^*$ entry of the table of Section 1 and of Lemma lem:euler, and Step 3 of the proof of Theorem thm:quant.
\end{paperlink}

\begin{lemma}[Ihara's lemma]\label{lem:modgalrep:ihara}
Let $q\nmid N\ell$ be prime, $\alpha_1,\alpha_q\colon X_0(Nq)\to X_0(N)$ the two degeneracy maps ($z\mapsto z$ and $z\mapsto qz$), and $\mathfrak m$ a maximal ideal of $\T_N$ of residue
characteristic $\ell$ whose residual representation is irreducible. Then the map $H^1(X_0(N),\F_\ell)^2_{\mathfrak m}\to H^1(X_0(Nq),\F_\ell)$, $(x,y)\mapsto\alpha_1^*x+\alpha_q^*y$, is
injective. Equivalently (Ribet), the kernel of $J_0(N)^2\to J_0(Nq)$, $(x,y)\mapsto\alpha_1^*x+\alpha_q^*y$, is the Shimura subgroup $\mathcal S_N$ embedded antidiagonally,
which is Eisenstein.
\end{lemma}

We do not prove this; see Ribet \cite{Ribet84} and Diamond and Taylor \cite{DiamondTaylor94} (Ihara's original argument uses the congruence
subgroup property of $\SL_2(\Z[1/q])$). It is the input of \Cref{thm:modgalrep:level-raising}(ii) and of Theorem thm:dfg14 (\Cref{ch:congruence}).

\begin{example}[level raising for $f_{11a}$ and $f_{20a}$]\label{ex:modgalrep:level-raising}
(Computed with PARI/GP on the new subspace of level $11q$, resp.\ $20q$: each newform given by \texttt{mfeigenbasis}, with coefficient field from
\texttt{mffields}, was reduced modulo every prime above $\ell$ of that field and compared with $a_p(f)$ for all primes $p<200$, $p\nmid\ell Nq$.) For $f=f_{11a}$ and $\ell=3$:
\begin{center}\small
\begin{tabular}{@{}rrrcp{5.9cm}@{}}
\toprule
$q$ & $a_q(f_{11a})$ & $(q+1)^2-a_q^2$ & $3\mid$? & newforms of level $11q$ congruent to $f_{11a}$ modulo a prime above $3$\\
\midrule
$2$ & $-2$ & $5$ & no & none (no newforms of level $22$)\\
$5$ & $1$ & $35$ & no & none\\
$7$ & $-2$ & $60$ & yes & $f_{77c}$ (rational), $a_7(f_{77c})=-1$\\
$13$ & $4$ & $180$ & yes & one form with quartic coefficient field, at one of its two primes above $3$\\
$17$ & $-2$ & $320$ & no & none\\
$19$ & $0$ & $400$ & no & none\\
$23$ & $-1$ & $575$ & no & none\\
$29$ & $0$ & $900$ & yes & two Galois orbits (coefficient fields of degrees $3$ and $7$)\\
$31$ & $7$ & $975$ & yes & one Galois orbit (coefficient field of degree $11$)\\
$37$ & $3$ & $1435$ & no & none\\
\bottomrule
\end{tabular}
\end{center}
So, for $q<40$, $q\ne3,11$, the newforms congruent to $f_{11a}$ at level $11q$ occur exactly when $3\mid(q+1)^2-a_q^2$, as \Cref{thm:modgalrep:level-raising} predicts; at
$q=7$, $a_7(f_{11a})=-2\equiv(7+1)a_7(f_{77c})=-8\pmod3$. At level $143$ the congruent form is not the rational newform $f_{143a}$: $a_2(f_{143a})=0\not\equiv-2\pmod3$. For
$f=f_{20a}$ and $\ell=5$ ($\bar\rho_{20a,5}$ is irreducible): $q=3$ gives $12$ and there are no newforms of level $60$; $q=7$ gives $60$, and $f_{140b}$ is congruent to
$f_{20a}$ modulo $5$, with $a_7(f_{20a})=2\equiv8a_7(f_{140b})=-8$; $q=11$ and $q=13$ give $144$ and $192$, and no newform of level $220$ or $260$ is congruent. The
congruence $f_{11a}\equiv f_{77c}$ modulo $3$ underlies the paper's Example for \curve{539c1}, whose twist by $\chi_{-7}$ is \curve{77c1}.
\end{example}

\section[The congruence criterion and the oldform lemma]{Congruences between eigenforms: the congruence criterion and the oldform lemma}\label{sec:modgalrep:congruences}

This section proves the two lemmas with which Section 3 of the paper begins. Lemma lem:crit identifies the congruence number with the
congruence ideal of $\phi_f$ in the full Hecke algebra and turns a congruence between $f$ and an eigenform orthogonal to $f$, \emph{for every
coefficient}, into the divisibility $\ell\mid r$. Lemma lem:old produces such an eigenform from a newform of lower level that is congruent to $f$ only at
the primes $p\nmid\ell N$, by choosing the eigenvalues of the operators $U_p$ on the oldforms prime by prime. The gap between ``every $n$'' and ``almost
every prime $p$'' is the whole point of the second lemma. The paper's statements are proved here in full (the book's \Cref{lem:modgalrep:crit} adds
$I=\Ann_\T(J)$ and the description of all congruence moduli); \Cref{ch:moddeg} relates $r$ to the modular degree, and \Cref{ch:congruence} restates
Lemma lem:crit in the language of congruence modules.

\begin{definition}[the congruence number]\label{def:modgalrep:congruence-number}
Let $f$ be a newform of weight $2$ on $\Gamma_0(N)$ with rational coefficients (so $f=f_E$ for an elliptic curve $E$ of conductor $N$). $S_2(\Gamma_0(N);\Z)$ is
the lattice of forms with integral $q$-expansion (\Cref{ch:modforms}), $f^\perp\subset S_2(\Gamma_0(N))$ the orthogonal complement of $\C f$ for the Petersson product,
and $L=f^\perp\cap S_2(\Gamma_0(N);\Z)$. The \emph{congruence number} $r=r_f=r_E$ is the largest integer $r\ge1$ such that $f\equiv g\pmod r$ coefficientwise for some
$g\in L$; by \Cref{lem:modgalrep:crit} it exists, and the moduli of such congruences are exactly the divisors of $r$. It depends only on $f$, that is, on
the isogeny class.
\end{definition}

\begin{theorem}[Agashe, Ribet and Stein]\label{thm:modgalrep:ars}
Let $E_0$ be the optimal curve of conductor $N$ and $r=r_{E_0}$. Then $\deg\varphi_{E_0}\mid r$, and $v_p(\deg\varphi_{E_0})=v_p(r)$ for every prime $p$ with $p^2\nmid N$.
\end{theorem}

We do not prove this here; it is \cite[Theorem 2.2]{ARS12}, stated and discussed in \Cref{ch:moddeg}. Example (computed with PARI/GP as the index
of $\Z f\oplus L$ in $S_2(\Gamma_0(N);\Z)$, by \Cref{lem:modgalrep:crit}(i)): for \curve{99a1} (the 99A1 of \cite{ARS12}, $N=3^2\cdot11$), $\deg\varphi=4$ and $r=12$, so the
$3$-adic valuations differ, as $9\mid N$ allows; for \curve{539d1}, $\deg\varphi=72$ and $r=504=7\cdot72$, with $7^2\mid539$; for \curve{726a1}, $\deg\varphi=144$ and $r=1584=11\cdot144$, with $11^2\mid726$. For \curve{11a1}, \curve{52a1}, \curve{77c1},
\curve{121a1}, \curve{121c1}, \curve{539c1} the same computation gives $r=1$, $3$, $6$, $6$, $6$, $288$, equal to $\deg\varphi$ in each case.

\begin{lemma}[the congruence criterion; the paper's Lemma lem:crit]\label{lem:modgalrep:crit}
Let $f$, $L$, $r$ be as in \Cref{def:modgalrep:congruence-number}, $\T$ the full Hecke algebra of level $N$, $\phi_f\colon\T\to\Z$, $T_n\mapsto a_n(f)$, $I=\ker\phi_f$ and
$J=\Ann_\T(f^\perp)$.
\begin{enumerate}[label=(\roman*)]
\item $I\cap J=0$, $J=\Ann_\T(I)$, $I=\Ann_\T(J)$, and $\T/(I+J)$ is a finite cyclic group of order $r$; equivalently $r\Z=\phi_f(J)=\phi_f(\Ann_\T I)$, the congruence ideal of $\phi_f$.
Moreover $[S_2(\Gamma_0(N);\Z):\Z f\oplus L]=r$, and $f\equiv g\pmod m$ for some $g\in L$ if and only if $m\mid r$.
\item Let $\mathcal{O}$ be the ring of integers of a number field, $\lambda\mid\ell$ a prime of $\mathcal{O}$, and $g\in(S_2(\Gamma_0(N))\otimes\Qbar)\cap f^\perp$ a normalised eigenform of \emph{all} of $\T$
with coefficients in $\mathcal{O}$ such that $a_n(g)\equiv a_n(f)\pmod\lambda$ for \emph{every} $n\ge1$. Then $\ell\mid r$.
\item The same holds over $\Zl$: with $\T_\ell=\T\otimes\Zl$, $I_\ell=I\otimes\Zl$, $J_\ell=J\otimes\Zl$, one has $\T_\ell/(I_\ell+J_\ell)\simeq\Z/\ell^{v_\ell(r)}$, $\phi_f(\Ann_{\T_\ell}I_\ell)=\ell^{v_\ell(r)}\Zl$,
and $\ell^{v_\ell(r)}$ is the largest power $\ell^t$ such that $f\equiv h\pmod{\ell^t}$ for some $h\in L\otimes\Zl$, the forms in $f^\perp$ with coefficients in $\Zl$.
\end{enumerate}
\end{lemma}

\begin{proof}
Write $S=S_2(\Gamma_0(N);\Z)$, $S_\Q=S\otimes\Q$; then $S_2(\Gamma_0(N))=S\otimes\C$ and $S$ is $\T$-stable (\Cref{ch:modforms,ch:hecke}). $\T$ is a subring of $\End_\Z(S)$, so it is
free of finite rank over $\Z$ and $\T\otimes\Q\subseteq\End_\Q(S_\Q)$ acts faithfully.

\emph{Step 1: $\T\otimes\Q=\Q\times\T'$.} First, $f^\perp$ is $\T$-stable. The adjoint $t\mapsto t^\dagger$ for the Petersson product is additive and reverses products, so
it suffices that $T_n^\dagger$ ($(n,N)=1$) and $U_p^\dagger$ ($p\mid N$) have $f$ as an eigenvector, these operators generating $\T$ as a ring
(\Cref{prop:modgalrep:newforms-recap}(i)). By \Cref{ch:hecke}, $T_n^\dagger=T_n$ for $(n,N)=1$ and $U_p^\dagger=w_NU_pw_N$, where $w_Nf=\varepsilon_N(f)f$ with $\varepsilon_N(f)=\pm1$;
so $U_p^\dagger f=\varepsilon_N(f)^2a_p(f)f=a_p(f)f$. Hence for $t\in\T$ and $g\in f^\perp$, $\langle tg,f\rangle=\langle g,t^\dagger f\rangle=0$.
Next, $\T\otimes\Q$ is a finite-dimensional commutative $\Q$-algebra, so $\T\otimes\Q=\prod_{\mathfrak n}(\T\otimes\Q)_{\mathfrak n}$ over its maximal ideals, and $S_\Q=\bigoplus_{\mathfrak n}(S_\Q)_{\mathfrak n}$,
the factor $(\T\otimes\Q)_{\mathfrak n}$ acting faithfully on the generalised eigenspace $(S_\Q)_{\mathfrak n}$ (\Cref{ch:modules}). Let $\mathfrak n_f=\ker(\phi_f\otimes\Q)$. The generalised
eigenspace $(S_\Q)_{\mathfrak n_f}\otimes\C$ lies in the generalised eigenspace of the $T_p$, $p\nmid N$, for the eigenvalues $a_p(f)$, which is $V_f=\C f$
(\Cref{prop:modgalrep:newforms-recap}(iv), strong multiplicity one: $f$ has level $N$, so $V_f$ is spanned by $f$ alone). So $(S_\Q)_{\mathfrak n_f}=\Q f$, and the factor
$(\T\otimes\Q)_{\mathfrak n_f}$ embeds in $\End_\Q(\Q f)=\Q$; it is therefore $\Q$, acting through $\phi_f$. Let $\T'=\prod_{\mathfrak n\ne\mathfrak n_f}(\T\otimes\Q)_{\mathfrak n}$ and $W=\bigoplus_{\mathfrak n\ne\mathfrak n_f}(S_\Q)_{\mathfrak n}$:
then $\T\otimes\Q=\Q\times\T'$, $S_\Q=\Q f\oplus W$, and $\T'$ acts faithfully on $W$ with its unit acting as the identity. Let $e_f=(1,0)$. It projects $S_\Q$ onto $\Q f$
along $W$. Since $f^\perp$ is $\T$-stable, $e_f(f^\perp)\subseteq f^\perp\cap\C f=0$, so $f^\perp\subseteq\ker e_f=W\otimes\C$, and the two have dimension $\dim S_\Q-1$; hence
$f^\perp=W\otimes\C$ and $L=S\cap W$.

\emph{Step 2: the ideals.} In $\T\otimes\Q=\Q\times\T'$, $\phi_f(a,t')=a$. An element $(a,t')$ kills $f^\perp=W\otimes\C$ iff $t'=0$. Hence $I=\T\cap(0\times\T')$ and $J=\T\cap(\Q\times0)$, and
$I\cap J=0$. Since $\T$ is a lattice in $\T\otimes\Q$, $I\otimes\Q=0\times\T'$ and $J\otimes\Q=\Q\times0$ (tensoring with $\Q$ is exact), and $J\ne0$. An element $(a,t')$ kills $I$ iff
$t'\T'=0$ iff $t'=0$, so $\Ann_\T(I)=J$; it kills $J$ iff $a=0$, so $\Ann_\T(J)=I$. $\T/I\simeq\phi_f(\T)=\Z$ (as $\phi_f(1)=1$), and $\phi_f$ is injective on $J$; so
$\T/(I+J)\simeq\Z/\phi_f(J)$ is cyclic, and finite because $\phi_f(J)\ne0$. Let $r'=\#\T/(I+J)$. Both $I$ and $J$ are saturated in $\T$ (if $mt\in J$ with $t\in\T$, $m\ne0$, then $t\in\T\cap(\Q\times0)$),
so $\T/I$ and $\T/J$ are free $\Z$-modules.

\emph{Step 3: the forms killed by $I$ and by $J$.} The pairing $S\times\T\to\Z$, $(g,t)\mapsto a_1(tg)$, is perfect (Ribet \cite[Theorem 2.2]{Ribet83}, Abbes and Ullmo
\cite[Lemma 2.1]{AU96}; proved in \Cref{ch:hecke}). For an ideal $\mathfrak a$ of $\T$, the orthogonal of $\mathfrak a$ is $S[\mathfrak a]=\{g:\mathfrak ag=0\}$: if $a_1(tg)=0$ for all $t\in\mathfrak a$ then
$a_n(tg)=a_1(T_ntg)=0$ for all $n$, since $T_nt\in\mathfrak a$, so $tg=0$. Under $S\simeq\Hom(\T,\Z)$, $S[\mathfrak a]\simeq\Hom(\T/\mathfrak a,\Z)$ (\Cref{ch:modules}, the algebra-module duality). A form
is killed by $I$ iff it is killed by $I\otimes\Q=0\times\T'$, iff its $W$-component vanishes (the unit of $\T'$ is the identity on $W$); so $S[I]=S\cap\Q f=\Z f$, because $a_1(f)=1$.
A form is killed by $J$ iff it is killed by $e_f$, iff it lies in $W$; so $S[J]=S\cap W=L$.

\emph{Step 4: the index.} The sequence $0\to\T\to\T/I\times\T/J\to\T/(I+J)\to0$, with first map $t\mapsto(t\bmod I,t\bmod J)$ (injective as $I\cap J=0$) and second map $(x,y)\mapsto x-y$, is
exact. Applying $\Hom(\cdot,\Z)$ and using $\Hom(\T/(I+J),\Z)=0$, $\operatorname{Ext}^1(\T/I\times\T/J,\Z)=0$ (free modules) and $\#\operatorname{Ext}^1(A,\Z)=\#A$ for finite $A$, we get
that $\Hom(\T/I,\Z)\oplus\Hom(\T/J,\Z)\to\Hom(\T,\Z)$ is injective with cokernel of order $r'$ (\Cref{ch:modules}, dualising an inclusion of lattices). Under the
pairing this map is the sum map $\Z f\oplus L\to S$. So $[S:\Z f\oplus L]=r'$.

\emph{Step 5: the congruences.} The map $S/(\Z f\oplus L)\to\Q f/\Z f$, $g\mapsto e_fg$, is well defined ($e_f$ kills $L$ and fixes $f$) and injective (if $e_fg=cf$ with $c\in\Z$ then
$g-cf\in S\cap W=L$). So $S/(\Z f\oplus L)$ is a finite subgroup of $\Q f/\Z f\simeq\Q/\Z$, hence cyclic of order $r'$, generated by the class of some $h\in S$ with $e_fh=f/r'$. Then
$g_0=r'h-f$ satisfies $e_fg_0=0$, so $g_0\in L$, and $f=r'h-g_0\equiv-g_0\pmod{r'}$ with $-g_0\in L$. Conversely, if $f\equiv g_1\pmod m$ with $g_1\in L$, then $h_1=(f-g_1)/m\in S$
and $e_fh_1=f/m$, whose class in $\Q f/\Z f$ has order exactly $m$; so $m\mid r'$. Hence the moduli of congruences between $f$ and elements of $L$ are exactly the
divisors of $r'$, the largest is $r=r'$, and (i) is proved.

\emph{Step 6: (ii).} Let $\phi_g\colon\T\to\mathcal{O}$, $t\mapsto$ the eigenvalue of $t$ on $g$; it is a ring homomorphism with $\phi_g(T_n)=a_n(g)\in\mathcal{O}$
(\Cref{prop:modgalrep:newforms-recap}(i), $a_1(g)=1$). It kills $J$: for $t\in J$, $\phi_g(t)g=tg=0$ as $g\in f^\perp\otimes\Qbar$. The two ring homomorphisms $\T\to\mathcal{O}/\lambda$ obtained
from $\phi_g$ and $\phi_f$ agree on every $T_n$, hence on $\T$; so $\phi_g(t)\equiv\phi_f(t)=0\pmod\lambda$ for $t\in I$. Hence $I+J$ lies in the kernel $\mathfrak m$ of $\T\to\mathcal{O}/\lambda$,
$t\mapsto\phi_g(t)\bmod\lambda$. Its image is a subring of the finite field $\mathcal{O}/\lambda$, hence a finite integral domain, hence a field of characteristic $\ell$; so $\mathfrak m$ is a
maximal ideal with $\ell\in\mathfrak m$, and $\T/(I+J)$ surjects onto $\T/\mathfrak m\ne0$, a group of exponent $\ell$. So $\ell\mid\#\T/(I+J)=r$.

\emph{Step 7: (iii).} $\Zl$ is flat over $\Z$, so $I_\ell=\ker(\phi_f\otimes\Zl)$, $\T_\ell/(I_\ell+J_\ell)=(\T/(I+J))\otimes\Zl\simeq\Z/\ell^{v_\ell(r)}$, $\Ann_{\T_\ell}(I_\ell)=J_\ell$ (annihilators of finitely
generated ideals commute with flat base change), and $\phi_f(J_\ell)=r\Zl=\ell^{v_\ell(r)}\Zl$. The pairing $S\otimes\Zl\times\T_\ell\to\Zl$ is perfect, $L\otimes\Zl=(S\otimes\Zl)\cap(W\otimes\Ql)$ (a
saturated sublattice stays saturated), and Steps 3 to 5 hold verbatim over $\Zl$, with $\Q_\ell f/\Zl f\simeq\Ql/\Zl$ in Step 5, whose finite subgroups are cyclic; the
index is now $\ell^{v_\ell(r)}$ (\Cref{ch:modules}, flat base change).
\end{proof}

\begin{paperlink}
Lemma lem:crit of the paper states (i) without $I=\Ann_\T(J)$ and the list of moduli, (ii), and (iii) in the words ``The same statements hold with $\Z$
replaced by $\Z_\ell$ and $r$ by its $\ell$-part''. Its proof is Steps 1 to 6 in compressed form: ``the $f$-line is $\T$-stable, and so is $f^\perp$, because $f$ is an
eigenvector of the adjoints of the $T_n$ as well'' is the first paragraph of Step 1; the decomposition $\T\otimes\Q=\Q\times\T'$ needs only that the
$\phi_f$-generalised eigenspace is the line of $f$, which is strong multiplicity one at level $N$ (the full Hecke algebra may contain nilpotents, but not
on that line); the finiteness of $\T/(I+J)$ is $\phi_f(J)\ne0$ (Step 2), which the paper's ``cyclic, of order $r'$'' takes for granted. The hypothesis that $g$ be an
eigenform of all of $\T$ and congruent to $f$ at every $n$, the operators $U_p$ ($p\mid N$) included, is used in Step 6 and cannot be weakened to the primes
$p\nmid\ell N$: that weakening is what \Cref{lem:modgalrep:old} repairs.
\end{paperlink}

\begin{lemma}[the operators $U_p$ on the span of an oldform]\label{lem:modgalrep:Up-oldspace}
Let $h$ be a newform of weight $2$ on $\Gamma_0(M)$, $M\mid N$, $V_h$ the span of the $h(dz)$, $d\mid N/M$, and for $p\mid N/M$ let $e_p=v_p(N/M)$ and $V_{h,p}$ the space with
basis $h_j$, $0\le j\le e_p$ (to be thought of as $h(p^jz)$).
\begin{enumerate}[label=(\roman*)]
\item $\bigotimes_{p\mid N/M}V_{h,p}\to V_h$, $\otimes_ph_{j_p}\mapsto h(\prod_pp^{j_p}z)$, is an isomorphism under which $U_p$ ($p\mid N/M$) acts on the $p$-th factor alone by $U_ph_j=h_{j-1}$
($j\ge1$) and $U_ph_0=a_p(h)h_0-p\,h_1$ if $p\nmid M$, $U_ph_0=a_p(h)h_0$ if $p\mid M$; $T_n$ ($(n,N)=1$) acts by $a_n(h)$, and $U_p$ for $p\mid M$ with $v_p(M)=v_p(N)$ by $a_p(h)$.
\item On $V_{h,p}$ the characteristic polynomial of $U_p$, its eigenvalues and one eigenvector for each are as follows, where $\alpha_p,\beta_p$ are the roots of
$X^2-a_p(h)X+p$:
\begin{center}\small
\begin{tabular}{@{}llll@{}}
\toprule
situation & characteristic polynomial & eigenvalue & eigenvector\\
\midrule
$p\nmid M$ & $X^{e_p-1}(X^2-a_pX+p)$ & $\alpha_p$; $\beta_p$ & $h_0-\beta_ph_1$; $h_0-\alpha_ph_1$\\
$p\nmid M$, $e_p\ge2$ & & $0$ & $h_0-a_p(h)h_1+p\,h_2$\\
$p\,\|\,M$ & $X^{e_p}(X-a_p)$ & $a_p(h)=\pm1$; $0$ & $h_0$; $h_0-a_p(h)h_1$\\
$p^2\mid M$ & $X^{e_p+1}$ & $0$ & $h_0$ ($a_p(h)=0$)\\
\bottomrule
\end{tabular}
\end{center}
$U_p$ is not semisimple on $V_{h,p}$ when $0$ is a repeated root of its characteristic polynomial.
\item For every choice of one eigenvalue $\mu_p$ from the table for each $p\mid N/M$, the image $g$ of the tensor product of the corresponding eigenvectors is
an eigenform of all of $\T$ with $U_pg=\mu_pg$ ($p\mid N/M$), $U_pg=a_p(h)g$ (other $p\mid N$), $T_ng=a_n(h)g$ ($(n,N)=1$); it has $a_1(g)=1$, and its coefficients lie in
$\mathcal{O}[\alpha_p,\beta_p:p\mid N/M]$ if those of $h$ lie in $\mathcal{O}$.
\end{enumerate}
\end{lemma}

\begin{proof}
(i) The $h(dz)$ are linearly independent (\Cref{prop:modgalrep:newforms-recap}(iv)). On $q$-expansions $a_n(U_pF)=a_{pn}(F)$; for $F=h(p^jd'z)$ with $p\nmid d'$ this gives
$U_pF=h(p^{j-1}d'z)$ if $j\ge1$, and $(U_ph)(d'z)$ if $j=0$, so $U_p$ acts on the $p$-th factor only. For $j=0$: if $p\mid M$, $U_p$ at level $N$ and at level $M$ have the
same effect on $q$-expansions, so $U_ph=a_p(h)h$; if $p\nmid M$, $a_n(T_ph)=a_{pn}(h)+p\,a_{n/p}(h)$ (the Hecke operator at level $M$), so $U_ph=T_ph-p\,h(pz)=a_p(h)h-p\,h(pz)$.
The statements about $T_n$ and about $U_p$ with $e_p=0$ are \Cref{prop:modgalrep:newforms-recap}(iv) and \Cref{ch:hecke}. (ii) In the basis $h_0,\dots,h_{e}$ ($e=e_p$) the
matrix of $U_p$ is a companion-type matrix; expanding $\det(X-U_p)$ along the last column gives the polynomials in the table. The eigenvectors are
checked directly: $U_p(h_0-\beta h_1)=a_ph_0-p\,h_1-\beta h_0=\alpha h_0-\alpha\beta h_1=\alpha(h_0-\beta h_1)$, using $a_p=\alpha+\beta$, $p=\alpha\beta$; $U_p(h_0-a_ph_1+p\,h_2)=a_ph_0-p\,h_1-a_ph_0+p\,h_1=0$;
$U_p(h_0-a_ph_1)=a_ph_0-a_ph_0=0$ when $p\mid M$. If $0$ is a repeated root, the $0$-eigenspace is the line spanned by the listed eigenvector (the matrix has
rank $e$), so $U_p$ is not diagonalisable. This is the shift-operator model of \Cref{ch:modules}; compare \Cref{ch:hecke}. (iii) A tensor product of
eigenvectors of operators acting on separate factors is a common eigenvector; $\T$ is generated by the $T_n$, $(n,N)=1$, and the $U_p$, $p\mid N$, which act as
stated. The coefficient of $h_0$ in each listed eigenvector is $1$, and $h(dz)$ with $d>1$ has $a_1=0$; so $a_1(g)=1$.
\end{proof}

\begin{lemma}[oldforms; the paper's Lemma lem:old]\label{lem:modgalrep:old}
Let $f=f_E$ be as in \Cref{def:modgalrep:congruence-number}, $\ell$ an odd prime with $\ell\nmid N$, and $\bar\rho=\bar\rho_{E,\ell}$ irreducible. Let $h\ne f$ be a newform of weight $2$
on $\Gamma_0(M)$, $M\mid N$, with coefficients in the ring of integers $\mathcal{O}$ of a number field, and $\lambda\mid\ell$ a prime of $\mathcal{O}$ with $a_p(h)\equiv a_p(f)\pmod\lambda$ for all
primes $p\nmid\ell N$. Assume that $p\nmid M$ for every prime $p\,\|\,N$ at which $\bar\rho$ is unramified. Then, after enlarging $\mathcal{O}$ and $\lambda$, the span $V_h$ of the $h(dz)$,
$d\mid N/M$, contains a normalised eigenform $g$ of all of $\T$ with coefficients in $\mathcal{O}$ and $a_n(g)\equiv a_n(f)\pmod\lambda$ for every $n\ge1$; it lies in $f^\perp\otimes\Qbar$.
Consequently $\ell\mid r$.
\end{lemma}

\begin{proof}
\emph{Step 1: orthogonality.} $V_h$ is $\T$-stable and orthogonal to $f$: if $M<N$ it is a different summand of \Cref{prop:modgalrep:newforms-recap}(iv) from $V_f=\C f$, and
if $M=N$ then $V_h=\C h$ with $h\ne f$. So any $g\in V_h\otimes\Qbar$ lies in $f^\perp\otimes\Qbar$.

\emph{Step 2: the residual representations.} By \Cref{thm:modgalrep:comparison-E}, $\bar\rho_f\simeq\bar\rho\otimes_{\F_\ell}k_\lambda$, which is absolutely irreducible
(\Cref{prop:modgalrep:lattices}(ii)). By \Cref{cor:modgalrep:congruent-newforms}(ii), $\bar\rho_h\simeq\bar\rho\otimes k_\lambda$.

\emph{Step 3: the choice of the $\mu_p$.} Enlarge $\mathcal{O}$ to contain the roots $\alpha_p,\beta_p$ of $X^2-a_p(h)X+p$ for the primes $p\mid N/M$ with $p\nmid M$, and replace $\lambda$ by a
prime above it (a congruence modulo the old $\lambda$ implies one modulo the new). For each $p\mid N/M$ choose $\mu_p$ as below, and let $g$ be the eigenform of
\Cref{lem:modgalrep:Up-oldspace}(iii); its coefficients lie in $\mathcal{O}$ and $a_1(g)=1$. We check $a_p(g)\equiv a_p(f)\pmod\lambda$ at every prime $p$.
\begin{itemize}
\item $p\nmid\ell N$: $a_p(g)=a_p(h)\equiv a_p(f)$ by hypothesis.
\item $p=\ell$: $\ell\nmid N$, so $a_\ell(g)=a_\ell(h)$; $f$ and $h$ have weight $2$ and level prime to $\ell$, and $\bar\rho_h\simeq\bar\rho_f$ (Step 2), so $a_\ell(h)\equiv a_\ell(f)$ by
\Cref{thm:modgalrep:deligne-fontaine}(iii).
\item $p^2\mid N$: $a_p(f)=0$ (\Cref{thm:modgalrep:local-types}). If $v_p(M)=v_p(N)$, then $p^2\mid M$ and $a_p(g)=a_p(h)=0$ by the same theorem applied to $h$. Otherwise $e_p\ge1$, and
$e_p\ge2$ if $p\nmid M$; in every case $0$ is in the list of \Cref{lem:modgalrep:Up-oldspace}(ii): take $\mu_p=0$, so $a_p(g)=0=a_p(f)$.
\item $p\,\|\,N$, $\bar\rho$ unramified at $p$: by hypothesis $p\nmid M$, so $e_p=1$ and the choices are $\alpha_p$, $\beta_p$. By \Cref{thm:modgalrep:local-types}(c) for $f$ (type (St)),
$\bar\rho(\Frob_p)$ has eigenvalues $a_p(f)p$, $a_p(f)$. Since $p\nmid M\ell$, $\bar\rho_h(\Frob_p)$ has characteristic polynomial $X^2-a_p(h)X+p$ modulo $\lambda$
(\Cref{prop:modgalrep:lattices}(i)), and $\bar\rho_h\simeq\bar\rho_f$; so $(X-\alpha_p)(X-\beta_p)\equiv(X-a_p(f)p)(X-a_p(f))\pmod\lambda$. Reducing modulo $\lambda$ (where $\mathcal{O}/\lambda$ is a field),
one of $\alpha_p,\beta_p$ is congruent to $a_p(f)$; take it as $\mu_p$.
\item $p\,\|\,N$, $\bar\rho$ ramified at $p$: then $p\mid M$, since otherwise $\bar\rho_h$ would be unramified at $p$ (\Cref{thm:modgalrep:eichler-shimura}(ii)(a), $p\ne\ell$); so
$v_p(M)=1=v_p(N)$ and $a_p(g)=a_p(h)=\pm1$. By \Cref{thm:modgalrep:local-types}(a), (b) for $f$ and for $h$ (both of type (St) at $p$), $\bar\rho_f|_{D_p}$ and $\bar\rho_h|_{D_p}$
are extensions of $\bar\varepsilon_f$ by $\bar\varepsilon_f\bar\chi_\ell$, resp.\ of $\bar\varepsilon_h$ by $\bar\varepsilon_h\bar\chi_\ell$, with $\bar\varepsilon_f(\Frob_p)=a_p(f)$, $\bar\varepsilon_h(\Frob_p)=a_p(h)$, and both are ramified,
hence non-split with a unique stable line. An isomorphism $\bar\rho_f|_{D_p}\simeq\bar\rho_h|_{D_p}$ carries the stable line to the stable line, hence induces an isomorphism of
the quotients: $\bar\varepsilon_f=\bar\varepsilon_h$, and $a_p(h)\equiv a_p(f)$. (This holds also when $p\equiv1\pmod\ell$, where the two characters of each extension coincide.)
\end{itemize}

\emph{Step 4: all $n$.} $f$ and $g$ are normalised eigenforms of all of $\T$, so both coefficient sequences satisfy $a_{mn}=a_ma_n$ for $(m,n)=1$, $a_{p^k}=a_p^k$ for $p\mid N$
($U_{p^k}=U_p^k$), and $a_{p^{k+1}}=a_pa_{p^k}-p\,a_{p^{k-1}}$ for $p\nmid N$, $p=\ell$ included (\Cref{prop:modgalrep:newforms-recap}(i)). Congruence at every prime therefore gives
congruence at every $n$, by induction on $n$. \Cref{lem:modgalrep:crit}(ii) now gives $\ell\mid r$.
\end{proof}

\begin{paperlink}
The paper's proof of Lemma lem:old follows the same plan and the same case division; the eigenvalue lists and eigenvectors are those of
\Cref{lem:modgalrep:Up-oldspace}. Two of its sentences are expanded here. ``$\bar\rho_h\simeq\bar\rho$ over a common finite extension of $\F_\ell$, by Chebotarev and
Brauer-Nesbitt, $\bar\rho$ being irreducible, hence absolutely irreducible'' is Step 2 (the original version of the paper had only ``$\bar\rho$ being irreducible'',
which gives only the semisimplifications). ``A non-split extension of two characters has a unique stable line'' is \Cref{thm:modgalrep:local-types}(b) with
\Cref{ch:reps}. The weight-two fact that $\alpha_p\ne\beta_p$ (Coleman and Edixhoven \cite{ColemanEdixhoven}) is not needed: the argument uses only that each listed
eigenvalue has an eigenvector.
\end{paperlink}

\begin{example}[Lemma lem:old on three pairs]\label{ex:modgalrep:old-examples}
(All congruences checked with PARI/GP for $n\le3000$, beyond the Sturm bounds $16$, $14$, $22$ at levels $77$, $52$, $121$, so they hold for every $n$ by \Cref{ch:modforms}.)
\begin{enumerate}[label=(\alph*)]
\item \emph{A prime $p\,\|\,N$ at which $\bar\rho$ is unramified.} $f=f_{77c}$, $\ell=3$, $h=f_{11a}$, $M=11$, $N/M=7$. \curve{77c1} has type $\kod{I}_3$ at $7$, and $3\mid n_7=3$, so $\bar\rho_{77c1,3}$ is
unramified at $7$ (\Cref{ch:galrep}, Lemma lem:tamagawa-exp); at $11$ it has type $\kod{I}_2$ and $\bar\rho$ is ramified, and $11\mid M$. Here $a_7(h)=-2$, $X^2+2X+7$ has roots
$\alpha,\beta=-1\pm\sqrt{-6}$, both congruent to $-1=a_7(f_{77c})$ modulo the prime $(3,\sqrt{-6})$ of $\Z[\sqrt{-6}]$, the ring of integers of $\Q(\sqrt{-6})$, above $3$ (which ramifies). So
$g=h-\beta h(7z)$ has $U_7$-eigenvalue $\alpha$ and $g\equiv f_{77c}$ modulo $(3,\sqrt{-6})$ for every $n$: modulo that prime $g$ has coefficients $a_n(f_{11a})+a_{n/7}(f_{11a})$. Hence
$3\mid r(\curve{77c1})$; indeed $r=\deg\varphi=6$.
\item \emph{An additive prime, at $q=2$.} $f=f_{52a}$, $\ell=3$, $h=f_{26b}$, $M=26$, $N/M=2$, $2\,\|\,M$, $e_2=1$. Here $\mu_2=0$, $g=h-a_2(h)h(2z)=f_{26b}(z)-f_{26b}(2z)$ ($a_2(f_{26b})=1$), and
$g\equiv f_{52a}\pmod3$ for every $n$. At $13$, $\bar\rho_{52a1,3}$ is ramified ($\kod{I}_2$) and $13\mid M$. Hence $3\mid r(\curve{52a1})$; indeed $r=\deg\varphi=3$. The other newform of level $26$,
$f_{26a}$, is not congruent to $f_{52a}$ modulo $3$.
\item \emph{An additive prime, at $q=11$.} $f=f_{121c}$, $\ell=3$, $h=f_{11a}$, $M=11$, $N/M=11$, $11\,\|\,M$, $e_{11}=1$: $\mu_{11}=0$, $g=f_{11a}(z)-f_{11a}(11z)$ ($a_{11}(f_{11a})=1$), the
$11$-depletion of $f_{11a}$, and $g\equiv f_{121c}\pmod3$ for every $n$; $r(\curve{121c1})=\deg\varphi=6$.
\end{enumerate}
\end{example}

\section[Proof of Theorem thm:qual]{Proof of Theorem thm:qual; the Example; the remark on \texorpdfstring{$\ell\mid N$}{l | N}}\label{sec:modgalrep:qual}

This section proves the paper's Theorem thm:qual in full, case by case, with the inputs of the previous sections: lower the level of $\bar\rho$ (case
(a)) or of $\bar\rho\otimes\chi_{q^*}$ (cases (b), (c)) to the Serre conductor, twist back to a newform $h\ne f$ of level dividing $N$, and apply Lemma lem:old and
Lemma lem:crit; the theorem of Agashe, Ribet and Stein passes from the congruence number to the modular degree. It then works the paper's
Example and the remark that follows it.

\begin{theorem}[the paper's Theorem thm:qual]\label{thm:modgalrep:qual}
Let $E/\Q$ have conductor $N$ and let $\ell$ be an odd prime not dividing $N$ such that $\bar\rho_{E,\ell}$ is irreducible and, if $\ell=3$, absolutely irreducible when
restricted to $\Gal(\Qbar/\Q(\sqrt{-3}))$ (equivalently, the mod $3$ image is not $N_s(\F_3)$). Let $q\ne\ell$ be a prime at which $E$ has additive reduction of one of the
following kinds:
\begin{enumerate}[label=(\alph*)]
\item type $\kod{IV}$ or $\kod{IV}^*$, with $q\ge5$ or $q=2$, and $\ell=3$;
\item type $\kod{II}$ or $\kod{II}^*$, with $q\ge5$ and $\ell=3$;
\item type $\kod{I}_n^*$ with $\ell\mid n$, $q$ odd.
\end{enumerate}
Then $\ell$ divides the modular degree of the $X_0(N)$-optimal curve $E_0$ in the isogeny class of $E$.
\end{theorem}

\subsection{The proof}\label{sec:modgalrep:proof-qual}

\begin{proof}
Let $f=f_E$, $E_0$ the optimal curve and $r=r_{E_0}=r_f$ (\Cref{def:modgalrep:congruence-number}); write $\bar\rho=\bar\rho_{E,\ell}$, which is irreducible, hence absolutely irreducible
(\Cref{prop:modgalrep:lattices}(ii)). Since $\ell\nmid N$, $\ell^2\nmid N$, and \Cref{thm:modgalrep:ars} gives $v_\ell(\deg\varphi_{E_0})=v_\ell(r)$; so it suffices to prove $\ell\mid r$. By
\Cref{lem:modgalrep:old} it suffices to find a newform $h\ne f$ of weight $2$ on $\Gamma_0(M_h)$, $M_h\mid N$, and a prime $\lambda\mid\ell$ of its coefficient ring with
\begin{enumerate}[label=(H\arabic*)]
\item $a_p(h)\equiv a_p(f)\pmod\lambda$ for all primes $p\nmid\ell N$, and
\item $p\nmid M_h$ for every prime $p\,\|\,N$ at which $\bar\rho$ is unramified.
\end{enumerate}
In all three cases $v_q(N)=2$: types $\kod{IV}$, $\kod{IV}^*$, $\kod{II}$, $\kod{II}^*$ at $q\ge5$ and $\kod{IV}$, $\kod{IV}^*$ at $q=2$ have conductor exponent $2$ (\Cref{ch:galrep}, Lemma lem:IV
and its corollary at $q=2$), and so has type $\kod{I}_n^*$ at an odd prime (\Cref{prop:modgalrep:twist-conductor}, row (RSt), with \Cref{prop:modgalrep:kodaira-dictionary}).

\emph{Case (a).} $\ell=3$, type $\kod{IV}$ or $\kod{IV}^*$ at $q$, $q\ge5$ or $q=2$.
\begin{enumerate}[label=(a\arabic*)]
\item By Lemma lem:IV (\Cref{ch:galrep}; at $q=2$ by its corollary there) the image of $I_q$ on $E[3]$ is a nontrivial unipotent group, $E[3]^{I_q}$ is a line, $\bar\rho$ is
tame at $q$, and the Serre exponent of $\bar\rho$ at $q$ is $2-1=1$.
\item Let $M=N(\bar\rho)$. Then $M\mid N$ (\Cref{prop:modgalrep:serre-conductor-divides}) and $v_q(M)=1<2=v_q(N)$, so $M\mid N/q$ and $M<N$.
\item \Cref{cor:modgalrep:optimal-level-E} (with $E'=E$; its hypotheses are those of the theorem) gives a newform $g$ of weight $2$ on $\Gamma_0(M)$ and $\lambda\mid3$ with
$\bar\rho_{g,\lambda}\simeq\bar\rho\otimes k_\lambda$, so $a_p(g)\equiv a_p(f)\pmod\lambda$ for $p\nmid3N$.
\item Take $h=g$, $M_h=M$. Then $h\ne f$, since the levels differ (\Cref{prop:modgalrep:newforms-recap}(iii)); (H1) holds; and (H2) holds because $N(\bar\rho)$ is prime to every prime at
which $\bar\rho$ is unramified. In Lemma lem:old the prime $q$ falls under ``$p^2\mid N$, $v_q(M)<v_q(N)$'', and the eigenvalue $\mu_q=0$ of $U_q$ is used.
\end{enumerate}

\emph{Case (b).} $\ell=3$, type $\kod{II}$ or $\kod{II}^*$ at $q\ge5$. Let $\chi=\chi_{q^*}$ and $E'=E^{(q)}$.
\begin{enumerate}[label=(b\arabic*)]
\item By Lemma lem:twist (i), (ii) (\Cref{ch:reduction}), $E'$ has type $\kod{IV}^*$, resp.\ $\kod{IV}$, at $q$, the same types as $E$ at the other primes, and conductor $N$; in particular
$3\nmid N(E')$.
\item $\bar\rho_{E',3}\simeq\bar\rho\otimes\chi$ (\Cref{ch:ellcurves}). It is irreducible, and, restricted to $\Gal(\Qbar/\Q(\sqrt{-3}))$, absolutely irreducible, since twisting by a character
preserves both properties.
\item By Lemma lem:IV for $E'$, $M=N(\bar\rho\otimes\chi)$ has $v_q(M)=1$, and $M\mid N(E')=N$; so $M\mid N/q$.
\item \Cref{cor:modgalrep:optimal-level-E} for $E'$ gives a newform $g$ on $\Gamma_0(M)$ and $\lambda\mid3$ with $\bar\rho_{g,\lambda}\simeq(\bar\rho\otimes\chi)\otimes k_\lambda$, so $a_p(g)\equiv\chi(p)a_p(f)\pmod\lambda$ for
$p\nmid3N$.
\item As $q\,\|\,M$, \Cref{cor:modgalrep:atkin-li} says that $h=g\otimes\chi=g^\chi$ is a newform on $\Gamma_0(Mq)$, and $Mq\mid N$. Its coefficients are $\pm$ those of $g$ or $0$, so $\lambda$
serves for $h$, and for $p\nmid3N$, $a_p(h)=\chi(p)a_p(g)\equiv\chi(p)^2a_p(f)=a_p(f)\pmod\lambda$: (H1).
\item $h\ne f$: otherwise $g=(g^\chi)^\chi=f^\chi=f_{E'}$ (\Cref{prop:modgalrep:twists}(iii), (iv)), which has level $N\ne M$.
\item (H2): if $p\,\|\,N$ and $\bar\rho$ is unramified at $p$, then $p\ne q$ (as $q^2\mid N$), so $\bar\rho\otimes\chi$ is unramified at $p$, $p\nmid M$, and $p\nmid Mq$.
\end{enumerate}

\emph{Case (c).} Type $\kod{I}_n^*$ at an odd $q\ne\ell$ with $\ell\mid n$. Let $\chi=\chi_{q^*}$ and $E'=E^{(q)}$.
\begin{enumerate}[label=(c\arabic*)]
\item By Lemma lem:twist(ii), $E'$ has type $\kod{I}_n$ at $q$, and by Lemma lem:twist(i) the same types as $E$ elsewhere; so its conductor is $N/q$, and $\ell\nmid N/q$.
\item By Lemma lem:tamagawa-exp for $E'$ (\Cref{ch:galrep}), $\bar\rho_{E',\ell}=\bar\rho\otimes\chi$ is unramified at $q$, because $\ell\mid n=n_q(E')$.
\item $M=N(\bar\rho\otimes\chi)$ is prime to $q$ and divides $N(E')=N/q$; so $M\mid N/q^2$.
\item \Cref{cor:modgalrep:optimal-level-E} for $E'$ (irreducibility and the condition at $\ell=3$ are twist-invariant, as in (b2)) gives a newform $g$ on $\Gamma_0(M)$ and $\lambda\mid\ell$
with $a_p(g)\equiv\chi(p)a_p(f)\pmod\lambda$ for $p\nmid\ell N$. At $q$ this lowering is the instance of Ribet's theorem (\Cref{thm:modgalrep:ribet}) in which $q$ exactly divides the
level $N/q$ of $E'$ and $\bar\rho\otimes\chi$ is unramified at $q$.
\item As $q\nmid M$, \Cref{cor:modgalrep:atkin-li} says that $h=g\otimes\chi=g^\chi$ is a newform on $\Gamma_0(Mq^2)$, and $Mq^2\mid N$; (H1) follows as in (b5).
\item $h\ne f$: otherwise $g=f^\chi=f_{E'}$, of level $N/q\ne M$ (as $q\mid N/q$ and $q\nmid M$).
\item (H2) holds as in (b7).
\end{enumerate}
In each case \Cref{lem:modgalrep:old} gives $\ell\mid r$, and so $\ell\mid\deg\varphi_{E_0}$.
\end{proof}

\begin{paperlink}
The paper's proof (Section sec:proof) is the one above. Three remarks on its wording, all about the twist dictionary of
\Cref{conv:modgalrep:twists}. (1) Case (b): ``its newform is $f'=f\otimes\chi$, a newform of weight $2$ on $\Gamma_0(N)$ (both have $a_q=0$)''. The naive twist is the newform
twist here because $a_q(f^\chi)=0$, which holds since $q^2\mid N$ (\Cref{thm:modgalrep:local-types}); Atkin and Li's theorem does not apply at this level ($q^2\mid N$), and is
used only for $h=g\otimes\chi$, where $q\,\|\,M$ (``As $q\,\|\,M$, Lemma lem:twist(iii) says that $h=g\otimes\chi$ is a newform on $\Gamma_0(Mq)$''). (2) The proof of $h\ne f$: the
revised paper argues by strong multiplicity one and notes ``Note that $h\otimes\chi$ is the $q$-depletion of $g$, not $g$, as $a_q(g)=\pm1$''; the original version wrote
$g=h\otimes\chi=f\otimes\chi=f'$, which is false as an identity of $q$-series. Steps (b6), (c6) give the argument with newform twists. (3) Case (c): the revised paper sets
``$f'=f_{E'}$ be its newform, of weight $2$ on $\Gamma_0(N/q)$ with $q\,\|\,N/q$, $a_q(f')=\pm1$ and $a_p(f')=\chi(p)a_p(f)$ for $p\ne q$ (so $f\otimes\chi$ is the $q$-depletion of $f'$ ...)''; the
original version called $f\otimes\chi$ itself a newform on $\Gamma_0(N/q)$, which is false, since its $q$-th coefficient vanishes while $a_q(f_{E'})=\pm1$
(\Cref{ex:modgalrep:twist-dictionary}(b)).
\end{paperlink}

\begin{remark}[case (a) at $q=2$, and the range of the hypotheses]\label{rem:modgalrep:qual-q2}
Case (a) at $q=2$ was added in the revised paper; the only input specific to $q=2$ is that for type $\kod{IV}$ or $\kod{IV}^*$ at $2$ the image of inertia on $T_3E$ is
cyclic of order $3$, so that the Serre exponent of $\bar\rho_{E,3}$ at $2$ is $1$ (\Cref{ch:galrep}). Case (b) is not claimed at $q=2$: there the Kodaira type does not
determine the image of inertia, and the twist by $\Q(\sqrt{q^*})$ is not available. In case (a) the hypothesis $3\nmid N$ excludes $q=3$. Below conductor $10{,}000$ there
are $1{,}051$ optimal curves with type $\kod{IV}$ or $\kod{IV}^*$ at $2$ (PARI codes $\pm4$), $3\nmid N$ and mod $3$ image neither reducible (labels \texttt{3B}, \texttt{3Cs}) nor $\Ns$, and
all of them have $3\mid\deg\varphi$ (Cremona's \texttt{degphi} and \texttt{galrep} files; none of them has complex multiplication); the first are \curve{52a1} ($\deg\varphi=3$),
\curve{76a1} ($6$) and \curve{92b1} ($6$). For $\ell\ge5$ only the irreducibility of $\bar\rho$ is used, as the paper says after the theorem.
\end{remark}

\begin{remark}[why exactly these types]\label{rem:modgalrep:why-these-types}
The proof needs a type at $q$ for which $\bar\rho$, or $\bar\rho\otimes\chi_{q^*}$, has Serre exponent at $q$ strictly smaller than the exponent of the curve it comes from, so that
the level can be lowered and raised back by the twist to a \emph{different} newform. By \Cref{ch:galrep} (inertia modulo $3$ by type) and
\Cref{prop:modgalrep:twist-conductor}:
\begin{enumerate}[label=(\roman*)]
\item Types $\kod{III}$, $\kod{III}^*$ at $q\ge5$: the Serre exponent of $\bar\rho_{E,3}$ and of its twist is $2$; nothing can be lowered. These are the types relevant to the image
$\Ns$ (Section sec:anomaly of the paper).
\item Type $\kod{I}_0^*$: $E'=E^{(q)}$ has good reduction at $q$ and $M=N(\bar\rho\otimes\chi)$ divides $N/q^2$. But \Cref{cor:modgalrep:optimal-level-E} may return $g=f_{E'}$ itself (it does
whenever $N(\bar\rho\otimes\chi)=N/q^2$), and then $h=g^\chi=f$: no new congruence. For \curve{539d1} at $\ell=3$ ($\kod{I}_0^*$ at $7$, twist \curve{11a3}), $N(\bar\rho\otimes\chi_{-7})=11$ and
$g=f_{11a}$, $h=f_{539d}=f$. The $\kod{I}_0^*$ contribution is reached by level raising instead (Theorem thm:quant and \Cref{thm:modgalrep:quant}).
\item Type $\kod{I}_n^*$ with $\ell\nmid n$: $\bar\rho\otimes\chi$ is ramified at $q$ with exponent $1$, so $M$ has $q^1$, $h=g^\chi$ has $q^2$, and \Cref{cor:modgalrep:optimal-level-E} may
return $g=f_{E'}$, giving $h=f$.
\item Types $\kod{II}$, $\kod{IV}$ and their duals at $\ell\ge5$: $\bar\rho_{E,\ell}$ has Serre exponent $2$ at $q$ ($\Phig_q$ has no $\ell$-torsion), and so has every twist.
\end{enumerate}
The paper's remark after Theorem thm:quant observes that the types of (ii) and (iii) are covered by Theorem thm:quant, with no hypothesis at $\Q(\sqrt{-3})$.
\end{remark}

\begin{example}[the paper's Example]\label{ex:modgalrep:paper-example}
The table lists, for the paper's four instances and the book's \curve{121c1}, the case of Theorem thm:qual, the type at $q$, the twist at curve level (an
\texttt{elltwist}, minimal model and table lookup), the Serre conductor that the proof lowers to and the newform there, the form $h$ of Lemma lem:old,
and the modular degree; $e_\ell$, $t_\ell$ are the paper's local terms (Definition def:terms; \Cref{ch:lfunctions}).
\begin{center}\small
\begin{tabular}{@{}lllllp{3.1cm}p{2.2cm}@{}}
\toprule
$E$ & case, $\ell$ & type at $q$ & twist & $M$, $g$ & $h$ & $\deg\varphi$\\
\midrule
\curve{52a1} & (a), $3$ & $\kod{IV}^*$ at $2$ & none & $26$, $f_{26b}$ & $f_{26b}$ & $3$; $e_3(E,2)=1$\\
\curve{121c1} & (a), $3$ & $\kod{IV}$ at $11$ & none & $11$, $f_{11a}$ & $f_{11a}$ & $6$; $e_3(E,11)=1$\\
\curve{121a1} & (b), $3$ & $\kod{II}$ at $11$ & \curve{121c2} ($\kod{IV}^*$) & $11$, $f_{11a}$ & $f_{121d}=f_{11a}\otimes\chi_{-11}$ & $6$; $e_3(E,11)=1$\\
\curve{539c1} & (c), $3$ & $\kod{I}_3^*$ at $7$ & \curve{77c1} ($\kod{I}_3$) & $11$, $f_{11a}$ & $f_{539d}=f_{11a}\otimes\chi_{-7}$ & $2^5\cdot3^2$; $e_3=t_3=1$\\
\curve{980i1} & (c), $5$ & $\kod{I}_5^*$ at $7$ & \curve{140b1} ($\kod{I}_5$) & $4\cdot5$, $f_{20a}$ & $f_{980g}=f_{20a}\otimes\chi_{-7}$ & $2^6\cdot3^2\cdot5$; $t_5=1$, $e_5=0$\\
\bottomrule
\end{tabular}
\end{center}
\begin{enumerate}[label=(\roman*)]
\item \emph{Every $n$.} Each congruence $f_E\equiv g'$ produced by Lemma lem:old holds for every $n$, not only at the primes the paper first listed: both sides lie in
$S_2(\Gamma_0(N))$ with integral coefficients, so a check up to Sturm's bound $\frac16N\prod_{p\mid N}(1+\frac1p)$ suffices (\Cref{ch:modforms}); the bounds are $14$, $22$, $22$, $112$, $336$
for $N=52$, $121$, $121$, $539$, $980$, and the congruences were checked with PARI/GP for all $n\le3000$. For the last three rows $M_h=N$, so $V_h=\C h$ and $g'=h$; for the first
two $g'=h-h(qz)$, the $q$-depletion of $h$ (\Cref{ex:modgalrep:old-examples}(b), (c)).
\item \emph{Curves versus classes.} The twist of \curve{121a1} by $\chi_{-11}$ is \curve{121c2}, of type $\kod{IV}^*$; the optimal curve \curve{121c1} has type $\kod{IV}$. The class 121d of $f_{11a}\otimes\chi_{-11}$
contains \curve{121d2}, the twist of \curve{11a1} ($\kod{I}_5^*$), and \curve{121d1}, the twist of \curve{11a3} ($\kod{I}_1^*$). The class 539d of $f_{11a}\otimes\chi_{-7}$ contains \curve{539d2}, the
twist of \curve{11a1}; the twist of \curve{20a1} by $\chi_{-7}$ is \curve{980g2}. Every quantity in the theorem depends only on the class: $r$ does, and $v_\ell(\deg\varphi)$ of
the optimal curve is what the theorem concerns.
\item \emph{$5\mid980$.} \curve{980i1} has multiplicative reduction at $5$, so $\ell=5$ divides $N$ and the example lies outside the hypotheses of Theorem thm:qual (and of Theorem
thm:main), as the paper says: ``As $5\mid980$, this last example illustrates the congruence, not the theorem.'' Accordingly the lowered level is
$N(\bar\rho_{140b1,5})\cdot5=20$, not $4$ (\Cref{ex:modgalrep:level-lowering}), and Lemma lem:old at $p=\ell$ and \Cref{thm:modgalrep:ars} at $\ell$ are not available.
\item \emph{Twists of newforms.} In the rows (b), (c), $h=g^\chi=g\otimes\chi$ by \Cref{cor:modgalrep:atkin-li}; the newform twists of $f_E$ itself are $f_{121a}^{\chi_{-11}}=f_{121c}$,
$f_{539c}^{\chi_{-7}}=f_{77c}$, $f_{980i}^{\chi_{-7}}=f_{140b}$, and only the first equals the naive twist (\Cref{ex:modgalrep:twist-levels}).
\end{enumerate}
\end{example}

\begin{paperlink}
The paper's Example lists the cases (a) at $q=2$ (\curve{52a1}), (b) (\curve{121a1}), (c) at $\ell=3$ (\curve{539c1}) and (c) at $\ell=5$ (\curve{980i1}), with the Sturm bounds $14$, $22$, $112$,
$336$; the row \curve{121c1} is the book's. All data in the table were recomputed (PARI/GP 2.17, Cremona's tables).
\end{paperlink}

\begin{remark}[the paper's remark on $\ell\mid N$]\label{rem:modgalrep:l-divides-N}
The remark after the paper's Example says that $\ell\nmid N$ is used in three places. (i) In the optimal-level theorem: \Cref{cor:modgalrep:optimal-level-E} needs $\bar\rho|_{D_\ell}$
finite, so that the level is $N(\bar\rho)$ and the character is trivial; for \curve{140b1} at $\ell=5$ the level is $N(\bar\rho)\cdot5$. (ii) In Lemma lem:old at $p=\ell$:
\Cref{thm:modgalrep:deligne-fontaine} needs the levels prime to $\ell$. (iii) In passing from $r$ to $\deg\varphi$, which needs only $\ell^2\nmid N$ and fails without it: for
\curve{99a1} (the 99A1 of \cite[Table 1]{ARS12}, $N=9\cdot11$), $\deg\varphi=4$ and $r=12$; for \curve{539d1}, $\deg\varphi=72$ and $r=504$; for \curve{726a1}, $\deg\varphi=144$ and $r=1584$ (the last three examples are the paper's, and all were
recomputed here as indices of $\Z f\oplus L$). The
remark ends: ``so do they at $q=3$ in cases (a) and (b), where $q=\ell$ and the theorem does not apply.'' Here $\ell=3$, so $q=3$ means $3\mid N$: the inertia group at $3$ then
acts on $T_3E$ through an infinite group (it contains the cyclotomic character on the determinant), and the local theory at $q=\ell$ is the $\ell$-adic one of
\Cref{sec:modgalrep:at-l}, not that of Lemma lem:IV.
\end{remark}

\section[Twists and levels]{Twists and levels in the paper's twist theorems}\label{sec:modgalrep:twisting}

Theorems thm:twist and thm:quant of the paper transport congruences from the newform $g=f_{E_S}$ of a twisted curve to $f=f_E$ by the twist map
$\iota$, and Theorem thm:quant reads the result as level raising. This section proves the facts about levels and twists on which both rest
and states the two theorems; their proofs are in \Cref{ch:congruence}, which uses the results below.

\begin{proposition}[the twist map]\label{prop:modgalrep:twist-map}
Let $S$ be a finite set of odd primes, $D=\prod_{q\in S}q$, $\chi_S=\prod_{q\in S}\chi_{q^*}$, and $\iota=\iota_S\colon h\mapsto h\otimes\chi_S=\sum_n\chi_S(n)a_n(h)q^n$. By \Cref{ch:hecke}, $\iota$
maps $S_2(\Gamma_0(M))$ into $S_2(\Gamma_0(\lcm(M,D^2)))$, preserves integrality, and $T_p\circ\iota=\chi_S(p)\,\iota\circ T_p$ for $p\nmid DM$. Moreover:
\begin{enumerate}[label=(\roman*)]
\item $\iota$ sends a generalised eigenform of the $T_p$, $p\nmid DM$, with eigenvalues $(a_p)$ to a generalised eigenform with eigenvalues $(\chi_S(p)a_p)$.
\item Let $g$ be a newform of level $M$, and suppose that $f=g^{\chi_S}$ is a newform of level $N$ with $\lcm(M,D^2)\mid N$. Then $\iota$ maps $g^\perp=\bigoplus_{h\ne g}V_h\subseteq S_2(\Gamma_0(M))$
(the $\T_M$-stable complement of $\C g$) into $f^\perp\subseteq S_2(\Gamma_0(N))$.
\item On newforms, $\iota(g)=(g^{\chi_S})^S$; so $\iota(g)=g^{\chi_S}$ iff $a_q(g^{\chi_S})=0$ for every $q\in S$, and $\iota(g)=\iota(g^S)$.
\item On Galois representations $\iota$ is $\rho\mapsto\rho\otimes\chi_S$; it preserves $\ad\rho$ and the true local factors $L_p(A_g,s)$.
\end{enumerate}
\end{proposition}

\begin{proof}
(i) From $T_p\iota=\chi_S(p)\iota T_p$ and $\chi_S(p)^2=1$: $(T_p-\chi_S(p)a_p)\iota=\chi_S(p)\,\iota\,(T_p-a_p)$, so a power of $T_p-a_p$ killing $h$ gives the same power of $T_p-\chi_S(p)a_p$
killing $\iota(h)$. (ii) For a newform $h\ne g$ of level dividing $M$, $\iota(V_h)$ consists of eigenforms of the $T_p$, $p\nmid DN$, with eigenvalues $\chi_S(p)a_p(h)$, which are
those of $h^{\chi_S}$; and $h^{\chi_S}\ne g^{\chi_S}=f$ by \Cref{prop:modgalrep:twists}(iii). By \Cref{prop:modgalrep:newforms-recap}(iv) at level $N$, $\iota(V_h)$ lies in $V_{h^{\chi_S}}$, which is
orthogonal to $V_f=\C f$. (iii) is \Cref{prop:modgalrep:twists}(ii) with $\Sigma(D)=S$, and $\iota$ kills the coefficients at $n$ with $(n,D)>1$, where $g$ and $g^S$ differ. (iv) is
\Cref{prop:modgalrep:twists}(ii) and \Cref{cor:modgalrep:delta-naive}.
\end{proof}

\begin{proposition}[levels of the twisted curve and of the depleted form]\label{prop:modgalrep:twist-levels}
Let $E/\Q$ have conductor $N$, $S$ a finite set of odd primes $q$ with $q^2\mid N$, $D=\prod_{q\in S}q$, $E_S=E\otimes\chi_S$ of conductor $N_S$, $f=f_E$ and $g=f_{E_S}$.
\begin{enumerate}[label=(\roman*)]
\item If $E$ has type $\kod{I}_n^*$ at every $q\in S$, then $N_S=N/D$; if type $\kod{I}_0^*$ at every $q\in S$, then $N_S=N/D^2$; if type $\kod{II}$, $\kod{III}$, $\kod{IV}$ or a dual at every $q\in S$, with $q\ge5$,
then $N_S=N$. In each case $\lcm(N_S,D^2)\mid N$, and in the third case also $\lcm(N,D^2)\mid N_S$.
\item $f=g^{\chi_S}$, and $\iota(g)=f$: the two $q$-series agree at every $n$.
\item In the first two cases the type of $g$ at $q\in S$ is (St), resp.\ (U), so $\delta_q(g)=1$, resp.\ $2$, and $N_S\prod_{q\in S}q^{\delta_q(g)}=N$: the $S$-depleted form $g^S$ lies in
$S_2(\Gamma_0(N))$, and $\iota(g^S)=f$.
\end{enumerate}
\end{proposition}

\begin{proof}
(i) At $p\notin S$ the character $\chi_S$ is unramified, so the exponents agree (Lemma lem:twist(i)). At $q\in S$, $v_q(N)=2$ ($q$ odd; \Cref{prop:modgalrep:twist-conductor}), and by
Lemma lem:twist(ii) the twist has type $\kod{I}_n$, $\kod{I}_0$, or a potentially good additive type, with exponent $1$, $0$, $2$
(\Cref{prop:modgalrep:twist-conductor}, rows (RSt), (RPS), last). The divisibilities are then immediate. (ii) $E=E_S\otimes\chi_S$, so $f=g^{\chi_S}$ by
\Cref{prop:modgalrep:twists}(iv); and $a_q(f)=0$ for $q\in S$ as $q^2\mid N$, so $\iota(g)=(g^{\chi_S})^S=f^S=f$ by \Cref{prop:modgalrep:twist-map}(iii). (iii) $E_S$ is multiplicative,
resp.\ good, at $q\in S$, so $g$ has type (St), resp.\ (U) (\Cref{prop:modgalrep:kodaira-dictionary}, \Cref{thm:modgalrep:local-types}), and the product is $N$ by (i). The depletion at
$q$ of a form of level $N'$ with $q\,\|\,N'$ is $g-a_q(g)g(qz)$, of level $N'q$, and with $q\nmid N'$ it is $g-a_q(g)g(qz)+q\,g(q^2z)$, of level $N'q^2$
(\Cref{lem:modgalrep:Up-oldspace}); doing this for each $q\in S$ puts $g^S$ in $S_2(\Gamma_0(N))$. Finally $\iota(g^S)=\iota(g)=f$.
\end{proof}

\begin{example}\label{ex:modgalrep:twist-levels-data}
\curve{539c1} ($\kod{I}_3^*$ at $7$): $E_S=\curve{77c1}$, $N_S=77=539/7$, and the depleted form
$g^S=f_{77c}(z)+f_{77c}(7z)$ (as $a_7(f_{77c})=-1$) lies at level $539$. \curve{539d1} ($\kod{I}_0^*$ at $7$): $E_S=\curve{11a3}$,
$N_S=11=539/49$, and $g^S=f_{11a}+2f_{11a}(7z)+7f_{11a}(49z)$ ($a_7(f_{11a})=-2$). \curve{121a1} ($\kod{II}$ at $11$): $N_S=121$.
\end{example}

\begin{theorem}[the paper's Theorem thm:twist]\label{thm:modgalrep:twist-monotonicity}
Let $E/\Q$ have conductor $N$, $S$ a set of odd primes, $D=\prod_{q\in S}q$, $\chi_S=\prod_{q\in S}\chi_{q^*}$, and $E_S=E\otimes\chi_S$ of conductor $N_S$. If $q^2\mid N$ for every $q\in S$ and
$\lcm(N_S,D^2)\mid N$, then $r_{E_S}\mid r_E$; if moreover $\lcm(N,D^2)\mid N_S$, then $r_{E_S}=r_E$.
\end{theorem}

The proof is in \Cref{ch:congruence}; in two lines: a congruence $g\equiv h\pmod m$ with $h\in g^\perp$ integral gives $f=\iota(g)\equiv\iota(h)\pmod m$ with $\iota(h)\in f^\perp$ integral
(\Cref{prop:modgalrep:twist-map}(ii), \Cref{prop:modgalrep:twist-levels}(ii)), so $m\mid r_E$ by \Cref{lem:modgalrep:crit}(i); and the second hypothesis makes the situation
symmetric. \Cref{prop:modgalrep:twist-levels}(i) shows that the hypotheses hold in the three cases listed after the theorem in the paper.

\begin{paperlink}
The paragraph after Theorem thm:twist (revised paper): ``In cases (a) and (b) Theorem thm:quant predicts a difference equal to the Euler term, which is positive in
case (a) at $\ell=3$ for every $q\ne3$, as $3\mid q^2-1$: for instance $v_3$ of the degrees is $2$ against $1$ for 539c1 and 77c1, $2$ against $1$ for 980i1 and 140b1, $1$ against
$0$ for 121d1 and 11a1, and $2$ against $0$ for 539d1 (type $\mathrm{I}_0^*$, $e_3(E,7)=2$) and 11a1.'' Recomputed from Cremona's \texttt{degphi} file: $\deg\varphi=288$, $6$, $2880$, $60$,
$24$, $1$, $72$ for \curve{539c1}, \curve{77c1}, \curve{980i1}, \curve{140b1}, \curve{121d1}, \curve{11a1}, \curve{539d1}, so the $3$-adic valuations are $2$, $1$, $2$, $1$, $1$, $0$, $2$, as stated, and each difference
is $e_3(E,q)$: $v_3(7^2-1)=1$, $v_3(7^2-1)=1$, $v_3(11^2-1)=1$, and $v_3(7-1)+v_3(8^2-(-2)^2)=2$ (\Cref{prop:modgalrep:level-raising-factor} with $a_7(f_{11a})=-2$). The original version
of the paper had claimed strict inequality ``only in case (b)'', which these pairs contradict.
\end{paperlink}

\begin{theorem}[the paper's Theorem thm:quant]\label{thm:modgalrep:quant}
Let $\ell$ be an odd prime with $\ell\nmid2N$ and $E[\ell]$ irreducible, $S$ the set of odd primes $q$ at which $E$ has type $\kod{I}_0^*$ or $\kod{I}_n^*$, and $E_S=E\otimes\chi_S$. Then
\[
v_\ell(\degphi)\ \ge\ v_\ell(r_{E_S})+\sum_{q\in S}e_\ell(E,q)\ \ge\ \sum_{q\in S}e_\ell(E,q).
\]
If moreover $\ell\notin S_f$ (automatic for $\ell\ge5$), the first inequality is an equality, and $v_\ell(r_{E_S})=v_\ell(\deg\varphi_{E'_S})$ for the optimal curve $E'_S$ of the class of
$E_S$.
\end{theorem}

The proof is in \Cref{ch:congruence}. Its four steps use, from this chapter: \Cref{prop:modgalrep:twist-levels}(iii) ($g^S$ at level $N$, $\iota(g^S)=f$),
\Cref{prop:modgalrep:twist-map}(i), (ii) (transport of complements), \Cref{lem:modgalrep:crit}(iii) and its proof (the congruence ideal of the eigencharacter of $g^S$ as a
largest congruence modulus, with the $\T$-stable complement in place of the Petersson complement), \Cref{lem:modgalrep:Up-oldspace} (the $U_q$-eigenvalues on
the span of the $g(dz)$: $0$ occurs once, the others are $\pm1$ or roots of $X^2-a_q(g)X+q$, units at $\ell\ne q$), and \Cref{thm:modgalrep:dfg14} with
\Cref{prop:modgalrep:level-raising-factor} (Step 3: $L_q^{\nv}(A_g,1)^{-1}=L_q(A_f,1)^{-1}$, the factor that $f$ omits).

\begin{remark}[Theorem thm:quant as level raising]\label{rem:modgalrep:quant-level-raising}
Theorem thm:qual lowers the level of $\bar\rho\otimes\chi$ and twists back. Theorem thm:quant goes the other way, as the paper's remark after it says: ``the level is
raised, from $N_S$ to $N_S\prod q^{\delta_q}=N$, and the price is the level-raising Euler factor''. The form $g=f_{E_S}$ of level $N_S$ is depleted at $S$, which places it at
level $N$ (\Cref{prop:modgalrep:twist-levels}(iii)); by Theorem thm:dfg14 the congruence ideal grows by $\prod_{q\in S}L_q^{\nv}(A_g,1)^{-1}$, which is $1-q^{-2}$ at the $\kod{I}_n^*$
primes and $\frac{q-1}q\cdot\frac{(q+1)^2-a_q(E_S)^2}{q^2}$ at the $\kod{I}_0^*$ primes (\Cref{prop:modgalrep:level-raising-factor}), exactly the true adjoint factors that $f$ omits
(\Cref{cor:modgalrep:delta-naive}); and $\iota$ transports the congruences to $f$. Examples at $\ell=3$ (degrees from \texttt{degphi}, congruence numbers computed as in
\Cref{thm:modgalrep:ars}): \curve{539d1}: $e_3(E,7)=v_3(6)+v_3(60)=2$, $v_3(r_{11a})=0$, $v_3(72)=2$; \curve{539c1}: $e_3(E,7)=v_3(48)=1$, $v_3(r_{77c})=v_3(6)=1$, $v_3(288)=2$; \curve{121d1}:
$e_3(E,11)=v_3(120)=1$, $v_3(r_{11a})=0$, $v_3(24)=1$. Equality holds in all three, as the second part of the theorem predicts ($3\notin S_f$ for these curves, whose mod $3$
images are surjective). The potentially good types $\kod{II}$, $\kod{III}$, $\kod{IV}$ and their duals are not in $S$: there $N_S=N$, $f$ is its own depletion, Theorem thm:dfg14 gives no
factor, and the omitted factor is reached only through Theorem thm:dfg37 (\Cref{ch:congruence}).
\end{remark}

\section[Congruence ideals in the literature]{The literature on congruence ideals and Tamagawa exponents; Taylor-Wiles in one page}\label{sec:modgalrep:literature}

Section 1.5 of the paper places its results among theorems of Ribet and Takahashi, Pollack and Weston, and Kim and Ota, which compare
congruence ideals at different levels with the ``Tamagawa exponent'' of $\rho$ at a multiplicative prime, and Remark rem:ko(2) derives one inequality of
the paper from the formula of Kim and Ota. This section states those results with their hypotheses, as far as the paper quotes them, and
sketches the Taylor-Wiles method behind Theorem thm:dfg37.

\begin{definition}[the Tamagawa exponent]\label{def:modgalrep:tamagawa-exponent}
Let $f$ be a newform of weight $2$, $\lambda\mid\ell$ with $\bar\rho_{f,\lambda}$ irreducible, and $q\ne\ell$ a prime dividing the level. The \emph{Tamagawa exponent} $t_q=t_q(\rho_{f,\lambda})$ is the
largest $t\ge0$ such that $\rho_{f,\lambda}\bmod\lambda^t$ (on any stable lattice; they are all homothetic, \Cref{prop:modgalrep:lattices}(iii)) is unramified at $q$. It is finite, as $\rho_{f,\lambda}$ is
ramified at $q$. This is the paper's $t_q$ and Kim and Ota's $t_f(q)$ \cite[Definition 2.12]{KimOta} (as the paper cites it).
\end{definition}

\begin{proposition}\label{prop:modgalrep:tamagawa-exponent}
For $f=f_E$ and a prime $q\,\|\,N$, $q\ne\ell$, $t_q=v_\ell(n_q)=t_\ell(E,q)$, whether the reduction is split or not.
\end{proposition}

\begin{proof}
This is Lemma lem:tamagawa-exp of the paper, proved in \Cref{ch:galrep}: in the basis of the Tate shape, $\rho(\sigma)=\sm{1}{n_qt_\ell(\sigma)}{0}{1}$ on $I_q$, with $t_\ell(I_q)=\Zl$, so
$\rho\bmod\ell^t$ is unramified at $q$ iff $\ell^t\mid n_q$; the twist by the unramified $\varepsilon$ does not change this.
\end{proof}

\begin{remark}[Ribet and Takahashi]\label{rem:modgalrep:ribet-takahashi}
Ribet and Takahashi \cite{RibetTakahashi97,Takahashi01} compare $\degphi$ with the degree of a parametrisation of $E$ by the Shimura curve attached to an indefinite
quaternion algebra of discriminant $N^-\,\|\,N$ ($N^-$ a product of an even number of primes of multiplicative reduction), and find the product of the $n_q$ over $q\mid N^-$,
up to primes $\ell$ at which $E[\ell]$ is reducible; Pasten \cite{Pasten} controls that error term. The precise statement is in \Cref{ch:moddeg}.
\end{remark}

\begin{theorem}[Pollack and Weston]\label{thm:modgalrep:pollack-weston}
Pollack and Weston \cite{PollackWeston11} conjectured that moving a prime $q$ into the ``new'' part of the level changes the congruence ideal by exactly
$\lambda^{t_q}$, and proved it for forms of weight $2$ and square-free level when $\bar\rho$ is ramified at two primes at least.
\end{theorem}

We do not prove this. The statement is the paper's summary in Section 1.5; we have not checked the hypotheses against \cite{PollackWeston11}.

\begin{theorem}[Kim and Ota]\label{thm:modgalrep:kim-ota}
Let $N=N^+N^-$ with $N^-$ square-free, $\ell\ge5$, $\ell\nmid N$, $\ell\notin S_f$, and $\bar\rho$ ramified at every $q\mid N^-$ with $q\equiv\pm1\pmod\ell$. Then
\[
\ord_\lambda\eta_f(N)=\ord_\lambda\eta_f(N^+,N^-)+\sum_{q\mid N^-}t_f(q),
\]
where $\eta_f(N)$ is the congruence ideal of $f$ in the full Hecke algebra of level $N$ and $\eta_f(N^+,N^-)$ that in the $N^-$-new quotient. Consequently, for $f=f_E$,
$v_\ell(r_E)\ge\sum_{q\mid N^-}t_\ell(E,q)$.
\end{theorem}

\begin{proof}
The formula is \cite[Theorem 1.1]{KimOta}, as the paper's Remark rem:ko(2) states it; we do not prove it. The consequence: $\eta_f(N^+,N^-)$ is an integral ideal, so
$\ord_\lambda\eta_f(N^+,N^-)\ge0$; $\ord_\lambda\eta_f(N)=v_\ell(r_E)$ by \Cref{lem:modgalrep:crit}(iii) (the congruence ideal of $\phi_f$ in $\T_\ell$ is $\ell^{v_\ell(r_E)}\Zl$); and $t_f(q)=t_\ell(E,q)$ by
\Cref{prop:modgalrep:tamagawa-exponent}, the primes of $N^-$ being multiplicative.
\end{proof}

\begin{paperlink}
Remark rem:ko(2): ``Since $\eta_f(N^+,N^-)$ is integral, this gives $v_\ell(r_E)\ge\sum_{q\mid N^-}t_\ell(E,q)$, which (eq:thm) contains, without the conditions $\ell\ge5$ and $q\not\equiv\pm1$;
the equality is theirs.'' Theorem thm:main of the paper, which assumes $\ell\notin S_f$, contains this inequality for every multiplicative prime, including the
primes $q\equiv\pm1\pmod\ell$ and $\ell=3$; it does not give the equality, which is the content of the theorem of Kim and Ota.
\end{paperlink}

\begin{remark}[Taylor-Wiles and $R=\T$ in one page]\label{rem:modgalrep:taylor-wiles}
Let $\bar\rho=\bar\rho_{f,\lambda}$ be absolutely irreducible, $\ell$ odd, and $\Sigma$ a finite set of primes not containing $\ell$. Mazur's \emph{deformation ring} $R_\Sigma$ is the
complete local $\mathcal{O}$-algebra that represents lifts of $\bar\rho$ (up to conjugation) which are minimally ramified outside $\Sigma\cup\{\ell\}$, unrestricted at the primes of
$\Sigma$, and flat (or Fontaine-Laffaille, \Cref{thm:modgalrep:crystalline}(ii)) at $\ell$, with determinant $\chi_\ell$. The localised Hecke algebra $\T_\Sigma=\T(N^\Sigma)_{\mathfrak m}$, at the
level $N^\Sigma$ obtained by raising the level at $\Sigma$, carries such a lift (by \Cref{thm:modgalrep:eichler-shimura} and \Cref{thm:modgalrep:carayol}), hence a surjection
$R_\Sigma\to\T_\Sigma$. The theorem of Wiles and of Taylor and Wiles \cite{Wiles95,TaylorWiles}, with Diamond's refinements, is that this map is an isomorphism of
complete intersections. Its proof patches together Hecke modules at auxiliary levels $N^\Sigma Q$, where $Q$ is a set of \emph{Taylor-Wiles primes} $q\equiv1\pmod{\ell^n}$ at which
$\bar\rho(\Frob_q)$ has distinct eigenvalues, chosen by Chebotarev so that the dual Selmer group is killed; finding them needs $\bar\rho|_{\Gal(\Qbar/\Q(\sqrt{\ell^*}))}$
absolutely irreducible. The \emph{numerical criterion} then equates the length of $\wp/\wp^2$, the dual of the Selmer group $H^1_\Sigma(\Q,W)$ of the adjoint, with
the valuation of the congruence ideal $\eta^\Sigma$. What Theorem 3.7 of \cite{DFG04} (the paper's Theorem thm:dfg37) needs is exactly this: $\ell\notin S_f$, that is
$\ell\nmid2N$ and $\bar\rho|_{\Gal(\Qbar/\Q(\sqrt{\ell^*}))}$ absolutely irreducible, and $\Sigma\not\ni\ell$ with $M_{f,\lambda}$ minimally ramified outside $\Sigma$. For $\ell\ge5$, $\ell\nmid N$ and
$E[\ell]$ irreducible the condition at $\Q(\sqrt{\ell^*})$ holds automatically (\Cref{ch:galrep}), so in the range of Theorem thm:main the hypothesis bites only for $\ell=3$
and image $N_s(\F_3)$. Theorem thm:qual avoids it: ``Two statements hold under weaker hypotheses than Theorem thm:main, by level lowering and twisting
rather than by the Taylor-Wiles method'' (the paper's Section 1.5). Details are in \Cref{ch:congruence} and \cite{DDT}.
\end{remark}

\section{Exercises}\label{sec:modgalrep:exercises}

The exercises run from routine to harder. Those marked (machine) need PARI/GP and Cremona's tables; their answers were computed for this book.

\begin{exercise}[$\ell=2$]\label{exer:modgalrep:l-two}
Show that \Cref{prop:modgalrep:lattices}(ii) fails for $\ell=2$: for $E\colon y^2=x^3-3x+1$, the cubic has square discriminant $81$, so $\Q(E[2])$ is a cyclic cubic field; deduce that
$\bar\rho_{E,2}$ is irreducible over $\F_2$ but not absolutely irreducible. Which step of the proof of \Cref{prop:modgalrep:lattices}(ii) breaks?
\emph{Hint.} The image is the cyclic subgroup of order $3$ of $\GL_2(\F_2)\simeq S_3$, which fixes no line of $\F_2^2$ but is diagonalisable over $\F_4$; complex conjugation acts
trivially, as $\Q(E[2])$ is totally real, and in any case $1=-1$ in $\F_2$.
\end{exercise}

\begin{exercise}[naive and newform twists]\label{exer:modgalrep:naive-newform}
(machine) (a) Show that the twists of \curve{11a1} by $\Q(\sqrt{-3})$ and by $\Q(\sqrt5)$ are \curve{99d2} and \curve{275b2}, and deduce from \Cref{cor:modgalrep:atkin-li} that $f_{11a}\otimes\chi_{-3}=f_{99d}$ and
$f_{11a}\otimes\chi_5=f_{275b}$. (b) Show that $f_{77c}\otimes\chi_{-7}=f_{539c}$ but $f_{539c}\otimes\chi_{-7}\ne f_{77c}$, and compute $(f_{77c}\otimes\chi_{-7})\otimes\chi_{-7}$.
\emph{Answer to (b).} \curve{77c1} has type (St) at $7$, so the first is \Cref{cor:modgalrep:atkin-li}; the second has $a_7=0\ne-1$; and $(f_{77c}\otimes\chi_{-7})\otimes\chi_{-7}=f_{77c}^{\{7\}}$.
\end{exercise}

\begin{exercise}[twists of curves and optimality]\label{exer:modgalrep:twist-optimal}
(machine) Show that the twist of the optimal curve \curve{121c1} by $\Q(\sqrt{-11})$ is \curve{121a2}, not the optimal curve \curve{121a1}, and that \curve{11a3} twists to \curve{121d1}. Explain
why, when $E[\ell]$ is irreducible, $v_\ell$ of the degree of the minimal parametrisation of any curve of the class equals $v_\ell(\deg\varphi_{E_0})$, so that Theorem thm:qual can be
stated for any $E$ in the class. \emph{Hint.} The isogenies in the class have degree prime to $\ell$ (\Cref{ch:galrep}).
\end{exercise}

\begin{exercise}[ramification of $\bar\rho$ at a Steinberg prime]\label{exer:modgalrep:steinberg-unramified}
Let $f$ have type (St) at $p\ne\ell$. Using \Cref{thm:modgalrep:local-types}(a), show that $\bar\rho_L$ is unramified at $p$ for one stable lattice $L$ if and only if the corner of
$\rho|_{I_p}$ is divisible by $\lambda$ on that lattice. For $f=f_{11a}$ and $p=11$, find all $\ell$ for which some stable lattice has unramified reduction at $11$, and compare
$T_5(\curve{11a1})$ with $T_5(\curve{11a3})$. \emph{Answer.} Only $\ell=5$ ($n_{11}=5$ for \curve{11a1}); on $T_5(\curve{11a3})$, where $n_{11}=1$, the reduction is ramified. The two lattices
are not homothetic, which is possible because $\bar\rho_{11a,5}$ is reducible.
\end{exercise}

\begin{exercise}[the conductor of twists]\label{exer:modgalrep:twist-table}
(machine) Check \Cref{prop:modgalrep:twist-conductor} on the twists \curve{14a1}$\otimes\chi_{-7}$ and \curve{15a1}$\otimes\chi_{-3}$ (row (St)), \curve{11a1}$\otimes\chi_{-7}$ (row (U)), \curve{539c1}$\otimes\chi_{-7}$ (row
(RSt)) and \curve{539d1}$\otimes\chi_{-7}$ (row (RPS)) by identifying the twisted curves. \emph{Answer.} \curve{98a3}, \curve{45a4}, \curve{539d2}, \curve{77c1}, \curve{11a3}.
\end{exercise}

\begin{exercise}[Serre conductors]\label{exer:modgalrep:serre-conductors}
Compute $N(\bar\rho_{E,\ell})$ and $N(\bar\rho_{E,\ell}\otimes\chi_{-7})$ for $(E,\ell)=(\curve{539c1},3)$ and $(\curve{980i1},5)$, using the formula $a_p=f_p-\dim\Phig_p[\ell]$. \emph{Answer.} $539$ and $11$;
$196$ and $4$. Explain why in the second case the newform attached to $\bar\rho\otimes\chi_{-7}$ has level $20$.
\end{exercise}

\begin{exercise}[characters of order divisible by $\ell$]\label{exer:modgalrep:nebentypus}
Let $\psi$ be a Dirichlet character and $\lambda\mid\ell$ a prime of $\Z[\psi]$ with $\psi\equiv1\pmod\lambda$. Show that $\psi=1$ if the order of $\psi$ is prime to $\ell$. Give a character $\psi\ne1$ of
conductor $7$ and a prime $\lambda\mid3$ with $\psi\equiv1\pmod\lambda$. Where does the proof of \Cref{cor:modgalrep:optimal-level-E} use that the character given by
\Cref{thm:modgalrep:optimal-level} has order prime to $\ell$? \emph{Answer.} A cubic character modulo $7$ takes values in $\mu_3$, and $\zeta_3\equiv1$ modulo $(1-\zeta_3)$.
\end{exercise}

\begin{exercise}[level raising for $f_{37a}$]\label{exer:modgalrep:raising-37}
(machine) For $f=f_{37a}$ and $\ell=3$ ($\bar\rho_{37a,3}$ is surjective), find the primes $q<25$, $q\ne3,37$, with $3\mid(q+1)^2-a_q^2$, and check that newforms of level $37q$ congruent to $f$
modulo a prime above $3$ exist exactly for those $q$. \emph{Answer.} $q=7$ ($63$), $13$ ($192$), $17$ ($324$); at each of these levels there is a congruent rational newform (and
at $q=7$, $17$ also forms with larger coefficient fields); for $q=2,5,11,19,23$ there is none.
\end{exercise}

\begin{exercise}[the level-raising factor at $\ell=3$]\label{exer:modgalrep:factor-11}
Compute $v_3$ of $\frac{q-1}q\cdot\frac{(q+1)^2-a_q(f_{11a})^2}{q^2}$ for $q=7$ and $q=13$. Use Theorem thm:quant to predict a lower bound for $v_3(\deg\varphi)$ for the optimal curves of the classes of
\curve{11a1}$\otimes\chi_{-7}$ and \curve{11a1}$\otimes\chi_{13}$, and check it in the tables (machine). \emph{Answer.} $2$ and $1+2=3$; the classes are 539d ($\deg\varphi=72$) and 1859a
($\deg\varphi=432=2^4\cdot3^3$), so the bound is attained in both.
\end{exercise}

\begin{exercise}[a toy congruence module]\label{exer:modgalrep:toy}
Let $r\ge2$ and $\T=\{(a,b)\in\Z^2:a\equiv b\pmod r\}$, a ring under componentwise operations, acting on $S=\Hom_\Z(\T,\Z)$ by $(t\cdot s)(t')=s(tt')$. The two projections $\phi,\psi\colon\T\to\Z$
lie in $S$ and are eigenvectors of $\T$ (with eigencharacters $\phi$ and $\psi$), playing the roles of $f$ and of an eigenform orthogonal to it. Compute $I=\ker\phi$, $J=\Ann_\T(I)$,
$\T/(I+J)$, the index $[S:\Z\phi\oplus\Z\psi]$, and the largest $m$ with $\phi-\psi\in mS$. \emph{Answer.} $I=0\times r\Z$, $J=r\Z\times0$, $\T/(I+J)\simeq\Z/r$, the index is $r$, and $(\phi-\psi)/r\in S$ because
$a-b\in r\Z$ on $\T$, so $m=r$.
\end{exercise}

\begin{exercise}[the congruence number of \curve{99a1}]\label{exer:modgalrep:99a1}
(machine) Compute $[S_2(\Gamma_0(99);\Z):\Z f\oplus L]$ for $f=f_{99a}$ and find $g\in L$ with $f\equiv g\pmod{12}$. Which hypothesis of \Cref{thm:modgalrep:ars} fails at $3$? \emph{Answer.} $12$;
$v_3(12)=1\ne0=v_3(4)$, and $9\mid99$.
\end{exercise}

\begin{exercise}[the hypothesis of Lemma lem:old]\label{exer:modgalrep:old-hypothesis}
Let $f$, $h$, $\ell$, $\lambda$ be as in \Cref{lem:modgalrep:old}, $p\,\|\,N$ with $\bar\rho$ unramified at $p$, and suppose that $p\,\|\,M$ (the excluded case). Show that $a_p(h)\equiv a_p(f)\pmod\lambda$ unless
$p\equiv-1\pmod\ell$, and that when $p\equiv-1\pmod\ell$ the two signs may differ. Why does this not contradict the lemma? \emph{Hint.} Both $\{a_p(f)p,a_p(f)\}$ and $\{a_p(h)p,a_p(h)\}$
are the eigenvalues of $\bar\rho(\Frob_p)$; if $a_p(h)=-a_p(f)$ then $p\equiv-1$. When $p\mid M$ the operator $U_p$ acts on $V_h$ by the scalar $a_p(h)$, so no choice is left.
\end{exercise}

\begin{exercise}[case (a) at $q=2$: \curve{76a1}]\label{exer:modgalrep:76a1}
(machine) \curve{76a1} has type $\kod{IV}^*$ at $2$ and $\kod{I}_1$ at $19$, and surjective mod $3$ image. Compute $N(\bar\rho_{76a1,3})$, find the newform $g$ given by \Cref{cor:modgalrep:optimal-level-E},
the eigenform of Lemma lem:old, and check $3\mid\deg\varphi$. \emph{Answer.} $N(\bar\rho)=38$; $g=f_{38b}$ (not $f_{38a}$); $a_2(f_{38b})=1$, so the eigenform is $f_{38b}(z)-f_{38b}(2z)$;
$\deg\varphi=6$.
\end{exercise}

\begin{exercise}[the step $h\ne f$]\label{exer:modgalrep:h-ne-f}
In case (b) of the proof, suppose one argued ``if $h=f$ then $g=h\otimes\chi=f\otimes\chi=f_{E'}$''. Show on \curve{121a1} ($g=f_{11a}$, $h=f_{121d}$) that $h\otimes\chi_{-11}$ is $f_{11a}^{\{11\}}$, not $f_{11a}$, so
the first equality is false; then give the correct argument with newform twists. \emph{Hint.} \Cref{conv:modgalrep:twists} and \Cref{prop:modgalrep:twists}(iii).
\end{exercise}

\begin{exercise}[another instance of case (c)]\label{exer:modgalrep:1175b1}
(machine) \curve{1175b1} has type $\kod{I}_3^*$ at $5$ and surjective mod $3$ image. Show that its twist by $\Q(\sqrt5)$ is \curve{235a1}, with types $\kod{I}_3$ at $5$ and $\kod{I}_1$ at $47$, so
$M=N(\bar\rho\otimes\chi_5)=47$. Show that the newform $g$ of level $47$ (coefficient field $\Q(y)$, $y^4-4y^2-y+1=0$) satisfies $a_p(g)\equiv\chi_5(p)a_p(\curve{1175b1})$ modulo the prime of
residue degree $1$ above $3$, and not modulo the other. What is $h$, and what is its level? Check $3\mid\deg\varphi$. \emph{Answer.} $h=g\otimes\chi_5$, a newform of level $47\cdot25=1175$ with
quartic coefficient field; $\deg\varphi=288$.
\end{exercise}

\begin{exercise}[a type $\kod{I}_0^*$ curve]\label{exer:modgalrep:I0star}
Run the construction of the proof of Theorem thm:qual on \curve{539d1} at $\ell=3$ (type $\kod{I}_0^*$ at $7$, twist \curve{11a3}) as if it were case (c), and show that it returns $h=f$.
Then explain how Theorem thm:quant still gives $9\mid\deg\varphi_{539d1}=72$. \emph{Hint.} \Cref{rem:modgalrep:why-these-types}(ii) and \Cref{rem:modgalrep:quant-level-raising}.
\end{exercise}

\begin{exercise}[Theorem thm:quant at $\ell=5$]\label{exer:modgalrep:quant-5}
(machine) For \curve{361b1} ($\kod{I}_1^*$ at $19$) and \curve{363a1} ($\kod{I}_1^*$ at $11$), whose mod $5$ representations are irreducible, find $E_S$, compute $e_5(E,q)=v_5(q^2-1)$, and compare
$v_5(\deg\varphi_E)$ with $v_5(\deg\varphi_{E'_S})+e_5(E,q)$. \emph{Answer.} $E_S$ in the classes 19a and 33a; $e_5=1$ in both; degrees $120$, $1$ and $180$, $3$: equality.
\end{exercise}

\begin{exercise}[Tamagawa exponents depend on the lattice when $\bar\rho$ is reducible]\label{exer:modgalrep:tamagawa-lattice}
Show that $t_q$ of \Cref{def:modgalrep:tamagawa-exponent} does not depend on the stable lattice when $\bar\rho$ is irreducible. For $f=f_{11a}$, $\ell=5$, $q=11$ compute the largest $t$ with
$\rho\bmod5^t$ unramified at $11$ on $T_5(\curve{11a1})$ and on $T_5(\curve{11a3})$. \emph{Answer.} $1$ and $0$ ($n_{11}=5$, resp.\ $1$).
\end{exercise}

\section{Notes and references}\label{sec:modgalrep:notes}

\paragraph{Sources followed.} The construction of $\rho_{f,\lambda}$ in weight $2$ follows Shimura \cite[Chapter 7]{Shimura71} and Diamond and Shurman
\cite[Chapter 9]{DiamondShurman}, with the integral and residual facts as in Diamond and Im \cite{DiamondIm} and Darmon, Diamond and Taylor
\cite{DDT}; the general construction is Deligne's \cite{Deligne}, and irreducibility in general is Ribet's \cite{Ribet77}. The automorphic representation
of a newform is from Jacquet and Langlands \cite{JacquetLanglands70}, conductors of representations of $\GL_2(\Q_p)$ from Casselman \cite{Casselman73}, the local Langlands
correspondence from Bushnell and Henniart \cite{BushnellHenniart06}, and local-global compatibility from Langlands \cite{Langlands73}, Deligne
\cite{Deligne73} and Carayol \cite{Carayol86}; Rohrlich \cite{Rohrlich94} is the reference for Weil-Deligne representations of elliptic curves, and
Loeffler and Weinstein \cite{LoefflerWeinstein12} for computing local components. The local statements at $\ell$ are quoted from Edixhoven
\cite[Section 2]{Edixhoven92} (as the paper cites them), Serre \cite{Serre72, Serre87}, Wiles \cite{Wiles88}, Fontaine and Laffaille \cite{FontaineLaffaille} and
Berger \cite{Berger}. Serre's conjecture is \cite{Serre87}, proved by Khare and Wintenberger \cite{KhareWintenberger} (with Khare \cite{Khare} and Kisin \cite{Kisin09});
level lowering is due to Mazur \cite{Mazur77}, Ribet \cite{Ribet90,Ribet91}, Carayol \cite{Carayol89} and Diamond \cite{Diamond95}, and the conductor of a residual
representation was studied by Livn\'e \cite{Livne89}; level raising is Ribet's \cite{Ribet84}, in the general form of Diamond and Taylor \cite{DiamondTaylor94}. The
congruence criterion goes back to Ribet \cite{Ribet83} and the perfect pairing to Ribet and to Abbes and Ullmo \cite{AU96}; Agashe, Ribet and Stein \cite{ARS12}
relate the congruence number to the modular degree. Twists of newforms are in Atkin and Li \cite{AtkinLi78}. The congruence ideal at higher level is
in Diamond, Flach and Guo \cite{DFG04}; Kim and Ota \cite{KimOta}, Pollack and Weston \cite{PollackWeston11}, Ribet and Takahashi \cite{RibetTakahashi97}, Takahashi
\cite{Takahashi01} and Pasten \cite{Pasten} are quoted as the paper summarises them. The Taylor-Wiles method is \cite{Wiles95,TaylorWiles}, surveyed in \cite{DDT}.

\paragraph{Theorem numbers.} The numbers printed are those the paper itself cites: \cite[Theorems 3.1(e), 3.15]{DDT}, \cite[Theorem 2.2]{Ribet83},
\cite[Lemma 2.1]{AU96}, \cite[Theorem 3]{AtkinLehner70}, \cite[Theorem 3.1]{AtkinLi78}, \cite[Theorem 2.2]{ARS12}, \cite[Section 2]{Edixhoven92}, \cite[Proposition 1.4(c)]{DFG04},
\cite[Theorem 1.1, Definition 2.12]{KimOta}. ``Theorem 7.14'' of \cite{Shimura71} is an unverified pointer from the edition we know. All other references are cited by
chapter, section or as a whole; we have not verified the numbering of the main theorems of \cite{Carayol86}, \cite{Ribet84}, \cite{Ribet90} or \cite{DiamondTaylor94},
nor the precise hypotheses of \cite{Ribet84}, \cite{Carayol89} (\Cref{rem:modgalrep:carayol-lemma}) and \cite{PollackWeston11}; the statements given here are those the
paper uses, or special cases of the published ones.

\paragraph{Normalisations.} Two normalisations of the local Langlands correspondence for $\GL_2$ are in use, and some sources attach to a newform the
dual of $\rho_{f,\lambda}$; this chapter avoids the issue by stating everything it uses on the Galois side (\Cref{thm:modgalrep:local-langlands} and the paragraph after it).
Our $\rho_{f,\lambda}$ has arithmetic Frobenius trace $a_p$ and determinant $\chi_\ell$, as $\rho_{E,\ell}$ does; DFG04 work with cohomology, where the geometric Frobenius has
trace $a_p$ (\Cref{thm:modgalrep:eichler-shimura}(iii)).

\paragraph{The two versions of the paper.} The book follows the revised version. The original version of 10 October 2026 differed in this chapter's area
in four places, all corrected in the revision and discussed in the paperlinks above: the lattice $M_{f,\lambda}$ was said to be homothetic to $T_\ell E$ rather
than to its dual (\Cref{thm:modgalrep:comparison-E}); Lemma lem:old said ``$\bar\rho$ being irreducible'' without ``hence absolutely irreducible''
(\Cref{lem:modgalrep:old}, Step 2); the proof of Theorem thm:qual treated the naive twists $f\otimes\chi$ and $h\otimes\chi$ as newforms in cases (b) and (c)
(\Cref{sec:modgalrep:proof-qual}); and the sentence after Theorem thm:twist claimed strict inequality only in case (b). The revised version also adds case (a) at $q=2$ and
the example \curve{52a1}, proves that the Example's congruences hold for every $n$ by Sturm's bound, and gives the reason for excluding $q=3$ in the remark after the
Example as $q=\ell$.

\paragraph{Quotations and computations.} Quotations of the paper are verbatim from the revised version, except that its internal cross-references are
printed as label names (``Lemma lem:twist(iii)''). Every number in this chapter (coefficients, Kodaira types, twists, Serre conductors, congruences, Sturm bounds,
congruence numbers, modular degrees, the level-raising tables and the count of $1{,}051$ curves) was recomputed with PARI/GP 2.17.2 \cite{PARI} through the Python package
\texttt{cypari2} and Cremona's tables \cite{Cremona97} for conductors below $10{,}000$. The congruence numbers were computed as indices $[S_2(\Gamma_0(N);\Z):\Z f\oplus L]$
(\Cref{lem:modgalrep:crit}(i)), with $S_2(\Gamma_0(N);\Z)$ the saturation of the $q$-expansion lattice of \texttt{mfinit([N,2],1)} up to twice the Sturm bound and $L$ the saturated
image of $T_p-a_p(f)$ for a prime $p\nmid N$ at which $a_p(f)$ is a simple eigenvalue; for example:
\begin{lstlisting}[language=GP]
mf = mfinit([99,2],1); n = 2*mfsturm([99,2]) + 20;
C = mfcoefs(mf, n); E = ellinit("99a1");
f = concat([0], ellan(E, n))~; M = mfheckemat(mf, 5);
S = matrixqz(C, -2); L = matrixqz(C*(M - ellap(E,5)), -2);
abs(matdet(lift(matinverseimage(S, matconcat([f, L])))))   \\ 12
\end{lstlisting}
Levels and twists of curves were identified by \texttt{elltwist}, \texttt{ellminimalmodel} and lookup in \texttt{allcurves}; newforms with irrational coefficients by
\texttt{mfeigenbasis}, \texttt{mffields} and reduction modulo each prime above $\ell$ (\texttt{idealprimedec}, \texttt{nfeltval}).
